\chapter{Introduction}\label{sec:introduction}

\section{Background}\label{sec:intro-background}

We consider a problem that arises in the study of higher-order
influence function estimators \cite{RobinsEtAl2008} and semiparametric
minimax lower bounds \cite{RobinsEtAl2009EJS} for nonlinear
functionals. We observe independent copies
of $(X,Y)$, where $X\in[0,1]^d$ has an unknown density $p$ bounded
above and away from zero, and let
\[
 f(x)=\E(Y\mid X=x),\qquad \Var(Y\mid X=x)=V.
\]
The regression function belongs to a fixed $s$-H\"older ball. Let
$Q_x$ denote the conditional law of $\varepsilon=Y-f(X)$ given $X=x$.
It has mean zero and variance $V$, and its fourth moment is uniformly
bounded, but it may otherwise depend on $x$. No smoothness is imposed on $p$.
We aim to estimate $V$, which satisfies
\begin{equation}\label{eq:variance-quadratic-identity}
 V=\E(Y^2)-\int_{[0,1]^d}f(x)^2p(x)\,\mathrm dx.
\end{equation}
Under the uniform H\"older-norm and conditional fourth-moment bounds,
the empirical mean of $Y^2$ estimates $\E(Y^2)$ at the root-$n$ rate.
Thus the only potentially non-root-$n$ component in
\eqref{eq:variance-quadratic-identity} is the quadratic regression
functional.

For $0<s<1$ and $d>4s$, Robins et al.\ \cite[Section~4]{RobinsEtAl2008}
give a close-pair construction with convergence rate
$n^{-4s/(d+4s)}$, which uses no smoothness of $p$. They divide
$[0,1]^d$ into small cubes, select one pair of observations in each
cube that contains at least two, and use the identity
\[
 \E\left[\frac{(Y_i-Y_j)^2}{2}\,\middle|\,X_i=x,X_j=y\right]
 =V+\frac{(f(x)-f(y))^2}{2}.
\]
By this identity, it is natural to average the squared differences
over pairs whose covariates are within distance $h$. The expected number of such pairs
is of order $n^2h^d$. Averaging the corresponding squared differences
gives a stochastic fluctuation bounded in root mean square by
$C\{n^{-1/2}+(n^2h^d)^{-1/2}\}$, while $s$-H\"older smoothness gives
bias $O(h^{2s})$. Balancing the bias with the pair term gives
$h\asymp n^{-2/(d+4s)}$; since $d>4s$, the root-$n$ term is smaller,
and the stated rate follows. Robins et al.\ introduced their
construction to show that a random design can give a faster rate than
an equally spaced design even when the design density has no smoothness
\cite[p.~338]{RobinsEtAl2008}. For $s$-H\"older regression on an
equally spaced $d$-dimensional grid, the corresponding
minimax rate for the root-mean-square risk is $n^{-2s/d}\vee n^{-1/2}$. Robins et
al.\ attribute the rate $n^{-2s/d}$, for $s<d/4$, to Wang et al.\ and
Cai et al.\ \cite[p.~338]{RobinsEtAl2008}. In the lower-bound
constructions of these papers, alternatives with a constant variance
give the lower bound $n^{-2s/d}$, in one dimension
\cite[Section~6.2]{WangBrownCaiLevine2008} and in several dimensions
\cite[Section~6]{CaiLevineWang2009}.\footnote{The construction
in~\cite[Section~6]{CaiLevineWang2009} uses cone-shaped bumps, which do
not belong to $C^s$ when $s>1$. For \(s>1\), the same lower bound
follows by using rescaled \(C^\infty\) bumps that equal one at the grid
points and are supported in balls of radius less than half the grid
spacing. Taking heights \(cn^{-s/d}\) with a sufficiently small fixed \(c>0\)
gives uniformly bounded \(C^s\) norms, while leaving the
prescribed values at the design points, and hence the testing
calculation, unchanged.}

Richardson and Rotnitzky \cite[Section~6]{RichardsonRotnitzky2014}
record the following question about the differentiable case,
attributed to James Robins: when $s>1$ but $s/d<1/4$, does there still
exist an estimator attaining $n^{-4s/(d+4s)}$? They write that analogy
with other nonparametric estimation problems would suggest that the
answer is ``yes''. As they also note
\cite[Section~6, footnote~50]{RichardsonRotnitzky2014}, the
squared-difference estimator of Robins et al.\ does not attain this
rate for $s>1$, because it does not exploit the differentiability of
the regression function. For $s>1$, its worst-case bias is generally
of order $h^2$, rather than $h^{2s}$. This shows only that the
estimator fails to exploit the additional smoothness; it does not rule
out that another estimator may attain a better rate. Their question
concerns rates of convergence; we study its minimax formulation,
which requires a uniform bound on the root-mean-square risk over the
model class specified below.

For $s>1$ and $d>4s$, we determine the minimax root-mean-square risk up
to a power of $\log n$. Let
\begin{equation}\label{eq:lambda}
 \lambda=\frac{2(s+1)}{d+4}.
\end{equation}
For all sufficiently large $n$, the minimax root-mean-square risk lies
between \(c\Psi_n\) and
\(C\Psi_n(\log n)^{\Gamma}\), where
\[
 \Psi_n=n^{-\lambda}e^{-\kappa\sqrt{\log n}}(\log n)^{(s-1)/(d+4)}
\]
and the constants $\kappa>0$ and $\Gamma>0$ are given explicitly in
\cref{thm:main}. These bounds identify $\lambda$ as the minimax
root-mean-square exponent. The minimax risk is smaller than $n^{-\lambda}$
by the factor $e^{-\kappa\sqrt{\log n}}$, up to powers of $\log n$,
with the same constant $\kappa$ in the upper and lower bounds, and the
ratio of the upper and lower bounds is at most $C(\log n)^\Gamma$.
Thus the rate \(n^{-4s/(d+4s)}\) is not uniformly attainable over the
stated model class, which gives a negative answer to the minimax
formulation of Robins's question. The upper bound also improves on the
fixed-design minimax rate. For $0<s\le1$, we give a direct proof of the
exact minimax rate $n^{-1/2}\vee n^{-4s/(d+4s)}$. For every $s>0$ and
$d\le4s$, ordinary least-squares residuals on a fixed partition attain
the parametric rate $n^{-1/2}$. Some of these results appear in
earlier versions of this paper, and related results were obtained in
concurrent work; we describe both in \cref{sec:intro-related}.

\section{Main results}\label{sec:main-result}
Fix \(s>0\) and an integer \(d\ge 1\). Fix an open set
\(\mathcal U\subset\mathbb R^d\) containing \([0,1]^d\), and constants
\begin{equation}\label{e:fix-constants}
 0<p_-<1<p_+<\infty,\qquad
 L_s>0,\qquad
 0<v_-<v_+<\infty,\qquad
 v_-^2<C_4<\infty.
\end{equation}

We adopt the following notation for H\"older norms.
Set
\begin{equation}\label{eq:holder-indices}
 \ell=\lceil s\rceil-1,
 \qquad \alpha=s-\ell\in(0,1].
\end{equation}
For a scalar function \(F\) on an open set
\(\Omega\subset\mathbb R^d\), define
\begin{align}
 [D^\ell F]_{\alpha;\Omega}
 &=\max_{|\gamma|=\ell}
   \sup_{\substack{u,v\in\Omega\\u\ne v}}
   \frac{|\partial^\gamma F(u)-\partial^\gamma F(v)|}
        {\|u-v\|^\alpha},
 \label{eq:holder-seminorm}\\
 \norm{F}_{C^s(\Omega)}
 &=\sum_{|\gamma|\le\ell}
   \norm{\partial^\gamma F}_{L^\infty(\Omega)}
   +[D^\ell F]_{\alpha;\Omega},
 \label{eq:holder-norm}
\end{align}
where \(\|\cdot\|\) is the usual Euclidean norm on \(\mathbb R^d\).
When \(s\notin\mathbb N\), this is the usual
\(C^{\ell,\alpha}\) norm, while integer \(s\) is interpreted as
\(C^{s-1,1}\).

For each \(n\in\mathbb N\), one observes independent and identically
distributed pairs \((X_i,Y_i)_{i=1}^n\).
Their common distribution is indexed by a parameter
\(\theta=(p,f,V,Q)\) satisfying the following three conditions.
Here, \(p\) is a probability density on \([0,1]^d\),
\(f\colon[0,1]^d\to\mathbb R\) is a regression function,
\(V>0\), and \(Q=(Q_x)_{x\in[0,1]^d}\) is a measurable
probability kernel from \([0,1]^d\) to \(\mathbb R\).
Define the regression error by
\(\varepsilon_i=Y_i-f(X_i)\).

\begin{enumerate}
\item[\textup{(i)}]
Each \(X_i\) has Lebesgue density \(p\) on \([0,1]^d\), and \(p\)
satisfies
\[
 p_-\le p(x)\le p_+
 \quad\text{for Lebesgue-a.e. }x\in[0,1]^d.
\]

\item[\textup{(ii)}]
The function \(f\) has an extension \(\widetilde f\) to
\(\mathcal U\) such that
\[
 \norm{\widetilde f}_{C^s(\mathcal U)}\le L_s,
\]
where the norm is defined in \eqref{eq:holder-norm}.

\item[\textup{(iii)}]
For \(p(x)\,\mathrm dx\)-almost every \(x\in[0,1]^d\),
the conditional distribution of \(\varepsilon_i\) given \(X_i=x\)
is \(Q_x\), and
\[
 \int_{\mathbb R}u\,Q_x(\mathrm du)=0,\qquad
 \int_{\mathbb R}u^2\,Q_x(\mathrm du)=V,\qquad
 \int_{\mathbb R}u^4\,Q_x(\mathrm du)\le C_4,
\]
with
\[
 v_-\le V\le v_+.
\]
The conditional variance \(V\) is constant in \(x\), but the
conditional distribution \(Q_x\) may otherwise depend arbitrarily
on \(x\).
\end{enumerate}

Since \(V^2\le C_4\) by conditional Jensen's inequality, the
effective variance interval is
\[
 \bigl[v_-,\min\{v_+,\sqrt{C_4}\}\bigr].
\]
The assumption \(v_-^2<C_4\) ensures that this interval is
nondegenerate.

Let
\[
 \Theta
 =\Theta(s,d,\mathcal U,p_-,p_+,L_s,v_-,v_+,C_4)
\]
denote this class of parameters \(\theta\).
For an estimator \(\widehat V_n\) of \(V\), define
\[
 \mathcal R_n(\widehat V_n;\Theta)
 =
 \sup_{\theta\in\Theta}
 \E_\theta\bigl(\widehat V_n-V\bigr)^2,
 \qquad
 R_n(\Theta)
 =
 \inf_{\widehat V_n}
 \mathcal R_n(\widehat V_n;\Theta),
\]
where the infimum is over all measurable functions of the
observations \((X_i,Y_i)_{i=1}^n\).
We write
\[
 \mathfrak r_n
 =R_n(\Theta)^{1/2}
 =\inf_{\widehat V_n}\sup_{\theta\in\Theta}
   \left\{\E_\theta\bigl(\widehat V_n-V\bigr)^2\right\}^{1/2}
\]
for the minimax root-mean-square risk.

We now state our four main results. Their constants depend only on the
fixed model parameters; $c_\epsilon$ may also depend on $\epsilon$.
\begin{theorem}\label{thm:main}\label{thm:bracket}
Let $s>1$ and let $d>4s$ be an integer. Let $\lambda$ be as in
\eqref{eq:lambda} and $\ell=\lceil s\rceil-1$ as in
\eqref{eq:holder-indices}, and put
\begin{equation}\label{eq:main-constants}
 \begin{gathered}
 \tau=\log\frac{\sqrt{p_+}+\sqrt{p_-}}{\sqrt{p_+}-\sqrt{p_-}},
 \qquad
 \kappa=4\sqrt{\frac{2\tau(s-1)(d-4s)}{(d+4)^3}},\\
 q_s=\binom{d+\ell}{d},
 \qquad
 \Gamma=\frac{d\bigl(q_s-\frac12\bigr)}{d+4}.
 \end{gathered}
\end{equation}
There exist $c,C>0$ such that, for all sufficiently large $n$,
\begin{equation}\label{eq:main-bracket}
 c\,\Psi_n
 \le\mathfrak r_n\le C\,\Psi_n(\log n)^{\Gamma},
 \qquad
 \Psi_n=n^{-\lambda}e^{-\kappa\sqrt{\log n}}
 (\log n)^{\frac{s-1}{d+4}}.
\end{equation}
\end{theorem}
The upper bound in \eqref{eq:main-bracket} is \cref{thm:upper-sharp},
proved in \cref{sec:upper-sharp}, and the lower bound is
\cref{thm:lower}, proved in \cref{sec:lower}.

The two bounds in \eqref{eq:main-bracket} differ by the factor
$(\log n)^\Gamma$; we do not know the correct power of $\log n$.
The estimator of \cref{sec:upper-sharp} is nearly minimax in the sense
that, for all sufficiently large $n$, its root-mean-square risk is
within a factor $C(\log n)^\Gamma$ of the minimax risk. The bounds agree
in the exponent $\lambda$ and in the constant $\kappa$ of the
stretched-exponential factor. In particular,
\[
 \log\mathfrak r_n=-\lambda\log n-\kappa\sqrt{\log n}+O(\log\log n).
\]
The constants $\tau$ and $\kappa$ depend on the density bounds
only through the ratio $p_+/p_-$. Here $q_s$ is the dimension of the
space of polynomials of degree at most $\ell$ in $d$ variables.
In \cref{thm:main}, the constants $c$ and $C$ and the threshold on $n$
depend only on $(s,d,p_-,p_+,L_s,v_-,v_+,C_4)$, and not on $\mathcal U$.

Neither proof gives an explicit threshold on $n$, and apart from $\kappa$
we have made no attempt to optimize the constants. The upper-bound proof
requires $\sqrt{\log n}$ eventually to dominate fixed multiples of
$\log\log n$, with constants depending on $q_s$. Our analysis is
asymptotic and does not establish the estimator's performance at
moderate sample sizes.

These bounds determine the minimax exponent. We state this as a corollary.
\begin{corollary}\label{thm:exponent}
Let $s>1$ and let $d>4s$ be an integer. For every $\epsilon>0$, there
exist $c_\epsilon,C>0$ such that, for all sufficiently large $n$,
\[
 c_\epsilon n^{-\lambda-\epsilon}\le\mathfrak r_n\le Cn^{-\lambda}.
\]
The minimax root-mean-square exponent satisfies
\[
 \lim_{n\to\infty}\frac{-\log\mathfrak r_n}{\log n}
 =\frac{2(s+1)}{d+4}.
\]
\end{corollary}
\begin{proof}
Since
$e^{-\kappa\sqrt{\log n}}(\log n)^{(s-1)/(d+4)+\Gamma}\to0$ as
$n\to\infty$, the upper bound in \eqref{eq:main-bracket} implies that
$\mathfrak r_n\le Cn^{-\lambda}$ for all sufficiently large $n$. For
every $\epsilon>0$, we have
$e^{-\kappa\sqrt{\log n}}(\log n)^{(s-1)/(d+4)}\ge n^{-\epsilon}$ for all
sufficiently large $n$, and then the
lower bound in \eqref{eq:main-bracket} gives
$\mathfrak r_n\ge cn^{-\lambda-\epsilon}$ for all sufficiently large
$n$. Taking logarithms in these two bounds, dividing by $\log n$, and
letting first $n\to\infty$ and then $\epsilon\downarrow0$, we obtain the
limit.
\end{proof}
The separation from the exponent in the question recorded in
\cite[Section~6]{RichardsonRotnitzky2014} is
\begin{equation}\label{eq:intro-gap}
 \frac{4s}{d+4s}-\lambda
 =\frac{2(s-1)(d-4s)}{(d+4s)(d+4)}>0.
\end{equation}
The lower bound rules out a uniform bound of order
\(n^{-4s/(d+4s)}\) on the root-mean-square risk, even after
multiplication by any fixed power of \(\log n\). The exponent also satisfies
\[
 \lambda-\frac{2s}{d}=\frac{2(d-4s)}{d(d+4)}>0,
\]
which means that the upper bound is faster than the fixed-design
minimax rate.
In the other direction, the minimax risk is not of the exact order
$n^{-\lambda}$, since the upper bound in \eqref{eq:main-bracket} implies
that $n^{\lambda}\mathfrak r_n\to0$. In particular, the weaker bound
$Cn^{-\lambda}$ recorded in \cref{sec:upper} is not the minimax rate.
The bound $Cn^{-\lambda}$ was proved in the third arXiv version of
this paper~\cite{AronowLopatto2026v3} and, in concurrent work, by
Dobriban et al.~\cite[Proposition~3.4]{DobribanEtAl2026}; see
\cref{sec:intro-related}.

For \(0<s\le1\), the exact minimax rate can be determined without the
constructions used for \(s>1\); see \cref{sec:small-s}.
\begin{theorem}\label{thm:main2}
Let $0<s\le1$ and let $d\ge1$ be an integer. There exists $C>0$
such that, for every $n\ge1$,
\begin{equation}\label{eq:main-rate}
 C^{-1}\bigl(n^{-1/2}\vee n^{-4s/(d+4s)}\bigr)
 \le\mathfrak r_n\le
 C\bigl(n^{-1/2}\vee n^{-4s/(d+4s)}\bigr).
\end{equation}
\end{theorem}

The residual-projection estimator of \cref{lem:projection-upper}
attains the parametric rate when $d\le4s$. Together with the lower
bound in \cref{sec:prelim-parametric}, it yields the following theorem.
\begin{theorem}\label{thm:parametric}
Let $s>1$ and let $d\le4s$ be a positive integer. There exists $C>0$
such that, for every $n\ge1$,
\begin{equation}\label{eq:parametric-rate}
 C^{-1}n^{-1/2}\le\mathfrak r_n\le Cn^{-1/2}.
\end{equation}
\end{theorem}


At $s=2$, the threshold $d\le4s$ of \cref{thm:parametric} has a
precedent in~\cite[Section~3.1]{Spokoiny2002}. For regression functions
with bounded second derivatives, i.i.d.\ Gaussian errors and an
i.i.d.\ design with a continuous positive density on a compact support,
root-$n$ accuracy is obtained there
when $d\le8$, with bounds that hold conditionally on the design on
events of probability arbitrarily close to one.

\section{Ideas of the proofs}\label{sec:intro-ideas}

\emph{The common approximation problem.} Both bounds are governed
by the same problem of polynomial approximation. Let $\tau$ be as in
\eqref{eq:main-constants}. On $[p_-,p_+]$, the reciprocal $1/z$ is
approximated by polynomials of degree $m$ with uniform error of order
$e^{-\tau m}$. This classical fact follows from the Chebyshev expansion
of the reciprocal \cite[equation~(2.4)]{Mathar2006}, and the geometric
factor $e^{-\tau}$ cannot be improved
\cite[Corollary~2.2]{KrausEtAl2012}. Dually, a signed measure on
$[p_-,p_+]$ that represents the value at the exterior point zero of
every polynomial of degree at most $m$ has total variation at least of
order $e^{\tau m}$, by the growth of Chebyshev polynomials outside the
interval \cite[Section~2.7, equation~(2.37)]{Rivlin1990}, and the rule
at the Chebyshev nodes attains this order (\cref{lem:finite-rules}).
The upper bound uses the first statement and the lower bound uses the
second, and this is the source of the common constant $\kappa$.

\emph{The upper bound.} We modify the close-pair construction of
Robins et al. For $s>1$, the bias of half the squared difference of two
responses at distance $h$ comes from the regression increment
$f(x')-f(x)$, which is of order $h$. We subtract a prediction of this
increment from the difference of the responses, and we use two
predictions computed from independent samples, one in each factor of
the square. As in double cross-fitting \cite[Section~2]{NeweyRobins2018},
the mean of the product then contains the product of the two mean
prediction errors, and the fluctuations of the predictions contribute
no bias. As in \cite[Section~3]{McCleanEtAl2024}, their covariances
across evaluation points remain in the variance, which we control by
averaging over many pairs (\cref{prop:pair-risk}). If the predictions
resolve $f$ at $K$ cells, the bias is of order $h^2K^{-2(s-1)/d}$,
while the pair statistic fluctuates at order $n^{-1}h^{-d/2}$. At the
sampling scale $K=n$, balancing these two terms gives the rate
$n^{-\lambda}$. The two-scale estimator of Dobriban et al.\
\cite{DobribanEtAl2026}, which projects out the local polynomial trend
in cells of diameter of order $n^{-1/d}$ and uses one close pair per
retained cell, attains the same rate; see \cref{sec:intro-related}.

To improve on $n^{-\lambda}$, we resolve $f$ at $K=ne^T$ cells with
$T>0$ tending to infinity, so that most cells contain no observation
and a local fit cannot be computed in them. As in Park~\cite[Section~3.2, equations~(9)--(14)]{Park2026} and in
our companion paper with
Kallus~\cite[Sections~4.1--4.3]{AronowKallusLopatto2026}, we replace the inverse of each local Gram matrix by a polynomial in
that matrix. The population fit is then a polynomial in local moments
of $p$ and $pf$. Each monomial in
these moments can be estimated without bias by a U-statistic
\cite{Hoeffding1948}, that is, by an average over distinct observations,
even when most cells are empty; for cell counts, this is the
falling-factorial estimate of powers of cell probabilities
\cite[equation~(21)]{JiaoEtAl2015}. This allows resolution below the
sampling scale (\cref{sec:U-factorial}). Related estimates of
functionals of an inverse Gram or covariance matrix appear in the
higher-order corrections of Robins et al., which use U-statistics of
increasing order with inverse Gram matrices computed under initial
nuisance estimates
\cite[Theorems~3.14--3.15]{RobinsEtAl2008}, and in the polynomial
estimates of Kong and Valiant \cite[Section~2.2]{KongValiant2018} and
of Chen, Liu and Mukherjee \cite[Section~3.2]{ChenLiuMukherjee2024},
who estimate quadratic forms involving matrix powers by U-statistics.

We express the prediction as a sum of increments of anchored local fits
over nested dyadic cells (\cref{sec:U-anchored}). Since consecutive
fits reproduce the same Taylor polynomial of $f$, each increment is
small.
To keep this Taylor cancellation, we clear the determinant of the parent
Gram matrix before we approximate the remaining inverse
(\cref{sec:U-inverse}). After a known preconditioning, each local Gram
matrix has its spectrum in $[p_-,p_+]$, and an approximation of
degree $m$ gives an increment error of relative order $e^{-\tau m}$, up to
a factor polynomial in $m$ (\cref{lem:U5}). The variance of the
unbiased lift grows factorially in the degree when the cells hold few
observations, and this limits the usable degree. We therefore lower
the degree as the resolution increases, along a quadratic profile in
the logarithm of the number of cells per observation
(\cref{sec:U-estimator}). With this profile, we may take $T$ of order
$\sqrt{\log n}$, that is, a resolution gain of $e^{O(\sqrt{\log n})}$
cells per observation, and we gain the factor $e^{-2(s-1)T/(d+4)}$, which is
$e^{-\kappa\sqrt{\log n}}$ up to powers of $\log n$
(\cref{sec:U-allocation}). Our companion paper
\cite[Sections~4.1--4.4]{AronowKallusLopatto2026} and Park
\cite[Sections~3.2--3.4]{Park2026} also keep the cancellation in
projection increments over nested partitions before approximating the
inverses of Gram matrices, and they also lower the degree on finer
partitions. Here the increments are anchored at a design point,
evaluated on close pairs, and normalized by an estimate of the noise
coefficient.

\emph{The lower bound.} We build on the ternary law of
\cref{sec:prelim-ternary}, which already appears in the third version
of this paper \cite{AronowLopatto2026v3}. Its probabilities are
quadratic in the regression value. By \eqref{eq:LB-heat-one}, we see
that a centered perturbation of $f$ with variance $\eta^2$ changes the
law of one response exactly as an increase of $V$ by $\eta^2$ does. Hiding a
change of the variance by randomizing the regression mean is the
principle of the lower-bound proof in~\cite[Section~4.5]{Spokoiny2002},
which places Gaussian regression values on smooth bumps, and of the
fixed-design lower bounds of Wang et al.\
\cite[Section~6.2]{WangBrownCaiLevine2008}, which match finitely many
Gaussian moments. With the ternary law, only observations that share a
local perturbation of the mean distinguish the two hypotheses. For
$0<s\le1$, independent perturbations of size $h^s$ on patches of scale
$h$ suffice, as in the one-dimensional random-design construction of Shen et al.\
\cite[Sections~2.2 and~5]{ShenGaoWittenHan2020}, whose design density
vanishes between the plateaus of the bumps. Here the ternary law makes the
cancellation for one observation exact, with a design density bounded
away from zero (\cref{sec:small-s}). Derivatives of expectations of bounded statistics of the data along
the resulting path are represented by a centered score whose second
moment, for $nh^d\le1$, is at most a constant times $h^{4s}$ times the
expected number of pairs of observations in a common patch, of order
$n^2h^d$, and the choice $h\asymp n^{-2/(d+4s)}$ gives the separation
$n^{-4s/(d+4s)}$. For $s>1$, the same balance would give
only $n^{-4s/(d+4s)}$, which is smaller than $n^{-\lambda}$ by a power
of $n$, and we must reduce the contribution of these pairs.

The likelihood of the observations in a patch contains the product of
the values of the design density at their positions. Robins et al.\
use priors in which the design density and the regression function
vary together \cite[Section~3]{RobinsEtAl2009EJS}. Here this
correlation has two roles. First, a joint update of the density and
of $f$ acts on the observations in a patch through a function of their
positions, provided that function has a separated representation
(\cref{sec:interp}). As in the third version, we choose this function by cardinal
interpolation, so that the covariance of the perturbation of $f$,
evaluated at the observed points, is diagonal (\cref{lem:interp}); here
count two uses a separate rule (\cref{sec:pair-new}). Second, the density factor
selects the number of observations in a patch. We integrate the updates of the density against
a signed measure that represents evaluation at zero, as in the dual
problem above. Among counts at most $m$, an update then acts only on
patches with a prescribed number $r\le m$ of observations; larger counts
can give aliases. The leading dependence of exterior evaluation on the cutoff
$m$ is $e^{\tau m}$; the full signed update also has factors that depend
on the selected count, the degree of the derivative rule and the
separated representation (\cref{lem:finite-rules,lem:local-operator}). Signed
measures dual to polynomial approximation were used to construct priors
with matching moments by Lepski, Nemirovski and Spokoiny
\cite[Section~4.4]{LepskiEtAl1999}, and for the reciprocal by Wu and
Yang \cite[Appendix~A]{WuYang2019}. The diagonal covariance imitates an
increase of the variance at each observation, which we cancel by
decreasing $V$. What remains are the defects of the interpolation,
the remainders of the fixed Taylor rule and of the rule for count two,
and the aliases at counts beyond the exactness range of the selectors
(\cref{sec:LB-rows,sec:LB-alias}).

These updates are signed. We realize them by a path of probability
priors: we give polynomial weights to finite histories of updates,
which we organize as heaps of pieces in the sense of Viennot
\cite{Viennot1986}. Conditional annihilation \eqref{eq:LB-annihilation}
keeps the law of the design density fixed along the path, and the
centering of the reference marks (\cref{lem:centered-row}) makes the
mean of its total mass equal to one; the method of bounded differences
then concentrates this mass near one. The weights are positive by a
deletion induction, which adapts the argument for hard-core partition
functions credited to Dobrushin and presented by Scott and Sokal
\cite[Section~5.1]{ScottSokal2005} (\cref{sec:LB-heaps}). A tilt by the
$n$th power of the total mass of the density cancels the normalization
of the likelihood at a fixed sample size. The third version Poissonized
the sample instead and reweighted the priors to cancel the random mass
in the Poisson likelihood \cite[Section~1.2, Step~3]{AronowLopatto2026v3}. Removing the last update computes
the score of the path on every data observable, and the local
cancellation makes this score small (\cref{prop:local}). Using the
score-to-risk comparison of \cref{lem:path-risk}, we then obtain the
lower bound (\cref{sec:LB-fisher}). This comparison is related to the information
inequality \cite[Chapter~2, Theorem~5.10]{LehmannCasella1998} and to
the comparison of mixtures of Cai and Low \cite[Section~2]{CaiLow2011};
we integrate the derivative of the mean of an estimator along the path
and bound the two endpoint biases.

\emph{The common balance.} Both proofs balance the logarithm $x$ of the
number of cells, or patches, per observation against the degree $m$ of
the approximation, or the largest selected count. Relative to
$n^{-\lambda}$, a larger $x$ improves the upper bound and weakens the
lower bound by the factor $e^{-2(s-1)x/(d+4)}$, while the approximation
problem enters through $e^{\tau m}$. In the upper bound, $mx$ must be at
most a multiple of $\log n$, to leading order, to control the factorial
variance of the unbiased lifts. In the lower bound, it must be at least
a multiple of $\log n$, to leading order, to make patches with more than
$m$ observations rare. In each bound, the optimal choice takes $x$ and
$m$ of order
$\sqrt{\log n}$, and in both bounds it gives the factor
$e^{-\kappa\sqrt{\log n}}$ with the same constant $\kappa$. The third
version already used a balance of this kind in its lower bound, with an
unspecified constant. Park~\cite[Sections~2 and~3.4]{Park2026} and our
companion paper~\cite[Section~1.3]{AronowKallusLopatto2026} use the same
balance in both of their bounds. The constructions below explain
these mechanisms where they are used, including the logarithmic
corrections.

\section{Related work}\label{sec:intro-related}

\emph{Variance estimation in nonparametric regression.} Difference-based
estimators of the variance go back to the first-difference estimator that
Brown and Levine \cite[Section~1]{BrownLevine2007} attribute to
von Neumann and, later, Rice \cite{Rice1984}. Gasser, Sroka and Jennen-Steinmetz \cite{GasserEtAl1986}
estimate the residual variance from local linear pseudo-residuals, Hall,
Kay and Titterington \cite{HallEtAl1990} construct difference
sequences that are asymptotically optimal, and Dette, Munk and Wagner
\cite{DetteEtAl1998} compare quadratic-form estimators and their
sensitivity to the regression mean. Munk et al.\ \cite{MunkEtAl2005}
study the estimation of the noise variance by differences in
multivariate regression with fixed covariates, and Brown and Levine
\cite{BrownLevine2007} and Wang et al.\ \cite{WangBrownCaiLevine2008}
study the estimation of a variance function by differences under fixed
design. Residual-based estimators form a second family. Hall and Marron
\cite{HallMarron1990} estimate the variance from squared nonparametric
residuals, and Fan and Yao \cite{FanYao1998} estimate a conditional
variance function by local linear smoothing of squared residuals. The
residual projection of \cref{lem:projection-upper} belongs to this
family, and variance estimators based on pseudo-residuals that
annihilate local linear trends were studied by Spokoiny
\cite[Sections~2.1--2.2]{Spokoiny2002}. Under random design, M\"uller,
Schick and Wefelmeyer \cite{MullerEtAl2003} estimate the error variance
by a covariate-matched U-statistic of squared differences. Du and
Schick \cite{DuSchick2009} give a covariate-matched version of the
pseudo-residual estimator of Gasser et al., which attains the
asymptotic performance of residual-based estimators under their
conditions. The close-pair construction of Robins et al.\ recalled in
\cref{sec:intro-background} does not use smoothness of the design
density.

Shen et al.\ \cite{ShenGaoWittenHan2020} study the estimation of the
variance under random design, when the standardized errors have a
common law independent of $X$. In
the multivariate case, they estimate the variance by a ratio of two
U-statistics of squared differences weighted by a product kernel, and
they prove that the minimax mean squared error is of order
$n^{-8\alpha/(4\alpha+d)}\vee n^{-1}$ when every coordinatewise H\"older
index lies in $(0,1]$, where $\alpha$ is the harmonic mean of these
indices \cite[Section~4.1]{ShenGaoWittenHan2020}. For isotropic
smoothness $s\le1$, this is the square of the rate in \cref{thm:main2},
and the estimator of \cref{prop:elementary-upper} for $0<s\le1$ has
the same ratio form, with the indicator of a common dyadic cube in
place of the product kernel. Their multivariate upper bound
\cite[Section~4.1, Proposition~1]{ShenGaoWittenHan2020} assumes
H\"older smoothness at most one in each coordinate, so their results do
not resolve the higher-smoothness case considered here. The design
densities used for the nonparametric component of their lower bound
vanish on sets of positive measure, so that construction does not yield
a lower bound for the present class. Conversely, the present
constructions for the lower bounds at rates slower than $n^{-1/2}$ use
the ternary law, under which the error law $Q_x$ varies with $x$ through
$f(x)$, so they do not establish the same lower bounds for a model in
which the errors have a common law independent of $X$. The design densities of the
present constructions are bounded above and away from zero, as the
present class requires.

\emph{Quadratic functionals and higher-order influence functions.} By
\eqref{eq:variance-quadratic-identity}, estimating $V$ amounts to
estimating the quadratic functional $\int f^2p$. Donoho and Nussbaum
\cite{DonohoNussbaum1990} develop minimax theory for quadratic
functionals in Gaussian models, Laurent \cite{Laurent1996} constructs
efficient estimators of integral functionals of a density from linear
and quadratic approximation terms, and Gin\'e and Nickl
\cite{GineNickl2008} estimate the integrated square of a density by
kernel U-statistics. Quadratic regression functionals are studied in
\cite{HuangFan1999,LiuEtAl2021}. The non-root-$n$ results in
\cite[Section~3.3]{LiuEtAl2021} assume that the design density belongs
to a H\"older class of smoothness strictly greater than
$2(1+\epsilon)s$, in addition to being bounded above and away from
zero. With bounded responses and under these density conditions, Liu
et al.\ state that the
constant variance can be estimated, adaptively over $s$, with mean
squared error of order $(n/\sqrt{\log n})^{-8s/(4s+d)}$ when $s<d/4$
\cite[Remark~3.5]{LiuEtAl2021}. Higher-order influence function
estimators are developed in
\cite{RobinsEtAl2008,RobinsEtAl2017}, and semiparametric minimax lower
bounds by the comparison of two mixtures in \cite{RobinsEtAl2009EJS};
van der Vaart \cite{vanDerVaart2014} reviews the underlying
higher-order tangent spaces. Mukherjee, Newey and Robins
\cite[Section~2.3]{MukherjeeNeweyRobins2017} construct empirical
higher-order estimators that are semiparametric efficient, under their
conditions, without estimating a smooth density. In the model in which
the conditional variance may depend on $X$ and the design density has
H\"older smoothness $\beta_g>0$, Robins et al.\ conjecture that the
rate in their equation~(1.3) is minimax, up to logarithmic factors,
when $\beta_g$ is too small for their estimators to attain
$n^{-4s/(d+4s)}$. For $\E\{\Var(Y\mid X)\}$, this rate tends to
$n^{-2s/d}$, up to a logarithmic factor, as $\beta_g\to0$
\cite[p.~338]{RobinsEtAl2008}.

\emph{Functional estimation under rough design.} For doubly robust
functionals such as the expected conditional covariance, let $s$ be the
average H\"older smoothness of the two nuisance functions. Robins et al.\
conjecture that their higher-order estimators, built from an initial
estimate of the density that is uniformly consistent, attain the rate $n^{-2s/d}$
in probability, up to a power of $\log n$, when $s<d/4$
\cite[p.~381]{RobinsEtAl2008}. For a counterfactual mean, Liu,
Mukherjee and Robins describe an empirical higher-order construction
whose bias is of order $n^{-2s/d}$, up to a power of $\log n$, without
smoothness of the density \cite[Section~4]{LiuMukherjeeRobins2020}.
McClean et al.\ \cite[Section~4.2, Theorem~1]{McCleanEtAl2024} construct
a double cross-fit estimator of the expected conditional covariance,
based on undersmoothed local polynomial regression. Under their moment
and support conditions, they bound its expected absolute error by a
term of order $n^{-2s/d}$, up to logarithmic factors and stated with
$O_P$ notation, when $s<d/4$ and the density satisfies only upper and
lower bounds. In concurrent work, Park \cite{Park2026} and
our companion paper \cite{AronowKallusLopatto2026} show that the
minimax exponent of such functionals is $2s/d$ below $s=d/4$ under the
same assumption on the density; we compare these works with ours
below. The results of the companion include the
quadratic functional $\int f^2p$ and the average conditional variance
$\E\{\Var(Y\mid X)\}$ of a bounded response whose conditional variance
may depend on $X$, with exponent $2s/d$ below $s=d/4$
\cite[Section~3.2]{AronowKallusLopatto2026}. In the present model,
where the conditional variance is constant, the exponent is instead
$4s/(d+4s)$ for $0<s\le1$ and $\lambda$ for $s>1$, when $d>4s$. For
this model when $s>1$ and $d>4s$, with a design density of H\"older
smoothness above an explicit threshold, Dobriban et al.\ show that a
clipped higher-order influence function estimator of Robins et al.\ has
mean squared error at most $Cn^{-8s/(d+4s)}$ if, in addition, the
unclipped construction satisfies the uniform integrated bounds in their
equation~(50) on squared conditional bias and conditional variance
\cite[Theorem~2.1]{DobribanEtAl2026}. In our lower-bound priors for
$s>1$, the design densities are updated patch by patch and can be
discontinuous, so these priors are not supported on a fixed H\"older
ball. Our lower bounds therefore do not apply to the H\"older-density
classes of Dobriban et al., where their conditional bound is faster;
they note this limitation for an earlier version of this paper
\cite[Section~4]{DobribanEtAl2026}.

\emph{Polynomial approximation and moment matching.} Lower bounds for
nonsmooth functionals are often built from two priors whose moments
agree up to some order, through the duality between moment matching
and best polynomial approximation; see Lepski, Nemirovski and Spokoiny
\cite[Section~4.4]{LepskiEtAl1999}, Cai and Low \cite{CaiLow2011} and
Wu and Yang \cite{WuYang2019}. Wu and Yang use best polynomial
approximation of the reciprocal on an interval bounded away from zero
\cite[Appendix~A]{WuYang2019}. The corresponding estimators replace the
functional by a polynomial and estimate its terms without bias, as in
\cite[Section~4.1]{CaiLow2011} and \cite{JiaoEtAl2015}. In the present
lower bound, the signed measures act on the values of the design
density, and in the upper bound, polynomial approximation replaces the
inverses of local Gram matrices. Related polynomial estimators of
functionals of inverse covariance matrices appear in
\cite{KongValiant2018,ChenLiuMukherjee2024}.

\emph{Concurrent work.} Dobriban, Mukherjee, Robins and Wang
\cite{DobribanEtAl2026} study the present model with the additional
assumption that the design density is H\"older of some positive order.
For $s>1$ and $d>4s$, they construct a two-scale estimator. It
partitions $[0,1]^d$ into cells of diameter of order $n^{-1/d}$,
projects out the local polynomial trend of degree $\lceil s\rceil-1$ in
suitable cells, and averages the squared normalized contrasts of one
close pair per cell. They prove that its mean squared error is at most
$Cn^{-4(s+1)/(d+4)}$, that is, that its root-mean-square risk is at most
$Cn^{-\lambda}$ \cite[Section~3.3, Proposition~3.4]{DobribanEtAl2026}.
Their proof of this bound uses only the upper and lower bounds on the
design density \cite[Appendix~A.5]{DobribanEtAl2026}. Their paper was
posted on 8 September 2026, one day before the third version of this paper,
which contained an estimator with the same bound. \Cref{thm:upper-sharp} improves the bound
$Cn^{-\lambda}$ by the factor $e^{-\kappa\sqrt{\log n}}$, up to powers
of $\log n$, and \cref{thm:lower} shows that the improved bound is
optimal over $\Theta$ up to a power of $\log n$.

Park \cite{Park2026}, in a paper posted on 4 October 2026, studies the
expected conditional covariance
$\E\{\Cov(A,Y\mid X)\}$ for responses bounded by one in absolute value,
conditional means $\E(A\mid X)$ and $\E(Y\mid X)$ in H\"older classes of
orders $\alpha$ and $\beta$, and an unknown design density between known
bounds $0<g_-<1<g_+$. The conditional covariance may depend on $X$. His
Theorem~1 states that the minimax mean squared error is of order
$n^{-1}$ when $\alpha+\beta\ge d/2$. For $w=(\alpha+\beta)/d<1/2$, the
theorem bounds it below by a constant multiple of
$n^{-2w}(\log n)^we^{-4\sqrt{\tau w(1-2w)\log n}}$ and above by a
constant multiple of
$n^{-2w}(\log n)^{w+q_\alpha+q_\beta+1/2}e^{-4\sqrt{\tau w(1-2w)\log n}}$,
where $\tau$ is computed from $g_\pm$ as in \eqref{eq:main-constants},
and $q_\alpha$ and $q_\beta$ are defined as $q_s$, with $\alpha$ and
$\beta$ in place of $s$ \cite[Section~2, Theorem~1]{Park2026}. When
$\alpha=\beta=s$, his root-mean-square exponent below the threshold is
$2s/d$, and his coefficient of $\sqrt{\log n}$ is
$(2/d)\sqrt{2\tau s(d-4s)}$. For $s>1$ and $d>4s$, our exponent
$\lambda$ exceeds $2s/d$ by $2(d-4s)/\{d(d+4)\}$, and our coefficient is
$\kappa$. The difference reflects the constant conditional variance of
the present model, which gives the identity for squared differences
recalled in \cref{sec:intro-background}. Park also reports a lower
bound for the average conditional variance of a single binary response
\cite[Section~2]{Park2026}; in that model, the conditional variance may
depend on $X$, and the lower bound does not transfer to the present
model. The two works share the threshold $d=4s$, the constant $\tau$,
and the identification of the coefficient of $\sqrt{\log n}$ in the
upper and lower bounds.

Both works use polynomial approximation on the density interval. In
both, the approximation of the reciprocal controls the bias of the
estimator, and in Park's lower bound it controls the separation; our
lower bound uses signed rules for the evaluation of polynomials at an
exterior point. Park's estimator starts from a double cross-fit average of
products of local polynomial residuals and subtracts bias corrections
on nested partitions. The corrections replace the inverses of the
population Gram matrices by Chebyshev matrix polynomials, they keep the
cancellation of the parent polynomial trends, they are estimated from
distinct observations, and their degrees decrease on finer partitions
\cite[Section~3]{Park2026}. Our estimator uses polynomial estimates of
the regression increment inside a pair statistic, and it relies on the
constant conditional variance. Park's main text describes his lower
bound through two mixtures with fixed cube probabilities, which give the
same distribution to the observations in each cube that contains at
most $r+1$ observations, where $r\ge1$ is the number of matched
moments, and whose targets differ by a signed reciprocal
moment \cite[Section~2]{Park2026}. Our lower bound selects counts by
signed local updates, which we realize by positive priors at a fixed
sample size. Park cites the third version and describes its
bounds, which have an unspecified constant in the
stretched-exponential factor \cite[Sections~1 and~8]{Park2026}; the
present version identifies this constant as $\kappa$ in both bounds.
Our comparison is based on the main text of arXiv:2610.05006v1. Its
proofs are deferred to a supplement that is not part of the arXiv PDF
or source files, as checked on 8 October 2026.

In our companion paper with Kallus \cite{AronowKallusLopatto2026}, we
study missing-at-random means, treatment effects and expected
conditional covariances in fully nonparametric models under rough
design. For two nuisance functions with average H\"older smoothness $s$,
we obtain the exponent $2s/d$ below $s=d/4$, and we identify the
coefficient of $\sqrt{\log n}$ in the upper and lower bounds
\cite[Theorem~2.3]{AronowKallusLopatto2026}. For the expected
conditional covariance, the exponent and the coefficient agree with
those of Park. The estimator for the main theorem of the companion also approximates
the inverses of local Gram matrices by polynomials on the interval that
contains their eigenvalues and estimates the polynomial terms without
bias, but it estimates the target directly, without regression fits or
close pairs of observations. The lower bound of the companion
follows the strategy of the third version, with a Poissonized sample,
but its priors give the design density a random
phase, and the two mixtures are separated through a Fourier
coefficient of the reciprocal of the density.

\emph{Earlier versions of this paper.} Versions~1 and~2 of
arXiv:2607.13170 \cite{AronowLopatto2026v1,AronowLopatto2026v2}, posted
on 14 July and 17 August 2026 under the title ``On Rates Attainable
under Random Design: A Negative Answer to a Problem of Robins'', proved
that $\mathfrak r_n\ge cn^{-\beta}$ for all $n\ge1$ when $s>1$ and
$d>4s$, where $\beta=\{d(3s+1)+8s\}/\{(d+2s)(d+4)\}$. Since
$\beta<4s/(d+4s)$, this lower bound already shows that the rate in the
question recorded by Richardson and Rotnitzky is not uniformly
attainable. Version~2 also established the exact minimax rate of
\cref{thm:main2} for $0<s\le1$. Version~3 \cite{AronowLopatto2026v3},
posted on 9 September 2026 under the title ``Nearly-minimax variance
estimation under rough random design'', proved that
$cn^{-\lambda}e^{-C\sqrt{\log n}}\le\mathfrak r_n\le Cn^{-\lambda}$ for
$s>1$ and $d>4s$, with an unspecified constant $C$ in the
stretched-exponential factor. These bounds identified the exponent
$\lambda$, and version~3 also proved \cref{thm:parametric}. Its upper bound
used an interpolation estimator built from corrected response
differences, and its lower bound used a Poissonized sample and
approximate matching of local moments. The present version identifies the
constant $\kappa$ in both bounds, so that the upper bound improves to
$C\Psi_n(\log n)^\Gamma$ and the two bounds differ only by the factor
$(\log n)^\Gamma$. It gives different proofs of
\cref{thm:main2,thm:parametric}, and its lower bound for $s>1$ uses signed local
updates and positive priors at a fixed sample size.

\section{Organization and conventions}\label{sec:intro-organization}
\Cref{sec:upper} proves the parametric upper bounds by residual
projection and the upper bounds for low smoothness by pair
differences. Its common risk bound for pair statistics is also used in
\cref{sec:upper-sharp},
which proves the refined upper bound for $s>1$ and $d>4s$.
\Cref{sec:lower-prelim} collects the inputs common to the lower bounds:
the ternary law, one score-to-risk comparison, and the smooth
weights; it also contains the proof of the parametric lower bound.
\Cref{sec:small-s} proves \cref{thm:main2}, introducing the local
cancellation and score argument in its elementary form. \Cref{sec:lower}
develops this argument into the lower bound in \cref{thm:main}, using
interpolation, count selection and positive priors built from signed
local updates. \Cref{thm:exponent} is proved immediately after its
statement.

All logarithms are natural. Constants may increase between occurrences.
For a finite signed measure $\mu$, $\norm{\mu}$ is the total mass of
$|\mu|$, that is, the total variation of $\mu$ without a factor $1/2$.

The constants of \cref{thm:main} are recalled in
\cref{thm:upper-sharp}. Apart from these and from the
model of \cref{sec:main-result}, each of
\cref{sec:upper,sec:upper-sharp,sec:small-s,sec:lower}
fixes its own notation, and several letters, among them $h$, $m$, $q$,
$M$, $N$, $\rho$, $\mu$ and $\Gamma$, are also used there
with local meanings, declared where they occur.

\chapter{Pair differences and elementary upper bounds}
\label{sec:upper}
Least-squares residuals directly give all of the parametric upper
bounds. For low smoothness, we obtain the remaining upper bounds from
half the squared difference of two nearby responses. We first record
the resulting rates, which include a consequence of the upper-bound
construction for higher smoothness.

\begin{proposition}\label{prop:elementary-upper}
Let $s>0$, $d\ge1$, and $t=\min\{s,1\}$. There exists a measurable
estimator based on $n$ observations such that, for every $n\ge1$,
\begin{equation}\label{eq:upper-simple}
 \sup_{\theta\in\Theta}
 \bigl\{\E_\theta(\widehat V_n-V)^2\bigr\}^{1/2}
 \le C\left(n^{-1/2}\vee n^{-2(s+t)/(d+4t)}\right).
\end{equation}
\end{proposition}

For $s>1$, this proves the upper bound in \cref{thm:parametric}
and the simpler bound $Cn^{-\lambda}$ when $d>4s$. For $0<s\le1$, it
is the upper bound in \cref{thm:main2}.
We prove every case with $d\le4s$, including equality, in the next
lemma. The bound for low smoothness is proved below using pair
differences. For $s>1$ and $d>4s$, \cref{thm:upper-sharp} implies the
bound $Cn^{-\lambda}$ for all sufficiently large $n$, since its
stretched exponential dominates its logarithmic factor. The constant
estimator $v_-$ covers the remaining sample sizes.

We write $\mathbb O=[0,1]^d\times\mathbb R$ for the observation
space, equipped with its Borel $\sigma$-field, and $\mu_\theta$ for the
law of one observation under $\theta$.

\section{Parametric upper bounds by residual projection}
\label{sec:projection-upper}

When we fit cellwise polynomials, a fixed fraction of the observations
remains available as residual degrees of freedom. The approximation error is
small even if the design points are irregular or the local fit is
rank deficient. Related variance estimators based on normalized local linear
pseudo-residuals are studied in~\cite[Sections~2.1--2.2]{Spokoiny2002},
for fixed design and conditionally on an i.i.d.\ random design.

\begin{lemma}\label{lem:projection-upper}
For every $s>0$ and $d\ge1$, there exists a Borel estimator satisfying
\begin{equation}\label{eq:projection-risk}
 \sup_{\theta\in\Theta}\|\widehat V_n-V\|_{L^2}
 \le C\bigl(n^{-1/2}+n^{-2s/d}\bigr),\qquad n\ge1.
\end{equation}
In particular it attains $Cn^{-1/2}$ whenever $d\le4s$.
\end{lemma}

\begin{proof}
Let $\ell=\lceil s\rceil-1$ and $q=\binom{d+\ell}{d}$. For $n\ge2q$,
we set $k=\lfloor(n/(2q))^{1/d}\rfloor$ and partition $[0,1]^d$ into
$K=k^d$ cubes, assigning their faces by a fixed convention. Then
$K\asymp n$ and $qK\le n/2$. Let $B_X$ evaluate a fixed basis of
the cellwise degree-$\ell$ polynomials at $X_1,\ldots,X_n$, and let
$P_X$ be the orthogonal projection onto its column space. We define
\[
 A_X=I_n-P_X,\qquad d_X=\operatorname{tr}A_X\ge n-qK\ge n/2,\qquad
 \widetilde V_n=\frac{Y^\top A_XY}{d_X}.
\]
The projection is Borel even at changes of rank, since
$P_X=\lim_{\epsilon\downarrow0}
B_X(B_X^\top B_X+\epsilon I_{qK})^{-1}B_X^\top$.

Let $f_X=(f(X_i))_{i\le n}$. On each cube, the Taylor polynomial centered at a fixed corner of
that cube approximates $f$ uniformly to within $CK^{-s/d}$.
For $\ell=0$, this is the H\"older bound; otherwise, it follows by
integrating the order-$\ell$ H\"older remainder along a segment in
the cube. Since these polynomials give a vector in the range of $P_X$,
we have
\[
 \|A_Xf_X\|^2\le CnK^{-2s/d},\qquad
 \E(\widetilde V_n\mid X)-V=\frac{\|A_Xf_X\|^2}{d_X}
 \le CK^{-2s/d}.
\]
Conditional on $X$, the errors are independent and centered, with
variance $V$ and fourth moments at most $C_4$. Expanding the quadratic
form of a symmetric matrix $A$, we find
\[
 \Var(\varepsilon^\top A\varepsilon\mid X)
 =\sum_i A_{ii}^2\Var(\varepsilon_i^2\mid X)
   +4V^2\sum_{i<j}A_{ij}^2
 \le C\operatorname{tr}(A^2).
\]
We also have $\Var(2f_X^\top A_X\varepsilon\mid X)
=4V\|A_Xf_X\|^2$. Using $A_X^2=A_X$ and
$\Var(U+W)\le2\Var U+2\Var W$, we obtain
$\Var(\widetilde V_n\mid X)\le C/n$.
Then the conditional mean-square error is at most
$C(n^{-1}+K^{-4s/d})$, uniformly in the design points.
We clip $\widetilde V_n$ to $[v_-,v_+]$ and average over $X$ to obtain
\eqref{eq:projection-risk}. For $n<2q$, we use the constant estimator
$v_-$ and increase $C$.
\end{proof}

\section{A common risk bound for pair statistics}\label{sec:pair-risk}

We partition $[0,1)^d$ into dyadic cubes of side $h$ and set
$k_h(x,x')=h^{-d}$ when $x,x'$ belong to the same cube of this
partition, and $k_h(x,x')=0$ otherwise, in particular when $x$ or $x'$
lies outside $[0,1)^d$. On the pair set
$\{k_h>0\}$, we denote points by $\mathsf w=(x,x')$ and let
$M_h(d\mathsf w)=k_h(x,x')p(x)p(x')\,dx\,dx'$. Then
\begin{equation}\label{eq:U7-mass}
 1\le\int dM_h\le p_+,\qquad
 M_h\{(x,x'):x\in A\}\le p_+^2|A|
 \quad(A\subset[0,1)^d\text{ Borel}).
\end{equation}
Indeed, writing $p_I=\Pp(X\in I)$, we see that the total mass is
$h^{-d}\sum_Ip_I^2$. By Cauchy--Schwarz and $\sum_Ip_I=1$, we obtain
the lower bound, and $p_I\le p_+h^d$ gives the upper bound.
The second assertion follows by restricting the first integration
to $A$.

A vector field $c(\mathsf w)=(1,c_2(\mathsf w),c_3(\mathsf w))^\top$
represents the affine pair residual $y'+c_2(\mathsf w)y+c_3(\mathsf w)$.
We allow random coefficients to cover predicted increments; the constant
field $(1,-1,0)^\top$ gives the ordinary response difference. For this
field, the ratio $\widehat N/\widehat D$ of the statistics defined below
is the average of half the squared response differences over pairs of
observations in a common cube. M\"uller, Schick and Wefelmeyer
\cite{MullerEtAl2003} estimate the variance by a covariate-matched
U-statistic of this kind, and Shen et al.\
\cite[Section~4.1, equation~(19)]{ShenGaoWittenHan2020} use the same
ratio form with a product kernel.

Independent pilots are useful because the mean of their product
contains the product of their mean prediction errors, without a pilot
variance term. The fluctuations of the pilots still matter, but the test-function
bound below controls them once many pairs are averaged. In the double
cross-fitting of Newey and Robins \cite[Section~2]{NeweyRobins2018},
two nuisance functions are likewise estimated from independent samples,
which removes the additional nonlinear bias that remains under single
cross-fitting. McClean et al.\
\cite[Section~3, Lemma~1]{McCleanEtAl2024} bound the resulting bias by
products of mean nuisance errors, and in their analysis the covariances
between evaluations of the same nuisance estimates at different points
remain in the variance. When the
observed response at the first design point of the ordered pair is used
in place of $f(x)$, the noise
coefficient changes, and the denominator estimates exactly that
normalization.

\begin{proposition}[A common pair-risk bound]\label{prop:pair-risk}
Let $c_{\mathcal A},c_{\mathcal B}$ be independent, identically
distributed jointly Borel random vector fields on the pair set,
with first coordinate one, independent of an evaluation sample of
$m\ge2$ observations. Suppose $m\ge c_*n$ for a fixed $c_*>0$.
Put $\bar c=\E c_{\mathcal A}$ and
$\zeta_{\mathcal S}=c_{\mathcal S}-\bar c$ for
$\mathcal S\in\{\mathcal A,\mathcal B\}$. Suppose that, uniformly over
$\theta\in\Theta$, for deterministic $b,W\ge0$ and a fixed $C_p<\infty$,
\begin{equation}\label{eq:pair-pilot-bounds}
\begin{gathered}
 \sup_{\mathsf w}|\bar c(\mathsf w)|\le C_p,\qquad
 \sup_{\mathsf w}\E|\zeta_{\mathcal S}(\mathsf w)|^2\le C_pW,\\
 \left\|\int\nu^\top\zeta_{\mathcal S}\,dM_h\right\|_{L^2}
 \le C_p\sqrt{W/n}\,\|\nu\|_{L^2(M_h)}
 \quad\text{for every Borel }\nu\in L^2(M_h;\mathbb R^3),\\
 r_h=\bar c^\top\mathbf f,\qquad
 \mathbf f(\mathsf w)=(f(x'),f(x),1)^\top,\qquad
 \|r_h\|_\infty\le b.
\end{gathered}
\end{equation}
For $O_i=(X_i,Y_i)$, let $\mathsf w_{ii'}=(X_i,X_{i'})$,
$z_{ii'}=(Y_{i'},Y_i,1)^\top$, $P_0=\operatorname{diag}(1,1,0)$,
and $(m)_2=m(m-1)$. Define
\begin{equation}\label{eq:pair-statistics}
\begin{aligned}
 \widehat N&=\frac1{2(m)_2}\sum_{i\ne i'}k_h(X_i,X_{i'})
 c_{\mathcal A}(\mathsf w_{ii'})^\top z_{ii'}z_{ii'}^\top
 c_{\mathcal B}(\mathsf w_{ii'}),\\
 \widehat D&=\frac1{2(m)_2}\sum_{i\ne i'}k_h(X_i,X_{i'})
 c_{\mathcal A}(\mathsf w_{ii'})^\top P_0c_{\mathcal B}(\mathsf w_{ii'}),
\end{aligned}
\end{equation}
reading a summand as zero when $k_h=0$. These statistics are square
integrable and $\E\widehat D\ge1/2$. If a fixed $a>0$ satisfies
$\E_\theta\widehat D\ge2a$ for every $\theta$, then, with clipping
denoting projection onto the displayed interval,
\begin{equation}\label{eq:pair-risk}
\begin{gathered}
 \widehat V_n=\operatorname{clip}_{[v_-,v_+]}
 \left(\frac{\widehat N}{\max\{\widehat D,a\}}\right),\\
 \sup_{\theta\in\Theta}\|\widehat V_n-V\|_{L^2}
 \le C\left[b^2+(n^{-1/2}+n^{-1}h^{-d/2})(1+W)\right].
\end{gathered}
\end{equation}
Here $C$ depends only on the model parameters and $(c_*,C_p,a)$.
In particular, $a=1/4$ always works.
\end{proposition}

\begin{proof}
Fix $\theta$, and write $\mu=\mu_\theta$. Using the assumptions and
\eqref{eq:U7-mass}, we obtain
$\sup_{\mathsf w}\E|c_{\mathcal S}(\mathsf w)|^2\le C(1+W)$ and,
by Tonelli's theorem, $\E\|\zeta_{\mathcal S}\|_{L^2(M_h)}^2\le CW$.
Then the random test functions used below belong to $L^2(M_h)$ almost
surely.

For an evaluation pair $o=(x,y),o'=(x',y')$, let $z=(y',y,1)^\top$.
We take either $\mathsf A=zz^\top-VP_0$ or $\mathsf A=P_0$, and define
$F=c_{\mathcal A}^\top\mathsf A c_{\mathcal B}$ and
$T_F=[2(m)_2]^{-1}\sum_{i\ne i'}k_h(X_i,X_{i'})F(O_i,O_{i'})$.
The two resulting statistics are $\widehat N-V\widehat D$ and
$\widehat D$, respectively. The assumption on the conditional fourth
moments and $|f|\le L_s$ imply that $\E[\|\mathsf A\|^2\mid X,X']\le C$.
Since the two pilots and the evaluation sample are independent, we
obtain
\begin{equation}\label{eq:U8-G2}
 B_F:=\iint k_hF^2\,d\mu^2,\qquad \E B_F\le C(1+W)^2.
\end{equation}
Note that this bound uses only one second moment from each pilot.
By finite mass and Cauchy--Schwarz, we also have
$\E\iint k_h|F|\,d\mu^2<\infty$. Since
$\E[zz^\top\mid X,X']=\mathbf f\mathbf f^\top+VP_0$, we have
\begin{equation}\label{eq:U8-cond}
 \E[T_F\mid\mathcal A,\mathcal B]
 =L_B:=\frac12\int c_{\mathcal A}^\top Bc_{\mathcal B}\,dM_h,
 \qquad B=\mathbf f\mathbf f^\top\ \text{or}\ P_0.
\end{equation}
Note that this calculation of the mean does not require independence
between pairs, since each ordered pair has law $\mu^2$. Averaging over the
pilots, we obtain
\begin{align}
 \E(\widehat N-V\widehat D)
 &=\frac12\int r_h^2\,dM_h\le Cb^2,\label{eq:U8-bias}\\
 \E\widehat D
 &=\frac12\int\bar c^\top P_0\bar c\,dM_h\ge\frac12.
 \label{eq:U8-denominator}
\end{align}
In the last inequality, we used that $\bar c_1=1$ and the lower bound
in \eqref{eq:U7-mass}.

Both $B$ and $\bar c$ are uniformly bounded. Expanding the
pilots, we find
\[
 L_B-\E L_B=\frac12\int
 \{\zeta_{\mathcal A}^\top B\bar c+
   \bar c^\top B\zeta_{\mathcal B}+
   \zeta_{\mathcal A}^\top B\zeta_{\mathcal B}\}\,dM_h.
\]
By the test-function bound, each linear term has $L^2$ norm at most
$C\sqrt{W/n}$. For the bilinear term, we condition on $\mathcal B$ and
apply the same bound with the independent test function
$B\zeta_{\mathcal B}$. Then the squared $L^2$ norm of the bilinear
term is at most
$(CW/n)\E\|B\zeta_{\mathcal B}\|_{L^2(M_h)}^2\le CW^2/n$.
Since $\sqrt W\le1+W$, we conclude that
\begin{equation}\label{eq:U11-pilot-fluctuation}
 \|L_B-\E L_B\|_{L^2}\le Cn^{-1/2}(1+W).
\end{equation}

For completeness, we include the order-two variance bound, which is
elementary. For a symmetric $H\in L^2(\mu^2)$, let
$H_1(o)=\int H(o,o')\mu(do')$ and $\bar H=\int H\,d\mu^2$. We decompose
$H=\bar H+h_1(o)+h_1(o')+h_2(o,o')$, where $h_1=H_1-\bar H$
and $h_2$ has mean zero in either argument. The two orders are
orthogonal; within each order, only equal sets of observation indices
contribute. Subtracting conditional means is an orthogonal projection.
The order-two Hoeffding decomposition \cite{Hoeffding1948} then gives
\begin{equation}\label{eq:U6-order-two}
 \Var\left(\frac1{(m)_2}\sum_{i\ne i'}H(O_i,O_{i'})\right)
 =\frac4m\|h_1\|_2^2+\frac2{(m)_2}\|h_2\|_2^2
 \le\frac4m\int H_1^2d\mu+\frac2{(m)_2}\int H^2d\mu^2.
\end{equation}
We now fix pilots with $B_F<\infty$, symmetrize
$\overline F(o,o')=(F(o,o')+F(o',o))/2$, and set
$H=k_h\overline F/2$. This symmetrization preserves $T_F$. By Jensen's
inequality and $k_h^2=h^{-d}k_h$, we have
$\int H^2d\mu^2\le h^{-d}B_F/4$. A weighted Cauchy--Schwarz
inequality gives
\[
 H_1(o)^2\le\frac14\left(\int k_h\,d\mu(o')\right)
                  \left(\int k_h\overline F(o,o')^2\,d\mu(o')\right),
 \qquad \int H_1^2d\mu\le p_+B_F/4.
\]
Then \eqref{eq:U6-order-two} implies that
\begin{equation}\label{eq:U11-pair-variance}
 \Var(T_F\mid\mathcal A,\mathcal B)
 \le\left(\frac{p_+}{m}+\frac1{2(m)_2h^d}\right)B_F.
\end{equation}
By Cauchy--Schwarz, we also have
$|\E[T_F\mid\mathcal A,\mathcal B]|^2\le p_+B_F/4$, and
\eqref{eq:U8-G2} then proves square integrability.
Since $(m)_2\ge m^2/2$ and $m\ge c_*n$, taking expectations yields
\begin{equation}\label{eq:U11-pair-fluctuation}
 \|T_F-L_B\|_{L^2}
 \le C(n^{-1/2}+n^{-1}h^{-d/2})(1+W).
\end{equation}
Finally, we note that clipping is $1$-Lipschitz and fixes $V$, and
that
\[
 \frac{\widehat N}{\max\{\widehat D,a\}}-V
 =\frac{\widehat N-V\widehat D-V(a-\widehat D)_+}
        {\max\{\widehat D,a\}}.
\]
Since $\E\widehat D\ge2a$, we have
$(a-\widehat D)_+\le|\widehat D-\E\widehat D|$, and it follows that
\begin{equation}\label{eq:U14-stab}
 |\widehat V_n-V|\le a^{-1}
 \{|\widehat N-V\widehat D|+v_+|\widehat D-\E\widehat D|\}.
\end{equation}
Combining this bound with the two fluctuation bounds and
\eqref{eq:U8-bias}, we obtain \eqref{eq:pair-risk}.
The reduction \eqref{eq:U14-stab} is similar to that of Shen et al.\
\cite[Appendix~A.1, equation~(28)]{ShenGaoWittenHan2020}, who control
the error of their ratio through the numerator minus the target times
the denominator, on an event where the denominator is bounded below.
\end{proof}

\begin{proof}[Proof of \cref{prop:elementary-upper} for $0<s\le1$]
For $n\ge2$, we apply \cref{prop:pair-risk} with all $m=n$
observations and the deterministic pilots $c_{\mathcal A}=c_{\mathcal B}=(1,-1,0)^\top$.
Then $W=0$, and \eqref{eq:pair-statistics} becomes
\[
 \widehat N=\frac1{2(n)_2}\sum_{i\ne i'}
 k_h(X_i,X_{i'})(Y_i-Y_{i'})^2,\qquad
 \widehat D=\frac1{(n)_2}\sum_{i\ne i'}k_h(X_i,X_{i'}).
\]
The ratio $\widehat N/\widehat D$ is a variant of the construction of
Robins et al.\ \cite[Section~4]{RobinsEtAl2008}, which selects one pair
of observations from each cube that contains at least two. Our ratio averages
half the squared differences over all pairs in common cubes, as in the
ratio estimator of Shen et al.\ \cite{ShenGaoWittenHan2020}.
Here $r_h=f(x')-f(x)$ has norm at most $Ch^s$. Since
$\E\widehat D=\int dM_h\ge1$, we may take $a=1/2$ in
\eqref{eq:pair-risk}, which gives RMS risk at most
$C(h^{2s}+n^{-1/2}+n^{-1}h^{-d/2})$ for the clipped ratio.
This statistic is Borel because its kernels are Borel and clipping
is continuous. Let $h$ be a dyadic side length such that
$n^{-2/(d+4s)}/2<h\le n^{-2/(d+4s)}$. Then the risk is at most
$C(n^{-1/2}\vee n^{-4s/(d+4s)})$, including $s=1$ and $d=4s$.
The remaining case $n=1$ is covered by the estimator $v_-$, after we
increase $C$.
\end{proof}
\chapter{A refined upper bound: resolving below the sampling scale}
\label{sec:upper-sharp}

Throughout this section, we assume that $s>1$ and that $d>4s$ is an
integer. We construct one estimator from two predictions of a regression
increment and an independent sample of nearby pairs. By resolving the
regression function below the sampling scale, we obtain the following
improvement.

\begin{theorem}\label{thm:upper-sharp}
Let $s>1$ and let $d>4s$ be an integer. Put $L=\log n$ and
\[
 \begin{gathered}
 \lambda=\frac{2(s+1)}{d+4},
 \qquad
 \tau=\log\frac{\sqrt{p_+}+\sqrt{p_-}}{\sqrt{p_+}-\sqrt{p_-}},
 \\
 \kappa=4\sqrt{\frac{2\tau(s-1)(d-4s)}{(d+4)^3}},
 \qquad
 q_s=\binom{d+\lceil s\rceil-1}{d}.
 \end{gathered}
\]
There exist $C<\infty$ and $n_0\in\mathbb N$, depending only on
$(s,d,p_-,p_+,L_s,v_-,v_+,C_4)$, such that for every integer
$n\ge n_0$
\begin{equation}\label{eq:upper-sharp}
 \mathfrak r_n
 \le C\,n^{-\lambda}e^{-\kappa\sqrt L}\,
 L^{\frac{s-1}{d+4}+\frac{d(q_s-1/2)}{d+4}} .
\end{equation}
\end{theorem}

Recalling that
$\Psi_n=n^{-\lambda}e^{-\kappa\sqrt L}L^{(s-1)/(d+4)}$ and
$\Gamma=d(q_s-\tfrac12)/(d+4)$, we see that the bound
\eqref{eq:upper-sharp} reads $\mathfrak r_n\le C\,\Psi_nL^{\Gamma}$.
The constants do not depend on the extension domain $\mathcal U$; the
only property of the domain used in bounding the Taylor remainder is the
convexity of $[0,1]^d$. We complete the proof in
\cref{sec:U-final}.

\section{Setting and fixed quantities}\label{sec:U-setup}

For pairs in one cube of side $h$, the estimator uses products of
two corrected response differences, with corrections given by
independent predictions of the increment $f(x')-f(x)$.
Resolving the regression function at $K$ cells yields a bias of order
$h^2K^{-2(s-1)/d}$ and a pair error $n^{-1}h^{-d/2}$. When $K=ne^T$,
balancing the two improves the rate $n^{-\lambda}$ at the sampling
scale by the factor $e^{-2(s-1)T/(d+4)}$. We achieve
$T\asymp\sqrt{\log n}$ using unbiased polynomials of local moments,
whose degrees are chosen to control their factorial variances.
Kong and Valiant \cite[Section~2.2]{KongValiant2018} approximate a
scalar functional of an unknown covariance matrix by a polynomial in its
powers, and they estimate the resulting quadratic forms without bias by
sums over distinct observations. Chen, Liu and Mukherjee
\cite[Section~3.2]{ChenLiuMukherjee2024} approximate quadratic forms in
an inverse covariance matrix by Chebyshev polynomials and estimate the
resulting quadratic forms in matrix powers by U-statistics. The higher-order influence function
estimators of Robins et al.\ \cite[Theorems~3.14--3.15]{RobinsEtAl2008}
use inverse Gram matrices computed under initial nuisance estimates,
together with U-statistics of increasing order, to correct the bias of
an estimator of a functional.

We write $|\cdot|$ for the Euclidean norm of a vector and
$\norm{\cdot}$ for the operator norm of a matrix; the matrix blocks of a
moment tuple are measured in the Frobenius norm. Let
$\ell=\lceil s\rceil-1$ and
\begin{equation}\label{eq:U-setup-dims}
 q_s=\binom{d+\ell}{d},\qquad r=q_s-1,\qquad L=\log n.
\end{equation}
The constants for the approximation of the inverse and for the
allocation are
\begin{equation}\label{eq:U-setup-rho}
 \tau=\log\frac{\sqrt{p_+}+\sqrt{p_-}}{\sqrt{p_+}-\sqrt{p_-}},
 \qquad \rho=e^{-\tau}\in(0,1),
\end{equation}
\[
 \vartheta=(s-1)/d,\qquad \beta=\vartheta/\tau,
\]
\begin{equation}\label{eq:U-setup-exponents}
 \varpi=\frac{d-4s}{d(d+4)},\qquad
 \eta=\frac{2(d+4s)}{d(d+4)},
\end{equation}
\begin{equation}\label{eq:U-setup-allocation}
 a_*=\sqrt{8\varpi/\beta}.
\end{equation}
In particular, we have
\begin{equation}\label{eq:U-setup-rate}
 \lambda=\frac{2(s+1)}{d+4},\qquad
 \kappa=\frac{2(s-1)}{d+4}a_*
 =4\sqrt{\frac{2\tau(s-1)(d-4s)}{(d+4)^3}}.
\end{equation}
Below, we write $C,c$ for fixed positive constants. They depend only on
$\mathfrak p=(s,d,p_-,p_+,L_s,v_+,C_4)$, and their values may change
between lines. Together with the model moment bounds, the feature constant $C_\psi$
of \cref{lem:U1} and the analytic constants $M_1,C_1$ of \cref{lem:U5}
determine the variance constant
$C_0\ge1$, which is fixed sufficiently large in \cref{lem:U7}.

\subsection*{Model bounds}
Let $\mu$ be the law of $O=(X,Y)$, and let $\varepsilon=Y-f(X)$.
By the assumptions of the model, we have
\begin{equation}\label{eq:U-M0}
 \sup_{[0,1]^d}|f|\le L_s,
\end{equation}
\begin{equation}\label{eq:U-M2}
 \E[(1+|Y|)^2\mid X]\le M_2:=2(1+L_s^2+v_+).
\end{equation}
For the second bound, note that $\E[Y^2\mid X]=f(X)^2+V$ and
$(1+|Y|)^2\le2(1+Y^2)$. By Taylor's formula, for all
$x,z\in[0,1]^d$, we have
\begin{equation}\label{eq:U-Tay}
 \left|f(z)-\sum_{|\gamma|\le\ell}
       \frac{\partial^\gamma f(x)}{\gamma!}(z-x)^\gamma\right|
 \le C_{\rm Tay}|z-x|^s,\qquad
 C_{\rm Tay}=d^\ell L_s/\ell!.
\end{equation}
To verify the endpoint convention for integer $s$, we apply the
integral form of Taylor's formula in one dimension to
$F(t)=\widetilde f(x+t(z-x))$. The remainder after degree $\ell$ is
\[
 \int_0^1\frac{(1-t)^{\ell-1}}{(\ell-1)!}
       \{F^{(\ell)}(t)-F^{(\ell)}(0)\}\,dt.
\]
By the H\"older bound and the identity
$\sum_{|\gamma|=\ell}\ell!/\gamma!=d^\ell$, its integrand is bounded
by $d^\ell L_s|z-x|^s$, and the weight has mass $1/\ell!$.
Since the segment stays in the convex cube, the extension domain does
not enter these constants. Finally, since $V^2\le C_4$, we have
\begin{equation}\label{eq:U-setup-IV}
 V\in I_V:=[v_-,\min\{v_+,\sqrt{C_4}\}].
\end{equation}

\subsection*{Cells and fixed samples}
For $j\ge0$, let $\mathcal D_j$ be the set of half-open dyadic cubes
of side $2^{-j}$ in $[0,1)^d$, let $z_I$ be a cell's corner, and set
$K_j=2^{jd}=|I|^{-1}$. We denote the cell containing $x$ by $I_j(x)$.
Any two cells are either disjoint or nested, and
\begin{equation}\label{eq:U-setup-cellprob}
 p_-K_j^{-1}\le\Pp(X\in I)\le p_+K_j^{-1}
 \qquad(I\in\mathcal D_j).
\end{equation}
We divide the first $3n_b$ observations into independent blocks
$\mathcal A,\mathcal B,\mathcal P$, where
\begin{equation}\label{eq:U-setup-Lambda}
 n_b=\lfloor n/3\rfloor,\qquad \Lambda=n/6.
\end{equation}
The first two blocks are used for the pilots, and the third for the
pairs; the remaining observations are discarded. For a block
$\mathcal S$ of size $m$, we define the factorial sums
\begin{equation}\label{eq:U-setup-Sigma}
 \Sigma_k(\phi;\mathcal S)
 =\sum_{\substack{1\le i_1,\ldots,i_k\le m\\
                    i_1,\ldots,i_k\text{ distinct}}}
             \phi(O_{i_1},\ldots,O_{i_k}),
 \qquad (m)_k=m!/(m-k)!.
\end{equation}
The sum is zero when $k>m$, and $\Sigma_0(c;\mathcal S)=c$. The
admissibility condition below ensures that the denominators are
positive.

\section{Anchored monomials and local fits}\label{sec:U-anchored}

We index coordinates by
$\mathcal I=\{\gamma\in\mathbb N_0^d:1\le|\gamma|\le\ell\}$, taken in a
fixed order, and let $q(v)=(v^\gamma)_{\gamma\in\mathcal I}$.
For $I\in\mathcal D_j$ and $x\in I$, we define
\begin{equation}\label{eq:Udef-features}
 \psi_{I,x}(z)=q(2^j(z-x)),\qquad
 B_I(x)=\int_{[0,1]^d}q(w-u)q(w-u)^\top dw,
 \quad u=2^j(x-z_I).
\end{equation}
We abbreviate $\psi_{j,x}=\psi_{I_j(x),x}$ and
$B_j(x)=B_{I_j(x)}(x)$. Note that the matrix $B_I(x)$ is known; its
$(\gamma,\delta)$ entry equals
\[
 \prod_{i=1}^d
 \frac{(1-u_i)^{\gamma_i+\delta_i+1}
                 -(-u_i)^{\gamma_i+\delta_i+1}}
      {\gamma_i+\delta_i+1}.
\]
The coordinate transport is the fixed diagonal matrix
\begin{equation}\label{eq:U-setup-transport}
 \mathcal T_j=\operatorname{diag}_{\gamma\in\mathcal I}(2^{-|\gamma|})
 \quad(j\ge1),\qquad \mathcal T_0=0.
\end{equation}

In part~(3) of the next lemma, we normalize the local Gram matrix by
its Lebesgue counterpart $B_I(x)$, so that the density bounds place its
spectrum in $[p_-,p_+]$. Park~\cite[Section~3.2, equation~(9)]{Park2026} obtains the
corresponding spectral interval with a block basis orthonormal for
Lebesgue measure. Our companion paper with
Kallus~\cite[Section~4.2, proof of Lemma~4.3]{AronowKallusLopatto2026}
uses a basis orthonormal for normalized Lebesgue measure on each cell
and includes the reciprocal cell volume in its Gram matrix.
\begin{lemma}\label{lem:U1}
There are fixed $c_B,C_B>0$ and $C_\psi\ge1$, depending only on
$(d,s)$, such that the following holds at all cells and anchors.
\begin{enumerate}
\item[\textup{(1)}] The components of $\psi_{I,x}$ form a basis of
the degree-at-most-$\ell$ polynomials vanishing at $x$, and
\begin{equation}\label{eq:U1-Gram}
 K_j\int_I\psi_{I,x}\psi_{I,x}^\top dz=B_I(x),
 \qquad c_BI_r\preceq B_I(x)\preceq C_BI_r.
\end{equation}
\item[\textup{(2)}] We have $\sqrt r,\sqrt{C_B/c_B}\le C_\psi$,
$\sup_I|\psi_{I,x}|\le C_\psi$,
\begin{equation}\label{eq:U1-lipschitz}
 |\psi_{I,x}(x')|\le C_\psi2^j|x'-x|\quad(x'\in I),
\end{equation}
and on $I$,
$|B_I(x)^{-1}\psi_{I,x}|,
 \|B_I(x)^{-1}\psi_{I,x}\psi_{I,x}^\top\|_F\le C_\psi^2$.
\item[\textup{(3)}] If $p_-\le p\le p_+$ almost everywhere on $I$ and
$H=K_j\int_Ip\psi_{I,x}\psi_{I,x}^\top dz$, then
\begin{equation}\label{eq:U1-spectrum}
 p_-B_I(x)\preceq H\preceq p_+B_I(x).
\end{equation}
Thus $B_I(x)^{-1}H$ is similar to a symmetric matrix with spectrum
in $[p_-,p_+]$, with similarity condition number at most
$\chi:=\sqrt{C_B/c_B}$.
\item[\textup{(4)}] On all of $\mathbb R^d$, for $j\ge1$,
\begin{equation}\label{eq:U1-transport}
 \psi_{j-1,x}=\mathcal T_j^\top\psi_{j,x},\qquad
 \norm{\mathcal T_j}\le1/2.
\end{equation}
\item[\textup{(5)}] The features and known inverse $B_I(x)^{-1}$
are continuous in their arguments on each fixed cell and globally
Borel; the transports are constant.
\end{enumerate}
\end{lemma}

\begin{proof}
The centered monomials form a basis of the space of anchored
polynomials, and the Gram identity follows from a change of variables
on the cell. If a coefficient vector is nonzero, then the corresponding
polynomial is nonzero, and its squared integral is positive. Since $u$
ranges over the compact cube $[0,1]^d$, the constants $c_B,C_B$ can be
chosen uniformly by continuity, and $B_I(x)^{-1}$ is continuous.
Part (2) holds after we increase the single fixed constant $C_\psi$,
since the monomials, their gradients and this inverse are bounded. Part (3) follows from the density bounds: writing $B=B_I(x)$, we see
that
$S=B^{-1/2}HB^{-1/2}$ has spectrum in $[p_-,p_+]$ and that
$B^{-1}H=B^{-1/2}SB^{1/2}$. Scaling each centered monomial gives (4),
because $|\gamma|\ge1$. Finally, (5) follows from the
explicit formulas on each cell and the fact that the cells of each
level form a finite Borel partition.
\end{proof}

\subsection*{Preconditioned moments and the anchored fit}
For a parameter $\theta=(p,f,V,Q)\in\Theta$, we define
\begin{equation}\label{eq:Udef-moments}
\begin{aligned}
 H_j(x)&=K_j\int_{I_j(x)}p(z)\psi_{j,x}(z)\psi_{j,x}(z)^\top dz,
 &G_j(x)&=B_j(x)^{-1}H_j(x),\\
 \beta_j(x)&=B_j(x)^{-1}K_j\int_{I_j(x)}p(z)\psi_{j,x}(z)f(z)dz,
 &\mathrm m_j(x)&=B_j(x)^{-1}K_j\int_{I_j(x)}p(z)\psi_{j,x}(z)dz.
\end{aligned}
\end{equation}
These quantities are bounded uniformly in the level and the anchor,
and they are Borel by dominated convergence on each cell. Using the
similarity in \cref{lem:U1}, we obtain
\begin{equation}\label{eq:U-setup-Gspec}
 \norm{G_j}\le\chi p_+,\qquad
 \norm{G_j^{-1}}\le\chi p_-^{-1},\qquad
 p_-^r\le\det G_j\le p_+^r.
\end{equation}
For a formal anchor value $y\in\mathbb R$, we set
\begin{equation}\label{eq:Udef-b}
 b_j(x;y)=\beta_j(x)-y\mathrm m_j(x)
 =B_j(x)^{-1}K_j\int_{I_j(x)}p(z)\psi_{j,x}(z)(f(z)-y)dz.
\end{equation}
At the root, the parent values are $G_{-1}=I_r$ and
$\beta_{-1}=\mathrm m_{-1}=0$. Since $p_-<1<p_+$, the spectrum of this
parent lies in the same interval. We define the moment tuples by
\begin{equation}\label{eq:Udef-theta}
 \theta_j(x)=(G_j,G_{j-1},\beta_j,\beta_{j-1},
                       \mathrm m_j,\mathrm m_{j-1})(x)\in E.
\end{equation}
The fixed finite-dimensional space $E$ carries the Frobenius norm on
its matrix coordinates and the Euclidean norm on its vector
coordinates. The population fit and the increments are defined by
\begin{equation}\label{eq:Udef-fit}
 a_j(x)=G_j(x)^{-1}b_j(x;f(x)),\qquad a_{-1}=0,\qquad
 \delta_j(x)=a_j(x)-\mathcal T_ja_{j-1}(x).
\end{equation}
We note that the known preconditioner cancels in the fit:
$a_j=H_j^{-1}K_j\int_{I_j(x)}p\psi_{j,x}(f-f(x))$.

\begin{lemma}\label{lem:U2}
Uniformly over parameters, levels and anchors, the following holds.
\emph{(1)} For $x'\in I_j(x)$,
\begin{equation}\label{eq:U2-approx}
 |f(x')-f(x)-\psi_{j,x}(x')^\top a_j(x)|
 \le C|x'-x|2^{-j(s-1)},
\end{equation}
\emph{(2)} The increments satisfy
\begin{equation}\label{eq:U2-increment}
 |\delta_j(x)|\le C2^{-js}.
\end{equation}
\emph{(3)} For every $J\ge0$ and $x'\in\mathbb R^d$,
\begin{equation}\label{eq:U2-telescope}
 \psi_{J,x}(x')^\top a_J(x)
 =\sum_{j=0}^J\psi_{j,x}(x')^\top\delta_j(x).
\end{equation}
The constants depend only on $(d,s,p_-,p_+,L_s)$.
\end{lemma}

\begin{proof}
The nonconstant part of the Taylor polynomial is
$P_x=\psi_{j,x}^\top c_j$, where
$c_{j,\gamma}=2^{-j|\gamma|}\partial^\gamma f(x)/\gamma!$.
Then $c_j=\mathcal T_jc_{j-1}$ for $j\ge1$. Let
$\varrho_x=f-f(x)-P_x$. By the normal equations, we have
\begin{equation}\label{eq:U2-ej}
 a_j=c_j+e_j,\qquad
 e_j=H_j^{-1}K_j\int_{I_j(x)}p(z)\psi_{j,x}(z)\varrho_x(z)dz.
\end{equation}
Then \eqref{eq:U-Tay} implies that $|\varrho_x(z)|\le C|z-x|^s$.
Combining this bound with $\norm{H_j^{-1}}\le(p_-c_B)^{-1}$ and the
boundedness of the features, we find
\begin{equation}\label{eq:U2-ej-bound}
 |e_j|\le C2^{-js}.
\end{equation}
The error of the fit equals $\varrho_x(x')-\psi_{j,x}(x')^\top e_j$.
Since $|x'-x|\le\sqrt d\,2^{-j}$, the Taylor bound, together with the
Lipschitz bound on the features, implies \eqref{eq:U2-approx}.
For $j\ge1$, we have $\delta_j=e_j-\mathcal T_je_{j-1}$, which proves
\eqref{eq:U2-increment}. At the root, the increment bound follows from
$|f-f(x)|\le2L_s$ and the bound on the inverse Gram matrix.
Finally, by \eqref{eq:U1-transport}, we have
$\psi_{j,x}^\top\delta_j=\psi_{j,x}^\top a_j-
\psi_{j-1,x}^\top a_{j-1}$, and the telescoping identity follows by
summation.
\end{proof}
\section{A polynomial estimator of a fitted increment}\label{sec:U-inverse}

Recall that the fits at consecutive levels reproduce the same Taylor
polynomial. To preserve this cancellation, we clear the determinant of
the parent before approximating the remaining inverse. Our companion
paper uses the same device for products of two projection
residuals: it clears the determinants of the parent Gram matrices and
approximates the remaining inverse and reciprocal determinants by one
joint series \cite[Sections~4.1--4.2]{AronowKallusLopatto2026}. Park
\cite[Section~3.3]{Park2026} uses matrix polynomials that annihilate
the polynomial trends of the parent exactly, in a multiresolution
estimator of the expected conditional covariance.
Let
$t(z)=(2z-p_--p_+)/(p_+-p_-)$ and $\kappa_0=(p_-p_+)^{-1/2}$.
For a free square matrix $M$, we define the formal series
\[
 g_\zeta(M)=\kappa_0(1-\rho^2\zeta^2)
     \{I+2\rho\zeta t(M)+\rho^2\zeta^2I\}^{-1},
\]
whose denominator has constant term $I$. Let
\begin{equation}\label{eq:Udef-Pm}
 \mathsf P_m(A,C)=\sum_{k=0}^m[\zeta^k]
                       \{g_\zeta(A)\det g_\zeta(C)\}.
\end{equation}
Its degree in the matrix entries is at most $m$, since the recurrence
for the inverse produces coefficients whose degree is at most their
order in the series, and the determinant and the product preserve this
bound.

For integers $m_j\ge0$, we set $D_j=m_j+q_s$. Writing
$G=G_j(x)$ and $G^-=G_{j-1}(x)$, we define
\begin{equation}\label{eq:Udef-increment}
\begin{aligned}
 \mathcal N_j(x;y)&=\det(G^-)b_j(x;y)
       -G\mathcal T_j\operatorname{adj}(G^-)b_{j-1}(x;y),\\
 F_j(x;y)&=\mathsf P_{m_j}(G,G^-)\mathcal N_j(x;y).
\end{aligned}
\end{equation}
Here $\operatorname{adj}(C)$ is the adjugate, which satisfies
$\operatorname{adj}(C)C=\det(C)I$, and $\operatorname{adj}(C)=1$ for
$r=1$. The root convention gives
$\mathcal N_0=b_0$ and $F_0=\mathsf P_{m_0}(G_0,I)b_0$.
We regard the right sides as polynomials $F_{j,x}(\theta;y)$ in the
free moment tuple $\theta\in E$. Their
coefficients are known and, since $b=\beta-y\mathrm m$, affine in the
formal variable $y$. For $x'\in I_J(x)$, the polynomial to be lifted
is
\begin{equation}\label{eq:U-affine-polynomial}
 \mathcal P_{x,x'}(\boldsymbol\theta;y)
 =\sum_{j=0}^J\psi_{j,x}(x')^\top F_{j,x}(\theta_j;y),
 \qquad\boldsymbol\theta=(\theta_0,\ldots,\theta_J).
\end{equation}
We write $D^k\Phi(\theta)[\gamma_1,\ldots,\gamma_k]$ for the
mixed derivative of $\Phi(\theta+\sum_i t_i\gamma_i)$ at $t=0$, with
the convention $D^0\Phi=\Phi$.

\begin{lemma}[The increment polynomial]\label{lem:U5}
There are fixed $C_{\rm inc},M_1,C_1>0$, depending only on
$(d,s,p_-,p_+,L_s)$, with the following properties at every population
tuple, level and anchor, for arbitrary nonnegative degrees $m_j$.

\emph{(1)} For every real $y$,
\begin{equation}\label{eq:U5-identity}
 G^{-1}\det(G^-)^{-1}\mathcal N_j(x;y)
 =G^{-1}b_j(x;y)-\mathcal T_j(G^-)^{-1}b_{j-1}(x;y),
\end{equation}
and hence
\begin{equation}\label{eq:U5-N}
 |\mathcal N_j(x;f(x))|
 \le\chi p_+^{r+1}|\delta_j(x)|\le C_{\rm inc}2^{-js}.
\end{equation}
\emph{(2)} The approximation error satisfies:
\begin{equation}\label{eq:U5-error}
 |F_j(x;f(x))-\delta_j(x)|
 \le C_{\rm inc}2^{-js}(m_j+1)^r\rho^{m_j}.
\end{equation}
\emph{(3)} The polynomial $F_{j,x}(\theta;y)$ has degree at most
$D_j=m_j+r+1$. Every monomial contains one vector-moment coordinate,
to first power, and has coefficient constant or linear in $y$.
Thus each observation's response enters its factorial lift at most
to first power.

\emph{(4)} Uniformly for $|y_0|\le L_s$,
$\Phi\in\{F_{j,x}(\cdot;0),-\partial_yF_{j,x},
F_{j,x}(\cdot;y_0)\}$, $k\ge0$ and all directions in $E$,
\begin{equation}\label{eq:U5-deriv}
 |D^k\Phi(\theta_j(x))[\gamma_1,\ldots,\gamma_k]|
 \le M_1 k!C_1^k\prod_i|\gamma_i|.
\end{equation}
The constants $M_1,C_1$ depend only on
$(r,p_-,p_+,C_\psi,L_s)$, not on the level or degrees.
\end{lemma}

\begin{proof}
For $z\in[p_-,p_+]$, let $\phi$ be such that $t(z)=\cos\phi$.
Summing two geometric series, we obtain the Chebyshev identity
\[
 g_\zeta(z)=\kappa_0\{1+2\sum_{k\ge1}(-\rho\zeta)^k\cos(k\phi)\},
 \qquad |\zeta|<\rho^{-1}.
\]
Its coefficients have modulus at most $2\kappa_0\rho^k$, and a direct
substitution of $\rho$ shows that $g_1(z)=z^{-1}$.
The identity is the weighted generating function of the Chebyshev
polynomials \cite[Eq.~18.12.7]{DLMF}; for generating functions, see also
\cite[Section~5.2.1]{MasonHandscomb2003} and
\cite[Section~1.5]{Rivlin1990}. At $\zeta=1$, after an affine change of
variables, it is the classical Chebyshev expansion of the reciprocal
\cite[equation~(5.14)]{MasonHandscomb2003},
\cite[equation~(2.4)]{Mathar2006}.
By \cref{lem:U1}(3), each population Gram matrix is similar to a
symmetric matrix with spectrum in this interval, through a similarity
with condition number at most $\chi$. The bound on its matrix
coefficients then gains only the fixed factor $\chi$.
Since the determinant is invariant under similarity, its $r$ scalar
factors satisfy the original bounds. Convolving these bounds, we find
\[
 \|[\zeta^k]\{g_\zeta(A)\det g_\zeta(C)\}\|
 \le\chi(2\kappa_0)^{r+1}\binom{k+r}{r}\rho^k.
\]
For these population matrices, the series at one equals
$A^{-1}\det(C)^{-1}$. Using
$\binom{m+1+i+r}{r}\le(r+1)(m+1)^r\binom{i+r}{r}$ and
$\sum_{i\ge0}\binom{i+r}{r}\rho^i=(1-\rho)^{-r-1}$, we obtain
\begin{equation}\label{eq:U4-approx}
 \norm{\mathsf P_m(A,C)-A^{-1}\det(C)^{-1}}
 \le C(m+1)^r\rho^m.
\end{equation}

We obtain \eqref{eq:U5-identity} from the adjugate identity, and at
$y=f(x)$ this identity gives $\mathcal N_j=\det(G^-)G\delta_j$.
Then \eqref{eq:U-setup-Gspec}, together with \eqref{eq:U2-increment},
implies (1).
For (2), we multiply the bound on the small numerator by
\eqref{eq:U4-approx}. Each monomial of the numerator contains at most $r$
matrix coordinates and one vector coordinate. Since $\mathsf P_m$
contains only matrix coordinates and has degree at most $m$, this
proves (3), including the substitution at the root.

For (4), let $R=(1+\rho^{-1})/2$, and consider the set
\[
 \mathcal G=\{B^{-1}H:\ B=B^\top,\ H=H^\top,\quad
 I_r\preceq B\preceq C_\psi^2I_r,\quad
 p_-B\preceq H\preceq p_+B\}
\]
This set is compact. To see that it contains every population Gram
matrix, we divide the known Lebesgue Gram and its density Gram by
$c_B$. It also contains the root parent $I_r$. Since its members are
similar to symmetric matrices with spectrum in $[p_-,p_+]$, the matrix
$I_r+2\rho\zeta t(A)+\rho^2\zeta^2I_r$ is invertible on
$\mathcal G\times\{|\zeta|\le R\}$.

Recall that the vector moments are uniformly bounded and that
$\|\mathcal T_j\|\le1/2$, including $\mathcal T_0=0$.
By compactness and the continuity of inversion, there exist fixed
$\epsilon,M>0$ such that the $y=0$, $y=y_0$ and $-\partial_y$
versions of $g_\zeta(G)\det g_\zeta(G^-)\mathcal N_j(\theta;y)$
are bounded by $M$
for $|\zeta|\le R$, $|\theta-\theta_j(x)|\le\epsilon$ and
$|y_0|\le L_s$, and are holomorphic on a neighborhood of this set.
We choose $\epsilon$ small enough that the closed neighborhood lies
in this holomorphic domain. Within this neighborhood, perturbations of the matrices may be
complex and need not be symmetric.

By Cauchy's formula, each $\zeta$ coefficient is bounded by
$MR^{-l}$, where $l$ is its order in the series. Summing the geometric series, we find that all partial sums
are bounded by $M'=M/(1-R^{-1})$, independently of $m_j$.
For nonzero directions, the Cauchy estimate in each direction on
$|t_i|=\epsilon/(k|\gamma_i|)$ gives
\[
 |D^k\Phi(\theta_j)[\gamma_1,\ldots,\gamma_k]|
 \le M'(k/\epsilon)^k\prod_i|\gamma_i|
 \le M'k!(e/\epsilon)^k\prod_i|\gamma_i|.
\]
In the last inequality, we used the fact that $k!\ge(k/e)^k$. The cases of
zero directions and $k=0$ are immediate. All compact sets and bounds
above depend only on $(r,p_-,p_+,C_\psi,L_s)$, and this proves (4).
\end{proof}
\section{Unbiased polynomials from a fixed sample}\label{sec:U-factorial}

Let $\mathcal S=(O_1,\dots,O_{n_b})$ consist of independent observations
with common law $\mu$. A polynomial
$\Phi:\mathbb R^N\to\mathbb R^r$ of degree $D\le n_b$ has a unique
homogeneous expansion
$\Phi(\theta)=\sum_{k=0}^D\Phi_{(k)}[\theta,\dots,\theta]$,
where $\Phi_{(k)}$ is symmetric and $k$-linear. Using the factorial sums
of \eqref{eq:U-setup-Sigma}, we define its unbiased lift by
\begin{equation}\label{eq:U6-hat}
 \widehat\Phi(\mathcal S;g)
 =\sum_{k=0}^D\frac{\Sigma_k(
       \Phi_{(k)}[g(\cdot),\dots,g(\cdot)];\mathcal S)}{(n_b)_k}.
\end{equation}
The lift is linear in $\Phi$. If the coefficients of the polynomial are
affine in a formal variable $y$, we lift the coefficients and leave $y$
unchanged. For the indicator of a cell and a monomial of degree $k$,
the lift is the falling factorial of the cell count divided by
$(n_b)_k$, the unbiased estimator of a power of a cell probability used
in \cite[equation~(21)]{JiaoEtAl2015}.

To compute the variance, we expand around the mean
$\theta=\int g\,d\mu$. This mean enters the analysis, while
\eqref{eq:U6-hat} is computed from the known coefficients of the
polynomial. For scalar kernels on $\mathbb O^k$, let $E_i$ denote
integration of the $i$th argument against $\mu$, and let
$\Pi_k=\prod_{i=1}^k(I-E_i)$. Since the operators $E_i$ commute and are
orthogonal projections in $L^2(\mu^k)$, we have
$\|\Pi_k H\|_2\le\|H\|_2$.

\begin{lemma}[Centered Taylor identity]\label{lem:U6}
Let $g\in L^2(\mu;\mathbb R^N)$, $\theta=\int g\,d\mu$ and
$Z(o)=g(o)-\theta$. The lift is square integrable, has mean
$\Phi(\theta)$, and satisfies
\begin{equation}\label{eq:U6-centered}
 \widehat\Phi-\Phi(\theta)
 =\sum_{k=1}^D\frac{1}{k!(n_b)_k}
 \Sigma_k\bigl(D^k\Phi(\theta)[Z(\cdot),\dots,Z(\cdot)];\mathcal S\bigr).
\end{equation}
For $a\in\mathbb R^r$, put
\begin{equation}\label{eq:U6-poly-kernel}
 \mathcal K_k(o_1,\dots,o_k)
 =a\cdot D^k\Phi(\theta)[g(o_1),\dots,g(o_k)].
\end{equation}
If $0<\Lambda\le n_b-D+1$, then
\begin{equation}\label{eq:U6-comparison}
 \Var(a\cdot\widehat\Phi)
 =\sum_{k=1}^D\frac{\|\Pi_k\mathcal K_k\|_{L^2(\mu^k)}^2}{k!(n_b)_k}
 \le\sum_{k=1}^D\frac{\|\mathcal K_k\|_{L^2(\mu^k)}^2}{k!\Lambda^k}.
\end{equation}
For a measurable field of scalar lifts, write $Z(w)$ for the centered
lift and $\mathcal K_{w,k}$ for its raw kernels, padding with zero up
to a common degree $D$. Let $M$ be a finite measure and suppose these
kernels have uniformly bounded $L^2(\mu^k)$ norms. For
$\nu\in L^2(M)$, if a finite partition into blocks $B$ has disjoint
observation supports for raw kernels from distinct blocks at every
order $k\ge1$, then
\begin{equation}\label{eq:U6-localized}
\begin{aligned}
 \left\|\int\nu(w)Z(w)\,dM(w)\right\|_{L^2}^2
 &\le\sum_{k=1}^D\frac{\left\|\int\nu(w)\mathcal K_{w,k}\,dM(w)\right\|_2^2}
                         {k!\Lambda^k}\\
 &\le\max_B M(B)\int\nu(w)^2
       \sum_{k=1}^D\frac{\|\mathcal K_{w,k}\|_2^2}{k!\Lambda^k}\,dM(w).
\end{aligned}
\end{equation}
The first inequality also holds without the support assumption.
\end{lemma}

\begin{proof}
By multilinearity and independence, every summand of the lift is
integrable, and the lift has the claimed mean. Each summand is square
integrable because $g\in L^2$. We expand each homogeneous term using $g=\theta+Z$.
For a term of degree $l$ and a choice of $k$ centered arguments,
summing over the unused observation indices gives the factor
$(n_b-k)_{l-k}$. Since
$(n_b)_l=(n_b)_k(n_b-k)_{l-k}$, its coefficient after averaging is
$\binom lk\Phi_{(l)}[Z,\dots,Z,\theta,\dots,\theta]$.
Summing over $l\ge k$, we obtain $D^k\Phi(\theta)[Z,\dots,Z]/k!$,
which proves \eqref{eq:U6-centered}.

Observe that each centered kernel integrates to zero in any one of
its arguments. Two summands in \eqref{eq:U6-centered} are therefore orthogonal
unless their tuples contain exactly the same observation indices,
since an index that appears only once can be integrated first.
At order $k$, there are $(n_b)_k$ choices of the first tuple and $k!$
reorderings for the second. Further, we have
$\Pi_k\mathcal K_k=a\cdot D^k\Phi(\theta)[Z,\dots,Z]$,
and this proves the variance identity. The inequality follows from the
fact that $\Pi_k$ is a contraction and from
$(n_b)_k\ge(n_b-D+1)^k\ge\Lambda^k$.

For the integrated version, since $M$ is finite and the kernel bounds
are uniform, we have $\int|\nu|\|\mathcal K_{w,k}\|_2\,dM<\infty$.
By Minkowski's inequality and Fubini's theorem, we may then integrate
the centered Taylor identity and commute the integral with $\Pi_k$. We
apply the variance identity and the contraction property of the
projections to the integrated kernels. At each order, the block
integrals are orthogonal because the raw supports are disjoint, and
Cauchy--Schwarz on each block gives
\[
 \left\|\int\nu\mathcal K_{w,k}\,dM\right\|_2^2
 =\sum_B\left\|\int_B\nu\mathcal K_{w,k}\,dM\right\|_2^2
 \le\sum_B M(B)\int_B\nu^2\|\mathcal K_{w,k}\|_2^2\,dM.
\]
Summing over $k$ completes the proof of \eqref{eq:U6-localized}. We
note that the localization takes place before centering; the actual
covariances between different blocks for the fixed sample need not
vanish.
\end{proof}

Since each centered kernel integrates to zero in any one of its
arguments, the centered Taylor identity is the Hoeffding decomposition
\cite{Hoeffding1948} of the lift. Our companion paper defines the same
lift and proves the same covariance identity
\cite[Section~4.3, Lemma~4.4]{AronowKallusLopatto2026}.
A variance identity of the same form, for polynomial combinations of
complete U-statistics estimating trace moments, appears in
\cite[Theorem~6.2]{Song2026}.
\section{Definition of the estimator}\label{sec:U-estimator}

The estimator uses the independent pilot and evaluation blocks and the
comparison parameter $\Lambda$ of \eqref{eq:U-setup-Lambda}.

Its parameters are a final level $J\ge0$, degrees
$m_0,\dots,m_J\ge0$ and a pair level $J_h\ge J$, all of which are
integers. We set
\[
  h=2^{-J_h},\quad K_*=K_J=2^{dJ},\quad D_j=m_j+q_s\ (0\le j\le J).
\]
We say that the parameters are \emph{admissible} if
\begin{equation}\label{eq:Udef-admissible}
 \begin{gathered}
 J\ge0,\quad m_j\ge0\ (0\le j\le J),\quad J_h\ge J,\\
 \max_{0\le j\le J}D_j\le n_b-\Lambda+1,\qquad n_b-1\ge\Lambda.
 \end{gathered}
\end{equation}
The condition on the degrees ensures that the unbiased polynomial
estimates below are defined and that their variances satisfy
\cref{lem:U6}. Here $\widehat V_n$ denotes the estimator with whichever
admissible parameters have been specified. The bounds in
\cref{sec:U-variance,sec:U-equation} apply to every such choice;
the allocation below yields the bound in \cref{thm:upper-sharp}.

\subsection*{Raw features and pilots}
For an observation $o=(z,y_o)$, we define the raw features
\begin{equation}\label{eq:Udef-rawfeatures}
 (\mathcal G_{j,x},\beta^{\rm f}_{j,x},\mathrm m^{\rm f}_{j,x})(o)
 =K_j\1_{I_j(x)}(z)
   \bigl(B_j(x)^{-1}\psi_{j,x}(z)\psi_{j,x}(z)^\top,
         B_j(x)^{-1}\psi_{j,x}(z)y_o,
         B_j(x)^{-1}\psi_{j,x}(z)\bigr).
\end{equation}
Since $\E[Y\mid X]=f(X)$, their means are
$(G_j,\beta_j,\mathrm m_j)(x)$. For the fixed parent of the root, we use
the features
$(\mathcal G_{-1,x},\beta^{\rm f}_{-1,x},\mathrm m^{\rm f}_{-1,x})(o)
=(I_r,0,0)\1_{[0,1)^d}(z)$. These have the required means, because
$X\in[0,1)^d$ almost surely. We collect the current and parent features
in the same order as the moment tuple:
\begin{equation}\label{eq:Udef-g}
 g_{j,x}=(\mathcal G_{j,x},\mathcal G_{j-1,x},
 \beta^{\rm f}_{j,x},\beta^{\rm f}_{j-1,x},
 \mathrm m^{\rm f}_{j,x},\mathrm m^{\rm f}_{j-1,x}),\qquad j\ge0.
\end{equation}
By \eqref{eq:U-M2} and \cref{lem:U1}, these known features are square
integrable. Further, $\E g_{j,x}=\theta_j(x)$.
Let $g_x=(g_{0,x},\ldots,g_{J,x})$.
For each independent pilot block
$\mathcal S\in\{\mathcal A,\mathcal B\}$, we apply the lift
\eqref{eq:U6-hat} directly to \eqref{eq:U-affine-polynomial}, keeping
$y$ as a formal variable:
\begin{equation}\label{eq:Udef-pilot}
 \widehat{\mathcal P}^{\mathcal S}_{x,x'}(y)
 =\widehat{\mathcal P_{x,x'}}(\mathcal S;g_x)(y)
 =\sum_{j=0}^J\psi_{j,x}(x')^\top
           \widehat{F_{j,x}(\cdot;y)}(\mathcal S;g_{j,x}).
\end{equation}
The second equality follows from linearity, so only degrees up to
$\max_jD_j$ are required. The pilot is affine in $y$, with coefficients given by
\begin{equation}\label{eq:U-affine-coefficients}
 u_{\mathcal S}=\widehat{\mathcal P}^{\mathcal S}_{x,x'}(0),\qquad
 v_{\mathcal S}=-\partial_y\widehat{\mathcal P}^{\mathcal S}_{x,x'}(y),
 \qquad \widehat{\mathcal P}^{\mathcal S}_{x,x'}(y)=u_{\mathcal S}-yv_{\mathcal S}.
\end{equation}

\subsection*{Pair statistic}
The pair statistic uses the kernel $k_h$ of \cref{sec:pair-risk},
with $h=2^{-J_h}$.
Since $J_h\ge J$, its pair points satisfy $x'\in I_J(x)$.
On the pair set, we define the coefficient vector
$c_{\mathcal S}=(1,v_{\mathcal S}-1,-u_{\mathcal S})^\top$, so that
\[
 c_{\mathcal S}^\top(y',y,1)^\top
 =y'-y-\widehat{\mathcal P}^{\mathcal S}_{x,x'}(y).
\]
Let $\widehat N,\widehat D$ be given by \eqref{eq:pair-statistics},
with $m=n_b$ and the evaluation block $\mathcal P$. The estimator is
\begin{equation}\label{eq:Udef-estimator}
 \widehat V_n=\operatorname{clip}_{I_V}
       \left(\frac{\widehat N}{\max\{\widehat D,1/4\}}\right).
\end{equation}
As in \eqref{eq:pair-statistics}, a summand is zero off the pair set.
The denominator bound in \cref{prop:pair-risk} allows us to use the
threshold $1/4$, once we verify the pilot hypotheses of that
proposition in \cref{lem:U7}.

\subsection*{Parameters for the refined bound}
Recall that $L=\log n$. Approximation at the fine levels requires
roughly $\tau m_j\ge\vartheta(T-t_j)$, while the factorial variance costs
$\exp\{m_j(t_j+\log m_j)\}$. At the required degree, the leading logarithmic variance cost is
largest near $t_j=T/2$. The quadratic profile below keeps the factorial
variances within the budget specified below and gives a Gaussian bound for the approximation errors.
Park \cite[Sections~3.3--3.4]{Park2026} also lowers the degree on finer
partitions and balances the geometric approximation error against a
factorial variance, and our companion paper uses a related balance
\cite[Section~4.4]{AronowKallusLopatto2026}.
For $n\ge3$, we set
\begin{equation}\label{eq:Udef-S}
 C_S=1+\frac{2(\eta+2)}{\beta},\qquad S=a_*\sqrt L-C_S.
\end{equation}
Since $\beta a_*^2/4=2\varpi$, we eventually have $S\ge1$ and
\begin{equation}\label{eq:Udef-budget}
 \beta S^2/4+(\eta+2)S\le2\varpi L.
\end{equation}
When $S\ge1$, we define
\begin{equation}\label{eq:Udef-degree-constants}
 C_D=\max\{1,4\beta\},\quad \ell_S=\log(C_0C_DS),\quad
 \gamma_*=(r+1/2)/\vartheta,\quad
 T^\circ=S-\ell_S-\gamma_*\log S,
\end{equation}
where $C_0\ge1$ is the constant fixed in \eqref{eq:U7-constants}, and
the quadratic profile
\begin{equation}\label{eq:Udef-degree-profile}
 \mathcal Q_S(z)=S/4+(S-z)^2/S.
\end{equation}
We will need the identities
\begin{align}
 z+\mathcal Q_S(z)-S&=(z-S/2)^2/S,\label{eq:Udef-profile-slack}\\
 z\mathcal Q_S(z)&=S^2/4-(S-z)(z-S/2)^2/S\le S^2/4
 \quad(z\le S).\label{eq:Udef-profile-variance}
\end{align}
The candidate parameters are
\begin{equation}\label{eq:Udef-params}
\begin{aligned}
 J&=\left\lfloor\frac{L+T^\circ}{d\log2}\right\rfloor,
 \quad K_*=2^{dJ},\quad T=\log(K_*/n),\\
 t_j&=\log(K_j/n),\quad z_j=t_j+\ell_S,\quad u_j=z_j-S/2,\\
 D_j&=\lfloor\beta\mathcal Q_S(z_j)\rfloor,
 \quad m_j=D_j-q_s\quad(0\le j\le J),\\
 h^\circ&=n^{-2/(d+4)}K_*^{4\vartheta/(d+4)},
 \quad J_h=\lceil\log_2(1/h^\circ)\rceil,\quad h=2^{-J_h},\\
 R_*&=(h^\circ)^2K_*^{-2\vartheta}=n^{-1}(h^\circ)^{-d/2}.
\end{aligned}
\end{equation}
Level-dependent quantities are used only if $J\ge0$. Rounding gives
\begin{equation}\label{eq:Udef-params-rounding}
 T^\circ-d\log2<T\le T^\circ,\qquad h^\circ/2<h\le h^\circ.
\end{equation}
There are $O(L)$ levels, the maximum degree is $O(L^{3/2})$, and the
fine levels have degree $O(\sqrt L)$. The shift $\ell_S$ accounts for the
factorial variance, and the term $\gamma_*\log S$ accounts for the
prefactor in the polynomial approximation and for the $O(\sqrt S)$
contributing levels.
\Cref{sec:U-allocation} verifies these statements and admissibility.

For every $n$, the estimator is \eqref{eq:Udef-estimator} if $n\ge3$,
$S\ge1$
and \eqref{eq:Udef-admissible} holds, and $v_-$ otherwise.
The estimator defined by this convention is jointly Borel in the
observations. Indeed, the known features are Borel, the extraction of
polynomial coefficients is a finite algebraic operation, and all
factorial sums are finite. The pilots are jointly Borel in their
anchors, evaluation points and pilot samples. The pair indicators and
the zero convention off pairs are Borel, while the flooring of the
denominator and the clipping are continuous. The fallback branch
depends only on $n$ and on known constants.
\section{The pilot and its risk bound}\label{sec:U-variance}
\label{sec:U-equation}

Fix admissible parameters \eqref{eq:Udef-admissible}, and let $\mu$
be the law of one observation. For a pair point $\mathsf w=(x,x')$
with $x'\in I_{J_h}(x)$, we have $|x'-x|\le\sqrt d\,h$ and
\begin{equation}\label{eq:U8-hK}
 hK_*^{1/d}=2^{J-J_h}\le1.
\end{equation}
We denote by $M_h$ the pair measure of \eqref{eq:U7-mass}. For either
pilot block, we center the coefficient vector:
\begin{equation}\label{eq:U-pilot-vector}
 \bar c=\E c_{\mathcal S},\qquad
 \zeta_{\mathcal S}=c_{\mathcal S}-\bar c.
\end{equation}
We fix one constant
\begin{equation}\label{eq:U7-constants}
 C_0\ge1
\end{equation}
that is sufficiently large for the feature and derivative bounds below,
and we define
\begin{gather}
 w_j=\sum_{k=1}^{D_j}k!\left(\frac{C_0K_j}{n}\right)^k,
 \label{eq:U7-w}\\
 H_*=\max_{0\le j\le J}w_j\max\{1,n/K_j\},
 \qquad W=h^2K_*^{2/d}H_*,\label{eq:U7-envelope}\\
 b_h=h\left[K_*^{-\vartheta}
       +\sum_{j=0}^JK_j^{-\vartheta}(m_j+1)^r\rho^{m_j}\right].
 \label{eq:U8-bh}
\end{gather}
All constants in this subsection depend only on the fixed model
parameters, uniformly over levels, anchors and polynomial degrees.

\begin{proposition}[Polynomial-pilot risk]\label{lem:U7}\label{lem:U8}
For every admissible choice, the pilots and pair statistics are square
integrable, $\E_\theta\widehat D\ge1/2$, and
\begin{equation}\label{eq:U-master-risk}
 \sup_{\theta\in\Theta}\|\widehat V_n-V\|_{L^2}
 \le C\{b_h^2+(n^{-1/2}+n^{-1}h^{-d/2})(1+W)\}.
\end{equation}
The pilot mean is
\begin{equation}\label{eq:U7-mean}
 \E\widehat{\mathcal P}^{\mathcal S}_{x,x'}(y)
 =\sum_{j=0}^J\psi_{j,x}(x')^\top F_j(x;y),
\end{equation}
and its coefficient vector satisfies
\begin{align}
 \sup_{\mathsf w}|\bar c(\mathsf w)|&\le C,\label{eq:U7-meanbound}\\
 \sup_{\mathsf w}\E|\zeta_{\mathcal S}(\mathsf w)|^2&\le CW,
 \label{eq:U7-pointwise}\\
 \left\|\int\nu^\top\zeta_{\mathcal S}\,dM_h\right\|_{L^2}
 &\le C\sqrt{W/n}\,\|\nu\|_{L^2(M_h)}
 \quad(\nu\in L^2(M_h;\mathbb R^3)).\label{eq:U7-linear}
\end{align}
\end{proposition}

\begin{proof}
The current and parent raw features are supported in
$I_j^\dagger(x)=I_{\max\{j-1,0\}}(x)$, and they satisfy
\begin{equation}\label{eq:U7-feature}
 |g_{j,x}(o)|\le CK_j(1+|y_o|)\1_{I_j^\dagger(x)}(z_o).
\end{equation}
We enlarge the fixed constant to cover the root-parent block $I_r$.
By \eqref{eq:U-M2} and $|I_j^\dagger(x)|\le2^d/K_j$, we have
\begin{equation}\label{eq:U7-feature-second}
 \int|g_{j,x}|^2d\mu\le C_gK_j
\end{equation}
for some fixed constant $C_g$. The lift lemma implies square
integrability and \eqref{eq:U7-mean}. On pair points, we have
\begin{equation}\label{eq:U7-psi}
 |\psi_{j,x}(x')|\le ChK_j^{1/d}.
\end{equation}
The order-zero bound in \cref{lem:U5} and
$\sum_{j=0}^JK_j^{1/d}\le2K_*^{1/d}$ imply that both coefficients of
the mean pilot are bounded by $ChK_*^{1/d}\le C$.

For either $\Phi_j=F_{j,x}(\cdot;0)$ or
$\Phi_j=-\partial_yF_{j,x}$, we consider
\[
 Z_j(\mathsf w)=\psi_{j,x}(x')^\top
 \{\widehat\Phi_j(\mathcal S;g_{j,x})-\Phi_j(\theta_j(x))\}.
\]
Its raw kernel of order $k$ is
$\psi_{j,x}(x')^\top D^k\Phi_j(\theta_j(x))
[g_{j,x}(o_1),\ldots,g_{j,x}(o_k)]$.
Using the derivative bound $M_1k!C_1^k$ and
\eqref{eq:U7-feature-second}, we find that the squared norm of this
kernel, divided by $k!\Lambda^k$, is at most
$Ch^2K_j^{2/d}k!(6C_1^2C_gK_j/n)^k$.
We choose $C_0\ge6C_1^2C_g$ before fixing the allocation.
We partition the pair points by the parent cells of their anchors. Raw
kernels from distinct cells have disjoint observation supports, and each
block of pair points has $M_h$-mass at most $C/K_j$. By
\eqref{eq:U6-localized}, we have
\[
 \E Z_j(\mathsf w)^2\le Ch^2K_j^{2/d}w_j,\qquad
 \left\|\int\nu Z_jdM_h\right\|_2^2
 \le Ch^2K_j^{2/d-1}w_j\|\nu\|_{L^2(M_h)}^2.
\]
Since $w_j\le H_*$ and $w_j/K_j\le H_*/n$, Minkowski's inequality and
the same geometric sum give \eqref{eq:U7-pointwise} and
\eqref{eq:U7-linear} for the last two coordinates of
$\zeta_{\mathcal S}$, and its first coordinate is zero. By Tonelli's
theorem, we have $\E\int|\zeta_{\mathcal S}|^2dM_h\le CW<\infty$, which
justifies the almost-sure integrals.

For $\mathbf f=(f(x'),f(x),1)^\top$, the mean residual is
\begin{equation}\label{eq:U8-gbar}
 r_h=\bar c^\top\mathbf f=f(x')-f(x)
       -\sum_{j=0}^J\psi_{j,x}(x')^\top F_j(x;f(x)).
\end{equation}
We subtract the telescoping finest fit of \cref{lem:U2}.
Its error is at most $ChK_*^{-\vartheta}$, and by \cref{lem:U5} and
\eqref{eq:U7-psi}, the remaining sum is at most
$Ch\sum_jK_j^{-\vartheta}(m_j+1)^r\rho^{m_j}$.
We conclude that $\|r_h\|_\infty\le Cb_h$. These estimates show that the
pilot hypotheses of \cref{prop:pair-risk} hold, with $b=Cb_h$ and
variance budget $W$.
The pilot coefficient fields are jointly Borel by the polynomial
construction.
Admissibility implies that $n_b\ge2$ and $n_b\ge n/6$. We apply the
proposition with $m=n_b$, $c_*=1/6$ and $a=1/4$, and it proves the
denominator bound, square integrability and \eqref{eq:U-master-risk}.
Further clipping to $I_V$ fixes $V$ and cannot increase the error.
\end{proof}
\section{One allocation for bias and variance}\label{sec:U-allocation}

We use the quadratic profile and parameters of
\eqref{eq:Udef-params}. All constants below depend only on
$(s,d,p_-,p_+,C_0)$. By the spatial balance and the dyadic rounding, we
have
\begin{equation}\label{eq:Urate-h}
 (h^\circ)^{(d+4)/2}=n^{-1}K_*^{2\vartheta},\qquad
 R_*=(h^\circ)^2K_*^{-2\vartheta}=n^{-1}(h^\circ)^{-d/2},
 \qquad h^\circ/2<h\le h^\circ.
\end{equation}
Since $d\vartheta=s-1$ and $K_*=ne^T$, we also have
\begin{equation}\label{eq:Urate-scale}
 (h^\circ)^2K_*^{2/d}=e^{-2\varpi L+\eta T},\qquad
 R_*=n^{-\lambda}e^{-2(s-1)T/(d+4)}.
\end{equation}

\begin{lemma}[Allocation]\label{lem:U12}\label{lem:U13}
For all sufficiently large $n$, the parameters \eqref{eq:Udef-params}
are admissible and satisfy
\begin{equation}\label{eq:U12-sum}
 \sum_{j=0}^J K_j^{-\vartheta}(m_j+1)^r\rho^{m_j}
 \le CK_*^{-\vartheta}.
\end{equation}
Consequently,
\begin{equation}\label{eq:U12-bias}
 b_h\le ChK_*^{-\vartheta},\qquad b_h^2\le CR_*.
\end{equation}
Moreover, with $c=(d-4s)/\{4(d+4)\}>0$,
\begin{equation}\label{eq:U13-sums}
 W\le1,\qquad n^{-1/2}\le n^{-c}R_*,
\end{equation}
and
\begin{equation}\label{eq:U13-pairs}
 n^{-1}h^{-d/2}\le2^{d/2}R_*.
\end{equation}
The resulting risk scale is
\begin{equation}\label{eq:Urate-R}
 R_*\asymp n^{-\lambda}e^{-\kappa\sqrt L}
 L^{\frac{s-1}{d+4}+\frac{d(q_s-1/2)}{d+4}}.
\end{equation}
\end{lemma}

\begin{proof}
By the choice of $S$, we have $S\asymp\sqrt L$ and, eventually,
$\beta S^2/4+(\eta+2)S\le2\varpi L$.
The floor in $J$ gives $T=T^\circ+O(1)>0$ and
\begin{equation}\label{eq:Urate-range}
 -L+\ell_S\le z_j\le T+\ell_S
 \le S-\gamma_*\log S\le S.
\end{equation}
Since $\mathcal Q_S(z)=S/4+(S-z)^2/S$, all degrees satisfy
$D_j\ge\beta S/4-1$ and $\max_jD_j=O(L^{3/2})=o(n)$.
Then $J\ge0$, $m_j\ge0$, and the sample-size conditions in
\eqref{eq:Udef-admissible} hold eventually. In addition, we have
\begin{equation}\label{eq:U13-reserve}
 h^2K_*^{2/d}\le e^{-2\varpi L+\eta T}
 \le e^{-\beta S^2/4-2S}<1.
\end{equation}
Since $hK_*^{1/d}=2^{J-J_h}$, this proves $J_h\ge J$ and hence
admissibility.

For the bias, recall that $u_j=z_j-S/2$. Using the identity
$z+\mathcal Q_S(z)-S=(z-S/2)^2/S$, the degree floor, and
$S-\ell_S-T\ge\gamma_*\log S$, we obtain
\begin{equation}\label{eq:U12-geom}
 K_j^{-\vartheta}\rho^{m_j}
 \le CK_*^{-\vartheta}S^{-r-1/2}e^{-\vartheta u_j^2/S}.
\end{equation}
Here we also used that $\tau\beta=\vartheta$ and
$m_j\ge\beta\mathcal Q_S(z_j)-q_s-1$.
For $S\ge1$, we have $\mathcal Q_S(z_j)\le S+3u_j^2/(2S)$, which implies
that $(m_j+1)^r\le CS^r(1+u_j^2/S)^r$.
Absorbing this polynomial factor into half of the Gaussian exponent, we
find
\begin{equation}\label{eq:U12-term}
 K_j^{-\vartheta}(m_j+1)^r\rho^{m_j}
 \le CK_*^{-\vartheta}S^{-1/2}e^{-\vartheta u_j^2/(2S)}.
\end{equation}
Since the $u_j$ have spacing $d\log2$, their Gaussian sum is at most
$C\sqrt S$. This proves \eqref{eq:U12-sum}, and
\eqref{eq:U12-bias} follows from $h\le h^\circ$ and
\eqref{eq:Urate-h}.

For the variance, we abbreviate $x_j=C_0K_j/n=e^{z_j}/(C_DS)$.
When $z\le0$ and $S\ge2$, differentiation shows that
$e^z\mathcal Q_S(z)\le5S/4$. This implies that
$D_jx_j\le5\beta/(4C_D)\le5/16<1/2$ for $z_j\le0$.
The factorial series then has ratios of successive terms at most $1/2$,
which gives $w_j\le2x_j$ and
$w_j\max\{1,n/K_j\}\le2C_0$.
For $0<z_j\le S$, the profile gives $D_j\le5\beta S/4\le C_DS$
and $D_jz_j\le\beta S^2/4$. Using $k!\le D_j^k$, we have
\begin{equation}\label{eq:U13-wj}
 w_j\le\sum_{k=1}^{D_j}e^{kz_j}
 \le C_DS e^{\beta S^2/4}.
\end{equation}
Since $\max\{1,n/K_j\}\le C_0C_DS$ on these fine levels,
$H_*\le2C_0C_D^2S^2e^{\beta S^2/4}$. The reserve gives
\begin{equation}\label{eq:U13-fine}
 W\le2C_0C_D^2S^2e^{-2S}\longrightarrow0.
\end{equation}
The pair-noise bound follows from \eqref{eq:Urate-h}.
Further, since $T\le S=O(\sqrt L)$ and $1/2-\lambda=2c$, we have
$n^{-1/2}/R_*\le n^{-2c}e^{C\sqrt L}\le n^{-c}$ eventually.
This proves \eqref{eq:U13-sums} and \eqref{eq:U13-pairs}.

Finally, since $\log S=\tfrac12\log L+O(1)$, we have
\[
 T=a_*\sqrt L-\frac{1+\gamma_*}{2}\log L+O(1).
\]
Substituting this into \eqref{eq:Urate-scale}, and using
$2(s-1)a_* /(d+4)=\kappa$ and
$\gamma_*=(q_s-1/2)/\vartheta$, we obtain \eqref{eq:Urate-R}.
\end{proof}
\section{Completion of the upper bound}\label{sec:U-final}

\begin{proof}[Proof of \cref{thm:upper-sharp}]
The estimator is Borel by the construction in \cref{sec:U-estimator}.

For all sufficiently large $n$, \cref{lem:U12} supplies admissible
parameters such that $b_h^2\le CR_*$, $W\le1$, and
$n^{-1/2}+n^{-1}h^{-d/2}\le CR_*$.
Inserting these bounds into the master bound \eqref{eq:U-master-risk},
we obtain
\[
 \mathfrak r_n\le\sup_{\theta\in\Theta}
 \|\widehat V_n-V\|_{L^2}\le CR_*.
\]
Using the rate identity \eqref{eq:Urate-R}, we see that this is
exactly the asserted bound, with constants that depend only on the stated model parameters.
\end{proof}
\chapter{Preliminaries for the lower bounds}\label{sec:lower-prelim}

In this section, we collect the inputs common to the lower bounds in
\cref{thm:main,thm:main2,thm:parametric}: a ternary law,
a score-to-risk comparison, and smooth weights whose squares
sum to one. We also prove the parametric $n^{-1/2}$ lower bound, which
is valid for every $s>0$. The lower bound in
\cref{thm:main} is proved in \cref{sec:lower}, and the lower
bound in \cref{thm:main2} is proved in \cref{sec:small-s}.

\section{The ternary law}\label{sec:prelim-ternary}

Every construction for the lower bounds uses the following response
law, which already appears in the third version of this paper
\cite[Section~3, equation~(3.1)]{AronowLopatto2026v3}. Fix constants
\[
 v_-<v<\min(v_+,\sqrt{C_4}),\qquad
 v<a_Y^2<C_4/v,\qquad a_Y>0.
\]
Such a choice is possible because $v_-^2<C_4$. For a regression value $f$
and conditional variance $V$, let $Y$ take values in
$\{-a_Y,0,a_Y\}$ with probabilities
\begin{equation}\label{eq:ternary}
 q_{\pm a_Y}(f,V)
 =
 \frac{V+f^2\pm a_Yf}{2a_Y^2},
 \qquad
 q_0(f,V)
 =
 1-\frac{V+f^2}{a_Y^2}.
\end{equation}
Then $\E Y=f$, $\Var Y=V$, and
\[
 \E(Y-f)^4=a_Y^2V+(6V-3a_Y^2)f^2+3f^4.
\]
At $(f,V)=(0,v)$, all three probabilities are positive and the
centered fourth moment is $a_Y^2v<C_4$. By continuity, there exist
$\rho,c>0$ such that, whenever $|f|\le\rho$ and $|V-v|\le\rho$,
we have $V\in(v_-,v_+)$, $q_y(f,V)\ge c$ for every response value,
and $\E(Y-f)^4\le C_4$. This means that the response constraints hold uniformly on this
neighborhood. Note also that the probabilities are quadratic in $f$
and affine in $V$, with $\partial_f^2q_y=2\partial_Vq_y$.
Randomizing the regression mean in order to hide a change of the
variance is also the device of the lower-bound proof
in~\cite[Section~4.5]{Spokoiny2002}, with Gaussian regression values on
smooth bumps, and of the fixed-design lower bounds of Wang et
al.\ \cite[Section~6.2]{WangBrownCaiLevine2008}, who use finite Gaussian moment
matching; for the ternary law, the last identity makes this
matching exact for one response.

\section{One score-to-risk comparison}\label{sec:prelim-risk}

An accurate estimator must follow the target along any path of mixtures.
The following lemma bounds how quickly its mean can move. The random
variable $R_t$ in the lemma need only represent derivatives of data
observables; it need not be the score of the joint law of the latent
parameter and the data. Let $J_V=[v_-,\min\{v_+,\sqrt{C_4}\}]$ and
$D_V=|J_V|$. The fourth-moment bound implies that every variance in
$\Theta$ belongs to $J_V$.

\begin{lemma}[Score-to-risk comparison]\label{lem:path-risk}
For $0\le t\le\delta$, let $\mathbb E_t$ describe a mixture of
$n$-observation laws with deterministic variance $V_t\in J_V$.
Suppose each mixture assigns probability at most $\epsilon$ to
parameters outside $\Theta$. Suppose also that, for every bounded
data-only statistic $T$, $t\mapsto\mathbb E_tT$ is absolutely continuous
and, almost everywhere,
\[
 \frac d{dt}\mathbb E_tT=\mathbb E_t[TR_t],\qquad
 \mathbb E_tR_t=0.
\]
If $B=\int_0^\delta\|R_t\|_{L^2(\mathbb E_t)}\,dt<\infty$ and
$\Delta=|V_\delta-V_0|$, then
\begin{equation}\label{eq:path-risk}
 \mathfrak r_n^2\ge\frac{\Delta^2}{(2+B)^2}-D_V^2\epsilon.
\end{equation}
\end{lemma}

\begin{proof}
We clip $T$ to $J_V$, which cannot increase its risk on $\Theta$, and
set $r_T^2=\sup_{\theta\in\Theta}\mathbb E_\theta(T-V)^2+D_V^2\epsilon$.
Since the exceptional states contribute at most $D_V^2\epsilon$, we have
$\mathbb E_t(T-V_t)^2\le r_T^2$. Then, for $m(t)=\mathbb E_tT$, we have
\[
 |m(t)-V_t|\le r_T,\qquad
 |m'(t)|=|\mathbb E_t[(T-V_t)R_t]|
 \le r_T\|R_t\|_{L^2(\mathbb E_t)}.
\]
Bounding the two endpoint errors and integrating the derivative, we
find that $\Delta\le(2+B)r_T$. It remains to square this bound and take
the infimum over clipped estimators.
\end{proof}

The lemma combines the covariance inequality behind the information
bound, which controls the derivative of the mean of a statistic
\cite[Chapter~2, Theorem~5.10]{LehmannCasella1998}, with a comparison
of the means of an estimator under two mixtures, as in the constrained
risk inequality for mixtures that Cai and Low
\cite[Section~2, Theorem~2 and Corollary~1]{CaiLow2011} develop from
that of Brown and Low \cite{BrownLow1996}. Here the comparison integrates the derivative
along a path of mixtures.

\section{The parametric lower bound}\label{sec:prelim-parametric}

\begin{proof}[Proof of the lower bound in \cref{thm:parametric}]
We use $p\equiv1$, $f=0$, and the ternary law with
$V_t=v-a_nt$, $0\le t\le1$, where $a_n=c_0n^{-1/2}$.
If the fixed constant $c_0>0$ is sufficiently small, then the whole
path is admissible for every $n$. The ordinary likelihood score of this
path is
\[
 R_t=-a_n\sum_{i=1}^n
 \frac{\partial_Vq_{Y_i}(0,V_t)}{q_{Y_i}(0,V_t)},\qquad
 \mathbb E_tR_t^2=
 \frac{na_n^2}{V_t(a_Y^2-V_t)}\le Cc_0^2.
\]
The summands are independent and centered, and the equality follows
directly by summing over the three outcomes. By decreasing $c_0$ if
necessary, we can ensure that $B\le1$. Applying \cref{lem:path-risk}
with $\epsilon=0$ and $\Delta=a_n$, we obtain
\begin{equation}\label{eq:parametric-lower}
 \mathfrak r_n\ge cn^{-1/2}.
\end{equation}
Since the zero function belongs to every H\"older class, this bound
holds for all $s>0$ and $d\ge1$.
\end{proof}

\section{Smooth weights}\label{sec:prelim-windows}

Both nonparametric lower bounds use one family of smooth weights. As
in the third version of this paper
\cite[Section~1.2, Step~3]{AronowLopatto2026v3}, their squares sum to
one at every point.
Let $Q=(-1,1)^d$, and fix a nonnegative
$\psi\in C_c^\infty(Q)$ that is positive on $[-1/2,1/2]^d$. Then
$\sum_{z\in\mathbb Z^d}\psi(\cdot-z)^2$ is smooth, $1$-periodic and
bounded below by a positive constant, and
\begin{equation}\label{eq:lattice-windows}
 w=\psi\Bigl(\sum_{z\in\mathbb Z^d}\psi(\cdot-z)^2\Bigr)^{-1/2}
 \in C_c^\infty(Q)
 \quad\text{satisfies}\quad
 \sum_{z\in\mathbb Z^d}w(u-z)^2=1,\qquad u\in\mathbb R^d.
\end{equation}
We write $\mathbb T^d=\mathbb R^d/\mathbb Z^d$ for the torus and
$\|x-x'\|_\infty$ for the sup-distance on it. Let $h=1/n_h$ with $n_h\ge4$
an integer, and let $\mathcal J=(\mathbb Z/n_h\mathbb Z)^d$. For
$j\in\mathcal J$, the \emph{patch} $Q_j$ is the set of $x\in\mathbb T^d$ with
$\|x-hj\|_\infty<h$; for $x\in Q_j$, $(x-hj)/h\in Q$ denotes the
rescaled representative of $x-hj$ in $(-h,h)^d$, and the weights are
$w_j(x)=w((x-hj)/h)$ for $x\in Q_j$ and $w_j(x)=0$ otherwise. Because
$n_h\ge4$, distinct points $z\in\mathbb Z^d$ with $\|x/h-z\|_\infty<1$
give distinct $j\in\mathcal J$. Then \eqref{eq:lattice-windows}
implies the partition identity below, and differentiation gives the
derivative bounds:
\begin{equation}\label{eq:smooth-windows}
 \sum_jw_j(x)^2=1,\qquad
 \|\partial^\gamma w_j\|_\infty\le C_\gamma h^{-|\gamma|}
 \quad\text{for every multi-index }\gamma.
\end{equation}
There are $h^{-d}$ patches, each of volume $(2h)^d$ and meeting at most
$3^d-1$ others. Each point of $\mathbb T^d$ lies in at most $2^d$ patches,
and two points of a common patch are at sup-distance less than $2h$.
\chapter{The case \texorpdfstring{$0<s\le1$}{0 < s <= 1}}
\label{sec:small-s}

This section gives a different proof of \cref{thm:main2}, whose rate
was established in the second version of this paper
\cite[Theorem~1.2 and Appendix~A]{AronowLopatto2026v2}.
Throughout, we
assume that $0<s\le1$. All constants may depend on the fixed parameters
that define the model class, but not on $n$.

\section{Upper bound}
The upper bound is the case $t=s$ of
\cref{prop:elementary-upper}. The estimator there averages half the
squared difference of paired responses whose design points are close,
and the proof also covers $s=1$ and the boundary $d=4s$. Shen et al.\
\cite[Section~4.1, Proposition~1]{ShenGaoWittenHan2020} prove this rate
for their product-kernel ratio estimator when the standardized errors
are independent of $X$; the proof of \cref{prop:pair-risk} allows the
error law $Q_x$ to depend on $x$.

\section{Lower bound}

\begin{proof}
If $d\le4s$, the parametric bound \eqref{eq:parametric-lower} already
has the required order. We now suppose that $d>4s$. We use the
smooth weights \eqref{eq:smooth-windows} at scale $h=1/n_h$ with
$n_h\ge4$, and set
\[
 \eta=bh^s,\qquad
 \xi_j\ \hbox{independent with law}\quad
 \nu_t=(1-t)\delta_0+\tfrac t2(\delta_{-1}+\delta_1),
 \qquad 0\le t\le1.
\]
Let $p\equiv1$, $f_\xi=\eta\sum_jw_j\xi_j$ and
$V_t=v-\eta^2t$, and let the responses follow the ternary law.
By bounded overlap, we have
\[
 \|f_\xi\|_\infty\le C\eta,\qquad
 |f_\xi(x)-f_\xi(x')|
 \le C\eta\min\{1,\|x-x'\|/h\}
 \le Cb\|x-x'\|^s.
\]
For a sufficiently small fixed $b>0$, every state is then admissible,
including at $s=1$. The periodic weights give the required extension
beyond the design cube. We will apply \cref{lem:path-risk} with
$\epsilon=0$ and separation $\Delta=\eta^2$.

Next, we represent derivatives of expectations of bounded statistics of
the data along this path, using the latent coefficients. We write
$\xi^{j,a}$ for $\xi$ with coordinate $j$ set to $a$, and define
\[
 D_jF(\xi)=\tfrac12F(\xi^{j,1})+
          \tfrac12F(\xi^{j,-1})-F(\xi^{j,0}).
\]
Since $\nu_t'=(\delta_{-1}+\delta_1)/2-\delta_0$, differentiating
the finite product prior replaces its $j$th factor by exactly this
signed measure. For a fixed sample, let
$L_\xi=\prod_{i=1}^nq_{Y_i}(f_\xi(X_i),V_t)$ and
\begin{equation}\label{eq:small-s-score}
 R_t=\sum_j r_j,\qquad
 r_j=\frac{D_jL_\xi-
       \eta^2\sum_i w_j(X_i)^2\partial_{V_i}L_\xi}{L_\xi},
\end{equation}
where $\partial_{V_i}$ differentiates only the $i$th factor.
Differentiating the prior and the likelihood, and then using
$\sum_jw_j^2=1$, we obtain
$\frac d{dt}\mathbb E_tT=\mathbb E_t[TR_t]$ for all bounded $T$.
Since all sums are finite and the likelihood is uniformly positive for
fixed $n$, the derivative and the integration interchange directly.

Only observations in the support patch $Q_j$ enter $r_j$, because the
other likelihood factors cancel. After this cancellation, summing the remaining numerator over all
configurations of the responses in the patch gives zero, since $D_j1=0$
and every response probability sums to one. This implies that
$\mathbb E_t[r_j\mid\xi,X_1,\ldots,X_n]=0$, and that scores from
disjoint patches are conditionally orthogonal. Further, the
quadratic ternary probabilities satisfy the exact identity
\[
 \tfrac12q_y(g+\eta w,V)+\tfrac12q_y(g-\eta w,V)-q_y(g,V)
 =\eta^2w^2\partial_Vq_y(g,V).
\]
The definition gives $r_j=0$ when $Q_j$ contains no observation. For
one observation $X_i\in Q_j$, this identity, with $w=w_j(X_i)$ and with
$g$ excluding the $j$th field, gives the same conclusion. Indeed,
$\partial_Vq_y$ does not depend on the mean, so the variance term in
\eqref{eq:small-s-score} also equals $\eta^2w^2\partial_Vq_y(g,V)$,
times the likelihood factors of the other observations.
This is the cancellation that makes the rate possible.
Independent local perturbations of the mean also mask a change of the
variance, through finite Gaussian moment matching, in the lower bound of
Shen et al.\ \cite[Sections~2.2 and~5]{ShenGaoWittenHan2020}, whose
design density vanishes between the bumps. Here the ternary law and the
identity $\sum_jw_j^2=1$ make this cancellation exact for one
observation with a design density bounded away from zero.

For the remaining counts, we need only a second-order Taylor
estimate, which we include for completeness. Given $k$ observations in
$Q_j$, their response product is $H(a)=\prod_{i=1}^kq_{y_i}(g_i+\eta w_i a,V_t)$, where
$g_i$ excludes the $j$th field. Along $-1\le a\le1$, all probabilities
are at most one, and their first two derivatives in $a$ are bounded by
$C\eta$ and $C\eta^2$, respectively. Then
$\sup|H''|\le Ck^2\eta^2$, which gives
$|\{H(1)+H(-1)\}/2-H(0)|\le Ck^2\eta^2$.
The variance derivative in \eqref{eq:small-s-score} is at most
$Ck\eta^2$. Since the denominator is at least $c^k$, this yields
$|r_j|^2\le C^k\eta^4$ for $k\ge2$, after increasing $C$.
For $\mu=nh^d\le1$, the binomial count and $|Q_j|\le Ch^d$ give
\[
 \mathbb E_t r_j^2\le
 \eta^4\sum_{k\ge2}\frac{(C\mu)^k}{k!}
 \le C\eta^4\mu^2.
\]
There are $O(h^{-d})$ patches, each meeting only a bounded number of
others. By conditional orthogonality and $2ab\le a^2+b^2$, we have
\[
 \mathbb E_tR_t^2\le C h^{-d}\eta^4\mu^2
 =Cn^2\eta^4h^d.
\]
We take $h=\lceil4n^{2/(d+4s)}\rceil^{-1}$. Then $\mu\le1$ because
$d>4s$, and $B\le Cb^2n h^{(d+4s)/2}\le Cb^2$.
We then decrease $b$ so that $B\le1$. Applying the score-to-risk
comparison, we obtain
\[
 \mathfrak r_n\ge\eta^2/3\asymp n^{-4s/(d+4s)}.
\]
Together with \eqref{eq:parametric-lower} and
\cref{prop:elementary-upper}, this proves \cref{thm:main2}, including
the boundary $d=4s$ covered at the outset.
\end{proof}
\chapter{The lower bound}\label{sec:lower}

Throughout this section, we assume that $s>1$ and that $d>4s$ is an
integer, and the model is that of \cref{sec:main-result}; in
particular, $p_-<1<p_+$.

\begin{theorem}\label{thm:lower}
Let $s>1$ and let $d>4s$ be an integer. Let $\lambda$, $\tau$ and
$\kappa$ be as in \eqref{eq:lambda} and \eqref{eq:main-constants}.
There are $c>0$ and $n_0$, depending only on
$(s,d,p_-,p_+,L_s,v_-,v_+,C_4)$, such that
\[
 \mathfrak r_n\ge c\,n^{-\lambda}e^{-\kappa\sqrt{\log n}}
 (\log n)^{(s-1)/(d+4)}
 \qquad\text{for all }n\ge n_0 .
\]
\end{theorem}

The constants below depend only on these fixed model parameters and
on the design constants in \cref{sec:LB-proof}; they never depend on
the sample, the state, $n,h,M,D,N$ or the nuisance values. They also
do not depend on $\mathcal U$, since the constructed means extend
periodically to $\mathbb R^d$. We write $L=\log n$.

The ternary law depends on the mean through $f$ and $f^2+V$.
A centered perturbation of the mean imitates an increase of the
variance through its diagonal covariance. We correlate the mean with
the design density in order to control this covariance at the observed
points, and then cancel its diagonal part by decreasing $V$. Priors
under which the covariate density and the regression functions are
perturbed together appear in the missing-data lower bound of Robins et
al.\ \cite[Section~3]{RobinsEtAl2009EJS}; here the perturbations of the
density also select the number of observations in a patch.

The third version of this paper used cardinal interpolation at every
count at least two and absorbed signed corrections into a common prior
for the design density \cite[Section~1.2, Steps~1--2, and
Appendix~B]{AronowLopatto2026v3}.
Here we construct the local updates for counts at least three by
cardinal interpolation. At count two, the unit and coarse kernels use ordinary
selectors, while the fine kernel uses a three-point selector with
bounded cell signs. We realize the signed generator through positive
weights on update histories, and an $n$th-power mass tilt cancels the
normalization of the likelihood. The risk criterion depends on the
activity and the local score energy of the construction, and balancing
these two quantities gives the stated rate.

\section{The sample, patches and raw states}\label{sec:LB-setup}

We use the ternary law \eqref{eq:ternary} with its constants
$v,a_Y,\rho$. Let $c_q$ denote its uniform positivity constant, and set
$\mathcal Y=\{-a_Y,0,a_Y\}$ and $\eta_0=\rho/2^{d+1}$. Let
$\mathsf P(\mathbf f,\mathbf V;\mathbf y)=\prod_{i\le k}q_{y_i}(f_i,V_i)$,
with $\mathsf P=1$ when $k=0$. When all variance coordinates equal
$V$, we denote the product by $\mathsf P(\mathbf f,V)$, and the derivative
$\partial_{V_i}$ is taken before this substitution. The ternary formulas give
\begin{equation}\label{eq:LB-heat-one}
 \partial_f^2q_y=2\partial_Vq_y,\qquad
 \sum_{y\in\mathcal Y}\partial_Vq_y=0.
\end{equation}
All mark spaces are standard Borel, and densities are Borel versions
of their almost-everywhere classes.

We identify $[0,1)^d$ with $\mathbb T^d$ and use the sample
space $\mathfrak X=(\mathbb T^d\times\mathcal Y)^n$ with the reference
measure $\Lambda_n=(dx\otimes\#_{\mathcal Y})^{\otimes n}$. For a
positive bounded density function $p$ and a mean $f$, we define
\begin{equation}\label{eq:LB-likelihood}
 \begin{gathered}
 \mathsf U(p,f,V;\xi)=\prod_{i\le n}p(x_i)q_{y_i}(f(x_i),V),\\
 \mathsf L(p,f,V;\xi)=m(p)^{-n}\mathsf U(p,f,V;\xi),\quad m(p)=\int p.
 \end{gathered}
\end{equation}
When $|f|\le\rho$ and $|V-v|\le\rho$, this is the likelihood of $n$
independent observations with design density $p/m(p)$ and the ternary
law, and $P_{p,f,V}$ and $\E_{p,f,V}$ denote the
corresponding law and expectation. For a patch $A$, $\xi_A$ consists
of the observations that lie in the patch, kept in their original
order.

We take the smooth weights \eqref{eq:smooth-windows} at $h=1/n_h$ with
$n_h\ge4$, and set $\mu=nh^d$ and $\omega=\log(1/\mu)$. The chart
$\chi_j:Q_j\to Q=(-1,1)^d$, $x\mapsto(x-hj)/h$, has Jacobian $h^{-d}$,
and $w_j=w\circ\chi_j$ on $Q_j$. Recall that each point lies in at most
$2^d$ patches. For distinct labels, we write $j\sim j'$ when
$Q_j\cap Q_{j'}\ne\emptyset$. Each label has exactly $3^d-1$
neighbors, given by the coordinate offsets
$\{-1,0,1\}^d\setminus\{0\}$ modulo $n_h$.

For integers $M\ge M_1=\lceil4/(p_+-p_-)\rceil$, let
$a_M=p_-+1/M$ and $b_M=p_+-1/M$. The interval between them has width
at least $(p_+-p_-)/2$, and
\begin{equation}\label{eq:LB-tauM}
 \tau_M=\operatorname{arcosh}\frac{b_M+a_M}{b_M-a_M}
 =\log\frac{\sqrt{b_M}+\sqrt{a_M}}{\sqrt{b_M}-\sqrt{a_M}},
 \qquad \tau\le\tau_M\le\tau+C_\tau/M.
\end{equation}
Here $C_\tau$ is a fixed constant, since differentiating the
displayed logarithm gives $\partial_a=\sqrt b/[\sqrt a(b-a)]$ and
$\partial_b=-\sqrt a/[\sqrt b(b-a)]$, and these derivatives are
bounded on the segment from $(p_-,p_+)$ to $(a_M,b_M)$.

For $D\ge1$, define
\begin{equation}\label{eq:frame}
 I_D=\{\beta\in\mathbb N^d:|\beta|\le D-1\},\qquad
 \phi_\beta(u)=(u/8)^\beta,\quad \phi=(\phi_\beta)_{\beta\in I_D},
\end{equation}
where $\mathbb N=\{0,1,\ldots\}$ and $\phi_0=1$. The fixed constant
\[
C_{\rm fr}=\max\{1,\sup_\beta\|w(\cdot)(\cdot/8)^\beta\|_{C^s(\mathbb R^d)}\}
\]
is finite, because on the compact support of $w$, the derivatives
through order $\ell+1$ are bounded by
$C(1+|\beta|)^{\ell+1}8^{-|\beta|}$, and these derivative bounds
control the H\"older seminorm. Let
\begin{equation}\label{eq:LB-coefficient-ball}
 \mathbb B=\{c\in\mathbb R^{I_D}:\|c\|_1\le C_{\rm fr}^{-1}\},
 \qquad F_c=\phi^\top c
\end{equation}
Then $|F_c|\le1$ and $\|G_c\|_{C^s(\mathbb R^d)}\le1$ for
$G_c=wF_c$, extended by zero outside $Q$.

A raw state $\vartheta=(p,(c_j)_j)\in\Theta_{\rm raw}$ consists of a
Borel function $p:\mathbb T^d\to[a_M,b_M]$ and coefficients
$c_j\in\mathbb B$. Its mean is
\begin{equation}\label{eq:LB-mean}
 f_\vartheta=\eta_f\sum_j w_j(F_{c_j}\circ\chi_j),\qquad
 \eta_f=c_fh^s,\quad 0<c_f\le1,
\end{equation}
where each term is zero outside $Q_j$. The vacuum
$\vartheta^\zeta_{\rm vac}$ is the raw state with $p\equiv\zeta$ and
$c_j=0$. The variance is not part of the raw state, and prior
probabilities will be placed on marked histories of these states.
For $V\in[v-\rho,v]$, we define
\begin{equation}\label{eq:LB-normalization}
 \begin{gathered}
 m_\vartheta=\int p_\vartheta,\qquad
 \theta(\vartheta,V)=(p_\vartheta/m_\vartheta,f_\vartheta,V,Q_{\vartheta,V}),\\
 Q_{\vartheta,V;x}
 =\sum_{y\in\mathcal Y}q_y(f_\vartheta(x),V)\delta_{y-f_\vartheta(x)}.
 \end{gathered}
\end{equation}
We abbreviate the corresponding laws and likelihoods by
$P_{\vartheta,V}$, $\E_{\vartheta,V}$, $\mathsf L(\vartheta,V;\xi)$ and
$\mathsf U(\vartheta,V;\xi)$. When $\eta_f\le\eta_0$, they exist for
every raw state, including raw states whose normalized density lies
outside the model, and the likelihoods satisfy
\begin{equation}\label{eq:LB-likelihood-bounds}
 (a_Mc_q/b_M)^n\le\mathsf L\le(b_M/a_M)^n,\qquad
 (a_Mc_q)^n\le\mathsf U\le b_M^n.
\end{equation}

Let
\begin{equation}\label{eq:LB-Cw}
 C_w=2^d\{\tbinom{d+\ell}{d}+2\}.
\end{equation}
\begin{lemma}[Legality of raw states]\label{lem:legality}
Let $n_h\ge4$, $M\ge M_1$, $0<\eta_f=c_fh^s\le\eta_0$,
$\vartheta\in\Theta_{\rm raw}$ and $V\in[v-\rho,v]$.

\textup{(a)} The periodic extension satisfies
$\|\tilde f_\vartheta\|_{C^s(\mathbb R^d)}\le C_wc_f$ and
$|f_\vartheta|\le2^d\eta_f\le\rho/2$.

\textup{(b)} The ternary law has mean $f_\vartheta$, variance
$V\in(v_-,v_+)$, fourth central moment at most $C_4$, and probabilities
in $[c_q,1]$.

\textup{(c)} If $0<c_m\le1/p_+$ and $|m_\vartheta-1|\le c_m/M$,
then $m_\vartheta>0$ and $p_-\le p_\vartheta/m_\vartheta\le p_+$; if $M\ge2$,
$m_\vartheta\ge1/2$.

Consequently, if $C_wc_f\le L_s$, $M\ge2$, and the hypothesis
of \textup{(c)} holds,
$\theta(\vartheta,V)\in\Theta$.
\end{lemma}
\begin{proof}
The periodic extension is
$\tilde f_\vartheta(x)=\eta_f\sum_{z\in\mathbb Z^d}G_{c_{[z]}}(x/h-z)$.
At each point, at most $2^d$ terms contribute, and for a difference at
two points, at most $2^{d+1}$ terms contribute. Scaling the profile
bounds, we obtain
\[
 \|\partial^\gamma\tilde f_\vartheta\|_\infty\le2^dc_f
 \quad(|\gamma|\le\ell),\qquad
 [D^\ell\tilde f_\vartheta]_\alpha\le2^{d+1}c_f.
\]
Adding the $\binom{d+\ell}d$ sup norms proves (a). These bounds also keep the
parameters of the response law in the ternary neighborhood, which implies (b).
For (c), we have $m_\vartheta\ge1-c_m/M>0$ and
\[
 \frac{p_-+1/M}{1+c_m/M}\ge p_-,
 \qquad
 \frac{p_+-1/M}{1-c_m/M}\le p_+.
\]
The first inequality uses that $p_-c_m\le1$, and the second uses
that $p_+c_m\le1$. When $M\ge2$, we have $m_\vartheta\ge1-1/M\ge1/2$.
Finally, the normalized density is Borel, the periodic extension
restricts to $\mathcal U$, and the displayed error kernel is measurable
because $f_\vartheta$ is continuous. These facts verify all the model
conditions.
\end{proof}

\medskip\noindent\emph{Parameter regime.}
We set $q_{\rm rsp}=\max\{2,\lfloor d/(4s)\rfloor\}$ and
$C_D=(d+8)/4$. For $K\ge1$, let $\mathsf{Reg}(K)$ consist of the
parameters with
\begin{equation}\label{eq:LB-regime}
 \begin{gathered}
 M\ge M_1\ \text{integer},\quad D=\lceil C_DM\rceil,\quad
 N\ge1,\quad 0<\mu<1,\quad \omega=\log(1/\mu),\\
 M/K\le\omega\le KM,\quad M^2/K\le\log N\le KM^2,\\
 0<\eta_f\le\eta_0,\quad \eta_f^{4q_{\rm rsp}}N^d\le1.
 \end{gathered}
\end{equation}
The last inequality bounds the Taylor error of the fixed response
rule. Define $C_\sharp=1/(p_-^2c_q)$ and
$M_1'=\max\{M_1,\lceil2/(1-p_-)\rceil,\lceil2/(p_+-1)\rceil\}$.
For $M\ge M_1'$, the density interval contains the value one with a
fixed margin.

The construction enters the statistical argument through its
activity $B_{\rm loc}$ and its local score energy $\mathcal E$. For
the admissible, annihilating and centered updates constructed below,
\cref{lem:score-risk} implies that
\begin{equation}\label{eq:LB-cost-risk}
 \mathfrak r_n^2\ge
 c_{\rm tr}\frac{\eta_f^4}{1+B_{\rm loc}^2+h^{-d}\mathcal E}
 -(v_+-v_-)^2\epsilon_n.
\end{equation}
Here $B_{\rm loc}$ bounds the signed mark weights, $\mathcal E$ is
defined in \eqref{eq:LB-score-energy}, and $\epsilon_n$ is the
mass-concentration bound \eqref{eq:LB-bad-mass}.
\section{Interpolation on one patch}\label{sec:interp}

In this section, we construct interpolation matrices whose evaluated
kernels are diagonal. Their defect and band estimates control the local
score energy and the aliases. We work with the frame \eqref{eq:frame}
and an integer $d>4$, and $\|A\|_{\ell^1}$ denotes the sum of the absolute
values of the entries of a matrix. A \emph{separated representation}
of a permutation-invariant
$E:Q^k\to\mathbb R^{I_D\times I_D}$ is
\[
 E(U)=\int_{\mathcal Z}
 \Sym\!\left[\prod_{l=1}^k b_l(U_l;\zeta)\right]A(\zeta)\,\sigma(d\zeta),
\]
where $\mathcal Z$ is standard Borel, $\sigma$ is finite and positive,
the functions are measurable, $|b_l|\le1$, and
$\int\|A\|_{\ell^1}d\sigma<\infty$. Here $\Sym$ averages over the
$k!$ permutations of the points. We call the last integral the
\emph{cost} of the representation. If $E$ is symmetric, then
symmetrizing $A$ does not increase the cost.

\begin{lemma}\label{lem:interp}
There are $C_A,C_I\ge1$, depending only on $d$, such that for integers
$D\ge2$, $2\le k\le D$ and real $T\ge1$ there are continuous,
permutation-invariant symmetric matrices $E_{k,T}(U)$ and continuous
weights $\widehat w_i(T;U)\in[0,1]$ satisfying
\begin{equation}\label{eq:interp-diag}
 \phi(U_i)^TE_{k,T}(U)\phi(U_l)
 =\widehat w_i(T;U)\1_{\{i=l\}}.
\end{equation}
The weights increase with $T$. Set $E_{k,T}=\widehat w_i(T;\cdot)=0$
for $0<T<1$. Each $E_{k,T}$ and each band $E_{k,T}-E_{k,L}$,
$0<L\le T$, has a separated representation of cost at most
\begin{equation}\label{eq:interp-cost}
 C_A^kT^2(1+\log T)^{k-2}.
\end{equation}
With $L_+=\max\{1,L\}$,
\begin{align}
 \int_{Q^k}\|E_{k,T}-E_{k,L}\|_{\ell^1}^2\,dU
 &\le C_I^kL_+^{4-d}(1+\log L_+)^{k-2},
 \label{eq:interp-l2}\\
 \int_{Q^k}\sum_i(1-\widehat w_i(T;U))\,dU
 &\le C_A^kT^{-d}(1+\log T)^{k-2}.
 \label{eq:interp-defect}
\end{align}
\end{lemma}

\begin{proof}
For a polynomial in $u,u'$, let $\|\cdot\|_{\rm coef}$ be the sum
of the absolute values of its coefficients in the monomials
$(u/8)^\beta(u'/8)^{\beta'}$. This norm is submultiplicative and equals
the $\ell^1$ norm of the coefficient matrix. In this proof, all
constants $C$ depend only on $d$. With $\lambda=(8d)^{-1}$ and
$\alpha(v)=\lambda|v|^2$, we define
\begin{equation}\label{eq:interp-partition}
 f_t(v)=t\alpha(v)^2e^{-t\alpha(v)},\qquad
 \int_0^\infty f_t(v)\,dt=1\quad(v\ne0).
\end{equation}
For $x\ne y$, let $v=y-x$ and
$L_{x,y}(u)=v\cdot(y-u)/|v|^2$. The polynomial
\[
 G_t(x,y)(u,u')=f_t(v)L_{x,y}(u)L_{x,y}(u')
 =t\lambda^2e^{-t\lambda|v|^2}
 [v\cdot(y-u)][v\cdot(y-u')]
\]
extends continuously by zero at $x=y$. If $Z$ is a standard Gaussian
vector, then differentiating its characteristic function gives
\begin{equation}\label{eq:interp-unary}
 G_t(x,y)(u,u')=\frac{\lambda}{2}\sum_{a,b}
 \E\!\left[(\delta_{ab}-Z_aZ_b)(y_a-u_a)(y_b-u'_b)
 \cos\bigl(\sqrt{2\lambda t}\,Z\cdot(y-x)\bigr)\right].
\end{equation}
Expanding the affine factors and the cosine of the difference
separates $x,y$ into coordinate monomials times sines or cosines, all
of which are bounded by one on $Q$. The total integrated coefficient
norm of these terms is at most $C$, uniformly in $t$, because
$\sum_{a,b}\E|\delta_{ab}-Z_aZ_b|<\infty$.
For these unary representations, the mark spaces are finite unions of
Gaussian spaces.

Gaussian integration also gives
\begin{equation}\label{eq:interp-spatial}
 \sup_x\int_Q f_t(y-x)\,dy\le C(1+t)^{-d/2-1},\qquad
 \sup_x\int_Qte^{-t\alpha(y-x)}\,dy\le C(1+t)^{1-d/2}.
\end{equation}
Indeed, since $\alpha(y-x)\le1/2$, we have
$e^{-t\alpha}\le e^{1/2}e^{-(1+t)\alpha}$. We then extend the integrals
to $\mathbb R^d$ and rescale by $\sqrt{\lambda(1+t)}$.

Let $q=k-1$. For each center $i$, we assign a scale $t_r\ge0$ to
every $r\ne i$ and write $S=\sum_{r\ne i}\log(1+t_r)$. We define
$E_{k,T}$ through
\begin{equation}\label{eq:interp-kernel}
 \begin{aligned}
 \phi(u)^TE_{k,T}(U)\phi(u')
 &=\sum_i\int_{S\le2\log T}
       \prod_{r\ne i}G_{t_r}(U_i,U_r)(u,u')\,d\boldsymbol t,\\
 \widehat w_i(T;U)&=\int_{S\le2\log T}
       \prod_{r\ne i}f_{t_r}(U_r-U_i)\,d\boldsymbol t.
 \end{aligned}
\end{equation}
The degree is $q\le D-1$ in each argument. The coefficients and the
weights are continuous because the scale domain is bounded. The
cardinal property of $L_{x,y}$ proves \eqref{eq:interp-diag}, and a
neighbor coincident with the center makes that center's summand zero. The range and
the monotonicity of the weights follow from
\eqref{eq:interp-partition}. Symmetry and permutation invariance are
immediate, and $E_{k,1}=0$.

Next, we multiply the separated representations
\eqref{eq:interp-unary}, integrate over the scale domain, and average
over permutations of the points. Products preserve the bound
$|b_l|\le1$. Absorbing the coefficient norms into the product Gaussian
and scale measures, we obtain a finite positive $\sigma$ and a
measurable $A$ on standard Borel spaces. With $s_r=\log(1+t_r)$, the
cost of these representations is bounded by
\begin{equation}\label{eq:interp-separated-cost}
 \frac{kC^q}{(q-1)!}\int_0^{2\log T}e^S S^{q-1}\,dS
 \le C^kT^2(1+\log T)^{k-2}.
\end{equation}
Restricting the scale domain to $2\log L_+<S\le2\log T$ gives a
representation of the band with the same bound. This proves
\eqref{eq:interp-cost}.

The remaining estimates use the fact that, for $b\ge1/2$ and
$A\ge0$,
\begin{equation}\label{eq:interp-gamma}
 \frac1{(q-1)!}\int_A^\infty e^{-bS}S^{q-1}\,dS
 =e^{-bA}\sum_{l=0}^{q-1}\frac{A^l}{l!b^{q-l}}
 \le C^q e^{-bA}(1+A)^{q-1}.
\end{equation}
The inequality follows by integration by parts, together with
$b^{-1}\le2$ and $\sum_l1/l!\le e$. Since
$\|L_{x,y}\|_{\rm coef}\le C/|y-x|$, we have
\[
 \frac{\|G_t(x,y)\|_{\rm coef}^2}{f_t(y-x)}
 \le Ct e^{-t\alpha(y-x)}.
\]
Here we define this quotient to be zero where $f_t=G_t=0$.
The measures
$d\nu_i=\prod_{r\ne i}f_{t_r}(U_r-U_i)\,d\boldsymbol t$
have total mass at most one. Applying the Cauchy--Schwarz inequality
weighted against their sum, and then \eqref{eq:interp-spatial}, we
obtain, with $A=2\log L_+$,
\[
 \int_{Q^k}\|E_{k,T}-E_{k,L}\|_{\ell^1}^2\,dU
 \le\frac{C^k}{(q-1)!}
      \int_A^\infty e^{-(d-4)S/2}S^{q-1}\,dS.
\]
In this bound, we used that the change of variables
$s_r=\log(1+t_r)$ has Jacobian $e^S$, and the $k^2$ center factors
were absorbed into $C^k$. Applying \eqref{eq:interp-gamma} with
$b=(d-4)/2\ge1/2$, we obtain \eqref{eq:interp-l2}.

Off the null set of coincident points, $1-\widehat w_i$ is the
complementary scale integral. Similarly, the first bound in
\eqref{eq:interp-spatial} gives
\[
 \int_{Q^k}\sum_i(1-\widehat w_i)\,dU
 \le\frac{C^k}{(q-1)!}
 \int_{2\log T}^\infty e^{-dS/2}S^{q-1}\,dS.
\]
Together with \eqref{eq:interp-gamma}, this implies
\eqref{eq:interp-defect}. It remains to choose fixed $C_A,C_I$ that
cover the preceding bounds.
\end{proof}
\section{Signed local updates}\label{sec:LB-local}

In this section, we work in one chart, with $u_i\in Q$, $w_i=w(u_i)$,
$p_i\in[a_M,b_M]$, $g_i\in[-\rho/2,\rho/2]$,
$c\in\mathbb B$, $V\in[v-\rho,v]$, $y_i\in\mathcal Y$ and
$0<\eta_f\le\eta_0$. We set
\[
 f_i(c)=g_i+\eta_fw_iF_c(u_i),\qquad
 \mathsf P(c)=\prod_{i\le k}q_{y_i}(f_i(c),V).
\]
All local identities below hold uniformly in these nuisance vectors,
before they are substituted from a global state in \cref{lem:score}.
We write $[k]=\{1,\ldots,k\}$ and list $\mathbf u_S$ in increasing
index order. Fix $D\ge M\ge\max\{M_1,2\}$, and let
\[
 \begin{gathered}
 c_{ab}=(a+b)/2,\quad R_{ab}=(b-a)/2,\quad
 \tau_{ab}=\operatorname{arcosh}(c_{ab}/R_{ab}),\\
 \vartheta(z)=\frac{(z-a_M)(b_M-z)}{b_M-a_M}.
 \end{gathered}
\]
Note that $0\le\vartheta(z)\le\min\{z-a_M,b_M-z\}$ for $z\in[a_M,b_M]$ and that
$\vartheta(0)=-a_Mb_M/(b_M-a_M)\ne0$. The following finite rules are
used to select counts and to extract responses. Signed measures that
represent polynomial functionals, obtained by duality from polynomial
approximation, underlie the moment-matching lower bounds of Lepski,
Nemirovski and Spokoiny \cite[Section~4.4]{LepskiEtAl1999} and of Wu
and Yang \cite{WuYang2016}; Wu and Yang
\cite[Appendix~A, Lemma~5]{WuYang2019} apply this duality to the
reciprocal on an interval away from zero. Here we select observation counts with the
rule in (a), which represents evaluation at the exterior point $0$ by a
signed measure on the density interval.

\begin{lemma}[Finite moment rules]\label{lem:finite-rules}
\textup{(a)} For $0<a<b$ and an integer $n\ge1$, at
$z_i=c_{ab}+R_{ab}\cos(i\pi/n)$, $0\le i\le n$,
let $\mathfrak l_i=\prod_{j\ne i}z_j/(z_j-z_i)$.
Then $\Lambda_n=\sum_i\mathfrak l_i\delta_{z_i}$ represents $P\mapsto P(0)$
through degree $n$, with $\|\Lambda_n\|=\cosh(n\tau_{ab})$.
Below $[a,b]=[a_M,b_M]$ and $\tau_{ab}=\tau_M$.

\textup{(b)} Let $n_D$ be the largest odd integer at most $D$,
and let $L_j$ be the cardinal polynomials at
$x_j=\cos(j\pi/n_D)$, $0\le j\le n_D$. The rule
$\mathsf J_D=\sum_jL_j'(0)\delta_{x_j}$ represents $P\mapsto P'(0)$
through degree $D$, with $\|\mathsf J_D\|=n_D\le D$ and
symmetric total variation.

\textup{(c)} At $z_j=j/(2q_{\rm rsp}+1)$, $0\le j\le2q_{\rm rsp}+1$,
let $H_j$ be the cardinal polynomials. The fixed rule
\begin{equation}\label{eq:LB-loc-nodes}
 \mathsf C_*=\sum_j([z^2]H_j(z))\delta_{z_j},\qquad
 \int\varsigma^a\,d\mathsf C_*=\1\{a=2\}\quad(0\le a\le2q_{\rm rsp}+1)
\end{equation}
has finite variation depending only on $(s,d)$.

\textup{(d)} For symmetric $A$, the fixed-atom signed kernel
\begin{equation}\label{eq:LB-covariance-measure}
 \nu_A=2C_{\rm fr}^2\sum_{\beta,\beta'}A_{\beta\beta'}
 \E_{\iota}\left[\iota_1\iota_2
 \delta_{(\iota_1e_\beta+\iota_2e_{\beta'})/(2C_{\rm fr})}\right],
\end{equation}
where $\iota_1,\iota_2$ are independent fair signs, is supported on
$\mathbb B$ and satisfies
\begin{equation}\label{eq:LB-covariance-moments}
 \int d\nu_A=0,\quad \int\mathbf v\,d\nu_A=0,\quad
 \int\mathbf v\mathbf v^\top d\nu_A=A,\quad
 \|\nu_A\|\le2C_{\rm fr}^2\|A\|_{\ell^1}.
\end{equation}
Its weights are linear in $A$.
\end{lemma}

\begin{proof}
Parts (a) and (c) follow from cardinal interpolation. The exterior
weights have sign $(-1)^{n+i}$. Applying (a) to the Chebyshev
polynomial $T_n((z-c_{ab})/R_{ab})$, whose node values are $(-1)^i$, we
find the variation in (a)
\cite[Section~2.7, Theorem~2.20 and equation~(2.37)]{Rivlin1990}. Differentiating the
interpolation formula proves (b) through degree $n_D$. The rule
vanishes on even monomials, because reflection pairs nodes with
opposite weights. This also gives exactness through degree $D$ when
$D=n_D+1$.
With $R(x)=\prod_j(x-x_j)$, the weight
$L_j'(0)=-R(0)/(x_j^2R'(x_j))$ has sign $(-1)^{(n_D-1)/2+j}$, and
exactness on $T_{n_D}$ gives the variation in (b). Finally, in (d),
averaging over the signs yields zero mass and zero mean. Each summand,
before multiplication by its prefactor, has second moment
$(e_\beta e_{\beta'}^\top+e_{\beta'}e_\beta^\top)/(4C_{\rm fr}^2)$.
Since $A$ is symmetric, this proves the covariance identity, including
repeated indices. Summing the absolute weights proves the bound, and
the fixed atoms give measurability.
\end{proof}

For a separated representation $(\mathcal Z,\sigma,\mathbf b,A)$ of $E$
as in \cref{sec:interp}, let
$\mathcal C(E)=\int\|A\|_{\ell^1}d\sigma$, and define
\begin{equation}\label{eq:LB-loc-packet}
 \Gamma_{r,m}=\frac{r^r}{r!}
 \Bigl(\frac{b_M}{\vartheta(0)}\Bigr)^r
 \Lambda_{m+r}\otimes\mathsf J_D^{\otimes r},\qquad
 \Xi_{r,m,E}=\Gamma_{r,m}\,\sigma(d\zeta)\,
                  \mathsf C_*(d\varsigma)\,\nu_{A(\zeta)}(d\mathbf v).
\end{equation}
For the density mark $e_{\rm d}=(z,\boldsymbol\epsilon,\zeta)$, the
update is
\begin{equation}\label{eq:LB-loc-update}
 p_i^e=z+\frac{\vartheta(z)p_i}{rb_M}
                    \sum_{l=1}^r\epsilon_l b_l(u_i;\zeta),\qquad
 c^e=(1-\varsigma)c+\varsigma\mathbf v.
\end{equation}
We keep the density marks separate from the response marks and
define
\begin{equation}\label{eq:LB-matrix-response}
 \mathcal B_A\Phi(c)=\int\Phi((1-\varsigma)c+\varsigma\mathbf v)
                     \,\mathsf C_*(d\varsigma)\,\nu_A(d\mathbf v).
\end{equation}
Fix a sufficiently large response constant $C_q\ge1$, as required
below, and let $C_{\rm hi}=12p_+C_q$.

\begin{lemma}[Local operator realization]\label{lem:local-operator}
Let $1\le r\le m$ and $r\le D$, with symmetric $A$ and finite
representation cost. The resets are measurable and preserve the
density interval and coefficient ball; the density slope has modulus
at most $(b_M-a_M)/(4b_M)<1/4$. Their weights are independent of
the incoming state. For every bounded measurable $b$,
\begin{equation}\label{eq:LB-operator-annihilation}
 \int b(e_{\rm d})\,d\Xi_{r,m,E}=0,\qquad
 \int b(e_{\rm d})c^e\,d\Xi_{r,m,E}=0.
\end{equation}
Moreover $\|\Gamma_{r,m}\|\le C^rD^re^{\tau_Mm}$ and
$\|\Xi_{r,m,E}\|\le C^rD^re^{\tau_Mm}\mathcal C(E)$.

For every bounded measurable $\Phi$ on $\mathbb B$ and $k\le D$,
\begin{equation}\label{eq:LB-loc-row}
 \int\prod_{i\le k}p_i^e\,\Phi(c^e)\,d\Xi_{r,m,E}
 =\mathfrak A_k^{r,m}\sum_{|S|=r}\prod_{i\in S}p_i\,
                  \mathcal B_{E(\mathbf u_S)}\Phi(c),
\end{equation}
with zero right side for $k<r$ and
\begin{equation}\label{eq:LB-loc-count}
 \begin{aligned}
 \mathfrak A_k^{r,m}&=\vartheta(0)^{-r}
 \int z^{k-r}\vartheta(z)^r\,d\Lambda_{m+r}(z)\quad(k\ge r),\\
 \mathfrak A_k^{r,m}&=\1\{k=r\}\quad(r\le k\le m),\qquad
 |\mathfrak A_k^{r,m}|\le C^ke^{\tau_Mm}\quad(k\ge r).
 \end{aligned}
\end{equation}

For every polynomial $\Phi$ of degree at most two,
$\mathcal B_A\Phi(c)=\tfrac12\operatorname{tr}(A D^2\Phi(c))$.
For the response product $\mathsf P$, define
\[
 \mathcal Q_A(c)=\mathcal B_A\mathsf P(c)
              -\tfrac12\operatorname{tr}(A D^2\mathsf P(c)).
\]
Then
\begin{equation}\label{eq:LB-response-error}
 |\mathcal Q_A|\le Ck^{2q_{\rm rsp}+2}\eta_f^{2q_{\rm rsp}+2}
                         \|A\|_{\ell^1},\qquad
 \mathcal Q_A=0\quad(k\le q_{\rm rsp}).
\end{equation}
If $[\phi(u_i)^TA\phi(u_l)]_{i,l}$ is diagonal with entries $d_i$,
\begin{equation}\label{eq:LB-loc-heat}
 \mathcal B_A\mathsf P(c)
 =\eta_f^2\sum_iw_i^2d_i\partial_{V_i}\mathsf P(c)+\mathcal Q_A(c).
\end{equation}
These response assertions hold for every symmetric $A$ and every count.
Uniformly for all $k\ge1$,
\begin{equation}\label{eq:LB-matrix-response-bound}
 |\partial_{V_i}\mathsf P|\le C_q,\qquad
 |\mathcal B_A\mathsf P(c)|\le Ck^2\eta_f^2\|A\|_{\ell^1}.
\end{equation}
More generally, any finite signed update measure $\Omega$ with legal
resets, density resets measurable in $e_{\rm d}$ alone, and the conditional annihilation in
\eqref{eq:LB-operator-annihilation} satisfies, for every count,
\begin{equation}\label{eq:LB-all-count-update}
 \left|\int\prod_{i\le k}p_i^e\,\mathsf P(c^e)\,d\Omega\right|
 \le2C_qk^2p_+^k\eta_f^2\|\Omega\|.
\end{equation}
All constants are independent of the frame dimension and counts.
\end{lemma}

\begin{proof}
The averaged modulation has modulus at most one, so the quadratic
margin keeps the density admissible, and the slope bound follows from
the maximum of the margin. The coefficient reset is a convex combination of two points in the
coefficient ball. Conditional
annihilation follows from the covariance rule. Multiplying the
variations, and using $r^r/r!\le e^r$ and the fixed rules, we obtain
the cost bounds, which also justify every application of Fubini's
theorem.

For $k\le D$, the density product has degree at most $D$ in each
$\epsilon_l$. Applying $\mathsf J_D$ in each coordinate selects
$r$ distinct observations and $r!$ assignments, and we obtain
\[
 \frac{r!}{r^r}\Bigl(\frac{\vartheta(z)}{b_M}\Bigr)^rz^{k-r}
 \sum_{|S|=r}\prod_{i\in S}p_i
 \Sym\!\left[\prod_l b_l(u_{S_l};\zeta)\right].
\]
Then \eqref{eq:LB-loc-row} follows from the separated
representation and the linearity of $\mathcal B_A$. The exterior
evaluation is exact when $k+r\le m+r$. For the other counts, the
variation bound and the identity
$\max\vartheta/|\vartheta(0)|=\sinh^{-2}\tau_M$ give
$|\mathfrak A_k^{r,m}|\le b_M^{k-r}\sinh^{-2r}(\tau_M)e^{\tau_M(m+r)}$.
Since $M\ge M_1$, \eqref{eq:LB-tauM} shows that $\tau_M$ is bounded by a
fixed constant and is at least $\tau$, which implies \eqref{eq:LB-loc-count}.

For a quadratic $\Phi$, the fixed response rule extracts the quadratic
coefficient of $\Phi(c+z(\mathbf v-c))$. The zero mass and mean of
$\nu_A$ leave only
$\tfrac12\int D^2\Phi(c)[\mathbf v,\mathbf v]\,d\nu_A(\mathbf v)$,
which is $\tfrac12\operatorname{tr}(A D^2\Phi(c))$.
For the response product, apply the same scalar rule to the Taylor
polynomial through degree $2q_{\rm rsp}+1$ of
$z\mapsto\mathsf P(c+z(\mathbf v-c))$.
Its quadratic coefficient, integrated against $\nu_A$, is
$\tfrac12\operatorname{tr}(A D^2\mathsf P(c))$.
The derivative of order $a=2q_{\rm rsp}+2$ along the reset is at most
$C_a k^a\eta_f^a$, since each factor has degree two and
$|F_{\mathbf v}-F_c|\le2$. The Taylor remainder and the variation
bound for $\nu_A$ prove \eqref{eq:LB-response-error}.
Since this scalar product polynomial has degree at most $2k$, the
discrepancy vanishes when $k\le q_{\rm rsp}$.
The heat identity follows from the chain rule and
$\partial_f^2q_y=2\partial_Vq_y$.

Finally, let $\Psi(\mathbf z)=\prod_iq_{y_i}(g_i+\eta_fw_i z_i,V)$
on $[-1,1]^k$. Since the ternary derivatives are bounded, we have
$|\partial_i\partial_l\Psi|\le C_q\eta_f^2$ and
$|\partial_{V_i}\mathsf P|\le C_q$. The remainder of the affine Taylor
approximation between any two field vectors is then at most
$2C_qk^2\eta_f^2$. The affine part is annihilated conditionally on the
density mark, and this proves both the matrix response bound and
\eqref{eq:LB-all-count-update}, including counts beyond $D$.
\end{proof}

\section{Cancellation for two observations}\label{sec:pair-new}

The unit and coarse kernels use the ordinary selector with $r=2$.
The fine kernel uses a three-point rule whose variation is independent
of $D$, and its higher even modulation terms are controlled by a
bounded field in the design density. The regression updates remain
smooth.

\begin{lemma}[A bounded hyperplane field]\label{lem:hyperplane}
For $T\ge0$ there is a jointly measurable field $\psi_T:Q\to\{-1,1\}$,
with a law measurable in $T$, such that
$\E\psi_T(u)\psi_T(v)=e^{-T|u-v|}$. Its odd moments vanish. Writing
$M_{T,j}(u_1,\ldots,u_j)=\E\prod_i\psi_T(u_i)$, for every even $j\ge4$,
\begin{equation}\label{eq:hp-moment}
 \left\|\E\prod_{i=1}^j\psi_T(u_i)\right\|_{L^2(Q^j)}
 \le C^j j^4T^{-d}\qquad(T>0).
\end{equation}
\end{lemma}

\begin{proof}
Let $\sigma$ be the uniform probability measure on $S^{d-1}$, and let
$\gamma_d=\int|v_1|\,d\sigma(v)$. We draw Poisson hyperplanes
\cite[Section~12.4]{LastPenrose2018} with
intensity $(T/\gamma_d)\sigma(dv)\,db$ on
$S^{d-1}\times[-\sqrt d,\sqrt d]$, and then give their cells
independent fair signs. There are finitely many hyperplanes. A mark
consists of their ordered list and a sign array indexed by the
halfspace patterns. With a fixed boundary convention, the evaluation
is Borel on a countable union of finite-dimensional spaces. Only the
law of the Poisson count depends on $T$.

Conditionally on the hyperplanes, averaging over the signs leaves one
precisely when every occupied cell contains an even number of the
indexed points. Two points share a cell when no hyperplane separates
them, and the intensity of the hyperplanes that separate them is
$(T/\gamma_d)\int|v\cdot(u-u')|\,d\sigma(v)=T|u-u'|$.
This implies the covariance and odd-moment statements. A related
representation of the exponential kernel, through the probability that
two points lie in a common cell, is derived for STIT tessellations by
O'Reilly and Tran \cite[Section~4.1, Theorem~4.1]{OReillyTran2022}.
Such an occupancy partition of $j\ge4$ points contains two disjoint
pairs of indices, with the points of each pair lying in the same cell. For any two chosen pairs, the union of their intervals of
separating offsets has length at least half the sum of their lengths.
The probability that each of the two pairs lies in a common cell is
therefore at most $e^{-T(|u_a-u_b|+|u_c-u_e|)/2}$, and two independent
spatial integrations show that the $L^2(Q^j)$ norm of this bound is at
most $C^jT^{-d}$. A union bound over the at most $j^4$ choices of
pairs proves \eqref{eq:hp-moment}.
\end{proof}

Let $T_0=Ne^{-c_0D}$, where $c_0=\tau+1$, and suppose that $T_0>1$.
For $D\ge3$, we choose fixed symmetric matrices $A_0,A_1$ with
kernels
\[
 \phi(u)^TA_0\phi(v)=1,\qquad
 \phi(u)^TA_1\phi(v)=-|u-v|^2.
\]
Indeed, $A_0=e_0e_0^\top$, and $-|u-v|^2=2u\cdot v-|u|^2-|v|^2$.
Their coefficient norms depend only on $d$. Let
$\mathrm{Rsp}(A)=\mathsf C_*\otimes\nu_A$, let $\mathcal L_T$ be the
law of the marked field, and define
\begin{equation}\label{eq:new-pair-packet}
 \begin{gathered}
 X=\tfrac12(\delta_{-1}+\delta_1)-\delta_0,\qquad
 \Gamma_T^{\rm fi}
 =\frac{b_M^2}{\vartheta(0)^2}
   \Lambda_{M+2}(dz)X(d\epsilon)\mathcal L_T(d\psi),\\
 p_i^e=z+\frac{\vartheta(z)}{b_M}p_i\epsilon\psi(u_i).
 \end{gathered}
\end{equation}
Note that $\|X\|=2$ and that its moments vanish at degree zero and
at every odd degree and equal one at every positive even degree. The fine
update stays in $[a_M,b_M]$ and has slope less than $1/4$, by the same
margin bound as for \eqref{eq:LB-loc-update}. We use three
separate row tags:
\begin{equation}\label{eq:new-pair-rows}
 \begin{gathered}
 \Xi_{\rm unit}=\Gamma_{2,D}\otimes\mathrm{Rsp}(A_0),\qquad
 \Xi_{\rm co}=\int_0^{T_0}T\,
   \Gamma_{2,D}\otimes\mathcal L_T\otimes\mathrm{Rsp}(A_1)\,dT,
 \\
 \Xi_{\rm fi}=\int_{T_0}^NT\,
   \Gamma_T^{\rm fi}\otimes\mathrm{Rsp}(A_1)\,dT.
 \end{gathered}
\end{equation}
For the first two rows, we use the ordinary update
\eqref{eq:LB-loc-update} with $b_1=b_2=1$ and $b_1=b_2=\psi$,
respectively. Their extraction and admissibility follow from
\cref{lem:local-operator}. We include $T$ in the mark, so that the
displayed integrals are finite signed measures on standard Borel
spaces. We denote their direct sum by $\Xi_{\rm pair}$.

\begin{lemma}[Pair activity and cancellation]\label{lem:new-pair}
In $\mathsf{Reg}(K)$, for sufficiently large $M$, the sum of these rows
has variation at most $CN^2e^{\tau_MM}$ and annihilates counts zero
and one. At count two, after subtracting the diagonal variance
derivative, its numerator is
\begin{equation}\label{eq:new-pair-defect}
 \eta_f^2w_1w_2p_1p_2(1+Nr)e^{-Nr}
 \partial_{f_1}\partial_{f_2}\mathsf P,
 \qquad r=|u_1-u_2|.
\end{equation}
This expression has squared norm at most $C\eta_f^4N^{-d}$ for every
$N\ge1$. At $3\le k\le D$ its contribution is
$\mathsf U_k+\mathsf V_k$, where $\mathsf U_k=0$ for $k\le M$ and
\begin{align}
 \sum_{k=M+1}^D\frac{(C_\sharp\mu)^k}{k!}\|\mathsf U_k\|_{(k)}^2
 &\le C\eta_f^4e^{2\tau_MM}N^{4-d}
       \frac{(C\mu)^{M+1}}{(M+1)!},\label{eq:new-pair-alias}\\
 \sum_{k=3}^D\frac{(C_\sharp\mu)^k}{k!}\|\mathsf V_k\|_{(k)}^2
 &\le C\eta_f^4e^{2\tau_MM}\mu^4T_0^{4-2d}
 \le C\eta_f^4\mu^2N^{-d}.\label{eq:new-pair-field}
\end{align}
\end{lemma}

\begin{proof}
The variations of the ordinary and fine packets give
\[
 C\{D^2e^{\tau_MD}(1+T_0^2)+N^2e^{\tau_MM}\}
 \le CN^2e^{\tau_MM}:
\]
eventually we have $\tau_M\le c_0$, so that the coarse term is at
most $CD^2N^2e^{-c_0D}$, and the unit term is absorbed because
$\log N\asymp M^2$.

For the ordinary pair selector, we write
$a_{k,m}=\mathfrak A_k^{2,m}$. This coefficient equals one at $k=2$,
vanishes for $3\le k\le m$, and is bounded by $C^ke^{\tau_Mm}$.
Expanding the fine density product, we obtain the same selected pair
term and, for even $4\le j\le k$, the coefficients
\[
 c_{k,j}=b_M^{2-j}\vartheta(0)^{-2}
 \int z^{k-j}\vartheta(z)^j\,d\Lambda_{M+2}(z),
 \qquad |c_{k,j}|\le C^ke^{\tau_MM}.
\]
For the fine row, the vanishing zeroth and first moments of $X$ prove
annihilation at counts zero and one. The ordinary rows annihilate these
counts by count selection. At count two, the response operator is
$p_1p_2(\mathcal B_{A_0}+W_N(r)\mathcal B_{A_1})\mathsf P$, where
\[
 W_N(r)=\int_0^NT e^{-Tr}\,dT
 =\frac{1-(1+Nr)e^{-Nr}}{r^2},\qquad W_N(0)=N^2/2.
\]
At count two, $q_{\rm rsp}\ge2$ and
\eqref{eq:LB-response-error} give exact extraction of the Hessian. The diagonal contribution is the heat derivative, and the cross entry
of the evaluated kernel is $1-r^2W_N(r)=(1+Nr)e^{-Nr}$, which proves
\eqref{eq:new-pair-defect}. The norm bound follows from
$\int_{\mathbb R^d}(1+N|x|)^2e^{-2N|x|}\,dx=CN^{-d}$, and this radial
calculation itself holds for every $N\ge1$.

For higher counts, we define
$W^{\rm fi}(r)=\int_{T_0}^NT e^{-Tr}\,dT$ and
\[
 \mathsf U_k=a_{k,M}\sum_{i<l}p_ip_lW^{\rm fi}(|u_i-u_l|)
 \mathcal B_{A_1}\mathsf P.
\]
Since $W^{\rm fi}(r)\le r^{-2}(1+T_0r)e^{-T_0r}$ and $d>4$, we have
$\|W^{\rm fi}(|\cdot-\cdot|)\|_2\le CT_0^{2-d/2}$. Using the full
response bound and the bounded densities, and summing over pairs and
responses, we find that
$\|\mathsf U_k\|_{(k)}\le C^k\eta_f^2e^{\tau_MM}T_0^{2-d/2}$.
The factorial tail starts at $M+1$. Replacing $T_0^{4-d}$ by
$N^{4-d}$ costs a factor $e^{c_0(d-4)D}$, which is absorbed into the
fixed $C^{M+1}$ because $D=O(M)$. This proves
\eqref{eq:new-pair-alias}.

For even $j\ge4$, let $H_j=\int_{T_0}^NT M_{T,j}\,dT$. By
Minkowski's inequality and \eqref{eq:hp-moment}, we have
$\|H_j\|_2\le C^jj^4T_0^{2-d}$. The remaining numerator is
\[
 \mathsf V_k=\sum_{\substack{4\le j\le k\\j\text{ even}}}
 c_{k,j}\sum_{|S|=j}\prod_{i\in S}p_i\,
 H_j(\mathbf u_S)\mathcal B_{A_1}\mathsf P.
\]
It follows that
$\|\mathsf V_k\|_{(k)}\le C^k\eta_f^2e^{\tau_MM}T_0^{2-d}$, and this
term vanishes for $k<4$. Summing these bounds gives the first
inequality in \eqref{eq:new-pair-field}. Up to a fixed multiplicative constant, the ratio of its right side
to $\eta_f^4\mu^2N^{-d}$ is
$e^{2\tau_MM}\mu^2N^{4-d}e^{c_0(2d-4)D}=o(1)$, since
$\log N\asymp M^2$, $D=O(M)$ and $d>4$.
\end{proof}

\medskip\noindent\emph{One cutoff for all target counts.}
Recall that $c_0=\tau+1$. For $N\ge1$ and $\mu\in(0,1)$, let
$H=1+\log N$, and define
\begin{equation}\label{def:rows}
\begin{gathered}
 \ell_\circ=Ne^{-c_0D},\qquad
 \nu_r=N(\mu H)^{(r-2)/d}\quad(2\le r\le D),\\
 E_r^{\rm co}=E_{r,\min\{\nu_r,\ell_\circ\}},\qquad
 E_r^{\rm fi}=E_{r,\nu_r}-E_r^{\rm co},\qquad
 m_r^{\rm co}=D,\quad m_r^{\rm fi}=M.
\end{gathered}
\end{equation}
As before, $E_{r,T}=0$ for $T<1$. We omit the cases $\nu_r<1$ and
the zero bands. For sufficiently large $M$ in $\mathsf{Reg}(K)$, we
have $\mu H\le1$ and $\ell_\circ>1$, and every nonzero fine band has
$r\le M$. The scale calculation in \cref{sec:LB-alias} proves these
guard conditions. The coarse guards equal $D$, so all aliases start at
count $M+1$.

For each selected band with $r\ge3$ and $b\in\{{\rm co},{\rm fi}\}$,
we fix a separated representation from \cref{lem:interp} with
symmetric $A(\zeta)$. For $r=1$, we take $b={\rm co}$, the constant
matrix $E_1=e_0e_0^T$ and the guard $D$. Its representation has one
point, with $b_1\equiv1$ and $A=e_0e_0^T$. For each such
representation, we form the operator of \cref{lem:local-operator}
with its prescribed guard:
\begin{equation}\label{eq:LB-row-measure}
 \Xi_{r,b}=\Xi_{r,m_r^b,E_r^b},
\end{equation}
where $E_1^{\rm co}=E_1$ and $m_1^{\rm co}=D$.

At target count two, the row is $\Xi_{\rm pair}$ from
\eqref{eq:new-pair-rows}. The singleton, the targets $r\ge3$ and the
three pair components have distinct tags. Let $E_0$ be the finite
disjoint union of their mark spaces, and define
\begin{equation}\label{eq:LB-signed-measure}
 \Xi=\Xi_{\rm pair}\oplus\bigoplus_{r\ne2,b}\Xi_{r,b},\qquad B_{\rm base}=\|\Xi\|.
\end{equation}
We also denote the component at target $r$ and band $b$ by
$\mathrm{Row}_r^b$. All mark spaces are standard Borel, all updates
are measurable and admissible, and the singleton measure is nonzero.
Hence $0<B_{\rm base}<\infty$.

We write $e_{\rm d}$ for the density mark, including its row tag. By
the conditional annihilation in \cref{lem:local-operator}, which also
holds for the fine pair row, we have, for every bounded
measurable $b$ and $c\in\mathbb B$,
\begin{equation}\label{eq:LB-annihilation}
 \int b(e_{\rm d})\,\Xi(de)=0,\qquad
 \int b(e_{\rm d})c^e\,\Xi(de)=0.
\end{equation}

\medskip\noindent\emph{A centered reference law for the same measure.}
The generator needs the signed measure $\Xi$, but the construction of
the prior also needs a positive mark law. We start with
$\pi_0=|\Xi|/B_{\rm base}$ and the weight
$\mathsf a_0=B_{\rm base}\,d\Xi/d|\Xi|$, where we choose a measurable
$\{-1,1\}$-valued version of the derivative. Then
$\Xi=\mathsf a_0\pi_0$ and $|\mathsf a_0|=B_{\rm base}$. Note that marks
with zero weight may be added without changing any signed integral. We use
this freedom to make the reference density have mean one. Let
\begin{equation}\label{eq:LB-centering-constants}
 r_{\rm ctr}=\tfrac12\min\{1-p_-,p_+-1\},\qquad
 \alpha_{\rm ctr}=\frac{r_{\rm ctr}}{2(p_+-p_-)},\qquad
 C_{\rm ctr}=\alpha_{\rm ctr}^{-1}.
\end{equation}

\begin{lemma}[The reference marks]\label{lem:centered-row}
For $M\ge M_1'$, extend $E_0$ by one zero-weight branch.
There is a probability law $\pi$ on the resulting mark space $E$
and a state-independent weight $\mathsf a^\circ$ such that
\begin{equation}\label{eq:LB-Bloc}
 \Xi=\mathsf a^\circ\pi,\qquad
 B_{\rm loc}:=\sup_E|\mathsf a^\circ|=C_{\rm ctr}B_{\rm base}.
\end{equation}
All updates remain measurable and admissible, with pointwise affine
density form $p_i^e=A(e_{\rm d},u_i)p_i+B(e_{\rm d},u_i)$, $|A|\le1$.
Their reference means satisfy
\begin{equation}\label{eq:LB-row-centered}
 \int A(e_{\rm d},u)\,\pi(de)=0,\qquad
 \int B(e_{\rm d},u)\,\pi(de)=1.
\end{equation}
Thus the expected density after an update is one for every incoming state.
\end{lemma}

\begin{proof}
Under $\pi_0$, the slope has mean zero, since the variation of each
ordinary packet factors over symmetric $|\mathsf J_D|$ modulations. The
fine pair packet has the same property because $|X|$ is symmetric. The
intercept has a constant mean $\bar z\in[a_M,b_M]$. With probability
$\alpha_{\rm ctr}$, we draw the mark from $\pi_0$ and multiply its
weight by
$C_{\rm ctr}$. Otherwise, we reset the density to
\[
 z_0=\frac{1-\alpha_{\rm ctr}\bar z}{1-\alpha_{\rm ctr}},
\]
leave $c$ fixed, and assign weight zero. Since both margins of the
interval around one are at least $r_{\rm ctr}$, the quantities
$(1-\alpha_{\rm ctr})(z_0-a_M)$ and
$(1-\alpha_{\rm ctr})(b_M-z_0)$ are at least $r_{\rm ctr}/2$, and the
reset is admissible. Under $\pi$, the slope has mean zero and the
intercept has mean one. Only the original branch has nonzero weight,
which proves \eqref{eq:LB-Bloc} and preserves \eqref{eq:LB-annihilation}.
All formulas are measurable, and the weight is independent of the
state.
\end{proof}

\section{The generator and its local cancellation}\label{sec:LB-rows}

We now transport these concrete updates to the patches. The update
$\mathsf T_{j,e}$ replaces $p_\vartheta(x)$ on $Q_j$ by
$A(e_{\rm d},\chi_j(x))p_\vartheta(x)+B(e_{\rm d},\chi_j(x))$
and $c_j$ by $c_j^e$, and it leaves everything else unchanged. It
preserves raw states. Updates at disjoint patches commute, since neither
changes the other's incoming state. For a bounded function $G$ of the
raw state, whenever the mark integrands are measurable, we define
\begin{equation}\label{eq:LB-generator}
 (\mathcal A_jG)(\vartheta)=\int G(\mathsf T_{j,e}\vartheta)\,\Xi(de)
 =\int G(\mathsf T_{j,e}\vartheta)\mathsf a^\circ(e)\,\pi(de),
 \qquad \mathcal A=\sum_{j\in\mathcal J}\mathcal A_j.
\end{equation}
In particular, we have $\|\mathcal A_jG\|_\infty\le B_{\rm loc}\|G\|_\infty$.
For the history notation used later, we also set
$\mathsf a_j(\vartheta,e)=\mathsf a^\circ(e)$.
Every mark integrand used below is measurable by the displayed update
formulas, and \cref{sec:LB-heaps} makes this explicit for finite
histories.

\emph{Local score.} The
operator acts on the raw product $\mathsf U$, before normalization.
For $\xi=((x_i,y_i))_{i\le n}$, we define $\partial_{V_i}\mathsf U$ by
replacing its $i$th response factor by $\partial_Vq_{y_i}$. The
\emph{score numerator} and the \emph{local score} are
\begin{equation}\label{eq:LB-score}
 \mathsf R_j=m_\vartheta^{-n}\left[
 \mathcal A_j\mathsf U(\cdot,V;\xi)(\vartheta)
 -\eta_f^2\sum_{i\le n}w_j(x_i)^2\partial_{V_i}\mathsf U(\vartheta,V;\xi)
 \right],
 \qquad \mathsf r_j=\mathsf R_j/\mathsf L.
\end{equation}
Here the suppressed arguments are $(\vartheta,V;\xi)$. Since
$\sum_jw_j^2=1$ and $\sum_i\partial_{V_i}\mathsf U=\partial_V\mathsf U$,
we have
\begin{equation}\label{eq:LB-score-sum}
 \sum_j\mathsf R_j
 =m_\vartheta^{-n}\left[
 \mathcal A\mathsf U(\cdot,V;\xi)(\vartheta)
 -\eta_f^2\partial_V\mathsf U(\vartheta,V;\xi)\right].
\end{equation}
We will see that the mass tilt in \cref{sec:LB-heaps} identifies the
derivative of the mixture with the prior average of this sum.

\emph{Local likelihood and energy.} For the finite vectors and
coefficients fixed in \cref{sec:LB-local}, we define
\begin{align}
 \mathsf L_k(\mathbf p,c,\mathbf g,V;\mathbf u,\mathbf y)
 &=\prod_{i\le k}p_i\,q_{y_i}(f_i(c),V),
 \label{eq:LB-chart}\\
 \mathsf R_k(\mathbf p,c,\mathbf g,V;\mathbf u,\mathbf y)
 &=\int_E\mathsf L_k(\mathbf p^e,c^e,\ldots)\,\Xi(de)
 -\eta_f^2\sum_{i\le k}w_i^2\,
 \partial_{V_i}\mathsf L_k(\mathbf p,c,\ldots),
 \label{eq:LB-numerator}
\end{align}
where the dots stand for $(\mathbf g,V;\mathbf u,\mathbf y)$,
$\mathsf L_0=1$, and
$\partial_{V_i}\mathsf L_k=\prod_lp_l\cdot(\partial_{V_i}\mathsf P)(c)$.
Since $|f_i(c)|\le\rho/2+\eta_0\le\rho$, \cref{sec:prelim-ternary}
implies that
\begin{equation}\label{eq:LB-chart-bounds}
 (a_Mc_q)^k\le\mathsf L_k(\mathbf p,c,\mathbf g,V;\mathbf u,\mathbf y)
 \le b_M^k.
\end{equation}
For a jointly Borel function $G$ that is continuous in the
finite-dimensional nuisance variables $(\mathbf p,c,\mathbf g,V)$, we define
\begin{equation}\label{eq:LB-norm}
 \norm G_{(k)}^2
 =\int_{Q^k}\sup_{\mathbf p,c,\mathbf g,V}
 \sum_{\mathbf y\in\mathcal Y^k}
 |G(\mathbf p,c,\mathbf g,V;\mathbf u,\mathbf y)|^2\,d\mathbf u.
\end{equation}
The supremum is taken over the compact ranges stated in
\cref{sec:LB-local}. It equals the supremum over a countable dense
subset, so the integrand is measurable. This norm applies to all
numerators below, since they are polynomials in these nuisance variables
with bounded Borel coefficients; the finite variation of $\Xi$ justifies
integrating their coefficients. Because the supremum is taken before the
integration, the nuisance values may be arbitrary at each tuple
$\mathbf u$. The triangle inequality for this norm follows from the
pointwise supremum and Minkowski's inequality.

\begin{lemma}[Chart scores and conditional centering]\label{lem:score}
For the constructed update measure, let $\vartheta\in\Theta_{\rm raw}$, $V\in[v-\rho,v]$ and
$j\in\mathcal J$.

\textup{(a)} Suppose $\xi_{Q_j}$ consists of $k$ observations, with chart
coordinates $\mathbf u\in Q^k$ and responses $\mathbf y\in\mathcal Y^k$.
Relabel these points as $x_i=\chi_j^{-1}(u_i)$ and substitute
\begin{equation}\label{eq:LB-background}
 p_i=p_\vartheta(x_i),\qquad c=c_j,\qquad
 g_i=f_\vartheta(x_i)-\eta_fw_j(x_i)F_{c_j}(u_i).
\end{equation}
These are admissible nuisance values, and
\begin{equation}\label{eq:LB-chart-score}
 \mathsf r_j=\frac{\mathsf R_k}{\mathsf L_k},\qquad
 |\mathsf r_j|\le B_{\rm loc}\Bigl(\frac{b_M}{a_Mc_q}\Bigr)^k
 +\frac{k\eta_f^2}{a_Y^2c_q}.
\end{equation}
In particular the score depends on the data only through $\xi_{Q_j}$.

\textup{(b)} $\mathsf R_0=0$ and
$\sum_{\mathbf y}\mathsf R_k(\mathbf p,c,\mathbf g,V;\mathbf u,\mathbf y)=0$
for every $k$, nuisance vector and design tuple.

\textup{(c)} The scores are square integrable and
$\E_{\vartheta,V}(\mathsf r_j\mid X)=0$, where
$X=(x_1,\ldots,x_n)$. If $j\ne j'$ are not adjacent, then
$\E_{\vartheta,V}(\mathsf r_j\mathsf r_{j'}\mid X)=0$.
\end{lemma}

\begin{proof}
The background is a sum of other profiles, so
$|g_i|\le(2^d-1)\eta_f\le\rho/2$, and the update does not change it.
With $I=\{i:x_i\in Q_j\}$, the raw likelihood factors as
\begin{equation}\label{eq:LB-lik-split}
 \mathsf U=O\mathsf L_k,\qquad
 O=\prod_{i\notin I}p_\vartheta(x_i)q_{y_i}(f_\vartheta(x_i),V).
\end{equation}
Note that the update leaves $O$ fixed and that $w_j=0$ outside the
patch. We conclude
that $\mathsf R_j=m_\vartheta^{-n}O\mathsf R_k$ and
$\mathsf L=m_\vartheta^{-n}O\mathsf L_k$, which proves the formula for
the ratio. The likelihood bounds and $|\partial_Vq_y|\le a_Y^{-2}$ give
the envelope in (a).
When we sum over responses, the integral term reduces to
$\int\prod_i p_i^e\,d\Xi=0$ by \eqref{eq:LB-annihilation}, and the
variance term also sums to zero. This argument includes $k=0$ and
proves (b). Conditional on $X$, the response probability is
$\mathsf L_k/\prod_i p_i$, and (b) implies that the score has conditional
mean zero. Disjoint patches involve disjoint sets of responses, which
are independent given $X$; the scores of such patches therefore have
zero conditional cross moment. Finally, the envelope in (a), with
$k\le n$, gives square integrability.
\end{proof}

\medskip\noindent\emph{The full response remainder.}
The full chart numerator also satisfies, for every $k\ge1$,
\begin{equation}\label{eq:LB-full-response}
 |\mathsf R_k(\mathbf p,c,\mathbf g,V;\mathbf u,\mathbf y)|
 \le C_{\rm hi}^k\eta_f^2(B_{\rm loc}+1).
\end{equation}

Indeed, \eqref{eq:LB-all-count-update} and the bound on the variance
derivative give the bound $3C_qk^2p_+^k\eta_f^2(B_{\rm loc}+1)$, and it
remains to use $k^2\le4^k$.

For $2\le k\le D$ and $\mathbf u\in Q^k$, we define
$\hat w_i^{[k]}(\mathbf u)=\hat w_i(\nu_k;\mathbf u)$, with value zero
when $\nu_k<1$. Recall that $C_\sharp=1/(p_-^2c_q)$.

\begin{proposition}[Local cancellation and energy]\label{prop:local}
Let $K\ge1$. There are $M_0(K)$ and constants $C_B$, $C_5$, $C_6$, $C_7$,
depending only on $(s,d,p_-,p_+,L_s,v_-,v_+,C_4)$, on the design
constants fixed in \cref{sec:LB-proof} and on $K$, such that for
every $(M,D,N,\mu,\eta_f)\in\mathsf{Reg}(K)$ with $M\ge M_0(K)$ the row
defined by \eqref{def:rows} and \eqref{eq:new-pair-rows} is well defined and has the following properties.
\begin{enumerate}
\item[\textup{(i)}] The update measure preserves raw states and has the
state-independent representation $\Xi=\mathsf a^\circ\pi$ constructed
above. It obeys \eqref{eq:LB-annihilation} and the centering identity
\eqref{eq:LB-row-centered}.
\item[\textup{(ii)}] $B_{\rm loc}\le C_BN^2e^{\tau_MM}$.
\item[\textup{(iii)}] $\mathsf R_0\equiv0$ and $\mathsf R_1\equiv0$.
\item[\textup{(iv)}] For $2\le k\le D$, uniformly in the prescribed
nuisance values, $\mathsf R_k=\mathsf R_k^{\rm def}
+\mathsf R_k^{\rm al}+\mathsf R_k^{\rm fld}$.
At $k=2$, the defect is \eqref{eq:new-pair-defect} and the other terms
vanish. For $3\le k\le D$, put
\begin{align*}
 \mathsf R_k^{\rm def}
 &=\prod_{l\le k}p_l\left\{
 \mathcal B_{E_{k,\nu_k}(\mathbf u)}\mathsf P(c)
 -\eta_f^2\sum_{i\le k}w(u_i)^2(\partial_{V_i}\mathsf P)(c)
 \right\},\\
 \mathsf R_k^{\rm al}
 &=\mathsf U_k+
 \sum_{\substack{3\le r\le\min\{k,D\}:\,E_r^{\rm fi}\ne0\\M<k}}
 \mathfrak A^{r,M}_k
 \sum_{S\subset[k],\,|S|=r}\prod_{l\in S}p_l\,
 \mathcal B_{E_r^{\rm fi}(\mathbf u_S)}\mathsf P(c),\\
 \mathsf R_k^{\rm fld}&=\mathsf V_k,
\end{align*}
where $\mathsf U_k,\mathsf V_k$ are the pair alias and field terms
in \cref{lem:new-pair}. The defect at $k\ge3$ includes the
interpolation error and the fixed response rule's Taylor remainder.
\item[\textup{(v)}] $\sum_{k=2}^D\frac{(C_\sharp\mu)^k}{k!}
\norm{\mathsf R_k^{\rm def}}_{(k)}^2\le C_5\,\eta_f^4\,\mu^2N^{-d}$.
\item[\textup{(vi)}] $\mathsf R_k^{\rm al}=0$ for $2\le k\le M$, and
$\sum_{k=M+1}^D\frac{(C_\sharp\mu)^k}{k!}
\norm{\mathsf R_k^{\rm al}}_{(k)}^2
\le C_6\,\eta_f^4\,e^{2\tau_MM}N^{4-d}\,\frac{(C_6\mu)^{M+1}}{(M+1)!}$.
\item[\textup{(vii)}]
$\sum_{k=3}^D(C_\sharp\mu)^k\|\mathsf R_k^{\rm fld}\|_{(k)}^2/k!
\le C_7\eta_f^4\mu^2N^{-d}$.
\end{enumerate}
\end{proposition}

Part (i) follows from the construction and \cref{lem:centered-row}.
The bounds on the activity and the aliases are proved in
\cref{lem:activity,lem:alias}, and the field bound is
\eqref{eq:new-pair-field}. For the count decomposition, we note that
centering preserves signed likelihood integrals. The score numerators at counts zero and one
vanish: only the singleton survives at count one, and its exact heat
response cancels the variance derivative. At count two, all other
targets vanish, and we are left with \eqref{eq:new-pair-defect}.

The remaining count decomposition follows directly from the
ordinary-row identity \eqref{eq:LB-loc-row}.
For $3\le k\le D$, the target-$k$ bands sum to $E_{k,\nu_k}$, and every
other coarse target vanishes. Every other fine target vanishes unless $k>M$; in
that case, their targets satisfy $r\le M<k$, and they give exactly the
stated aliases. The pair rows supply $\mathsf U_k+\mathsf V_k$.
Subtracting the variance derivative proves (iii)--(iv), including the
omitted targets where $E_{k,\nu_k}=0$.
We prove the remaining defect bound next. We emphasize that the activity
constant is free of $M$: an extra factor $C^M$ there would change
$\kappa$.

\begin{lemma}[Interpolation and Taylor defects]\label{lem:defect}
Let $N\ge1$, $\mu\in(0,1)$ and $0<\eta_f\le\eta_0$ satisfy
$\mu(1+\log N)\le1$ and $\eta_f^{4q_{\rm rsp}}N^d\le1$.
Let $\nu_k$ be as in \eqref{def:rows}, and let
$\mathsf R^{\rm def}_k$, $2\le k\le D$, be as in \cref{prop:local}(iv).
There exists a fixed constant $C_5$ such that
\[
 \sum_{k=2}^D\frac{(C_\sharp\mu)^k}{k!}\,
 \norm{\mathsf R^{\rm def}_k}_{(k)}^2
 \le C_5\,\eta_f^4\mu^2N^{-d}.
\]
\end{lemma}

\begin{proof}
By the radial integral
$\int_{\mathbb R^d}(1+N|x|)^2e^{-2N|x|}\,dx=CN^{-d}$, the explicit
expression at count two has squared norm at most $C\eta_f^4N^{-d}$ for every
$N\ge1$. This part needs no additional regime assumption.
For $k\ge3$, the heat identity decomposes the defect into the sum of
\[
 \mathsf D_k=-\eta_f^2\prod_lp_l\sum_iw_i^2
 (1-\widehat w_i^{[k]})\partial_{V_i}\mathsf P,\qquad
 \mathsf Q_k=\prod_lp_l\mathcal Q_{E_{k,\nu_k}(\mathbf u)}.
\]
If $\nu_k\ge1$, interpolation gives
\[
 \int\sum_i(1-\widehat w_i^{[k]})
 \le C^k\nu_k^{-d}(1+\log\nu_k)^{k-2}
 \le C^kN^{-d}\mu^{-(k-2)}.
\]
If $\nu_k<1$, the left side equals $k|Q|^k$, and the same bound
follows from $N^d\mu^{k-2}<1$. Since the defect weights lie in $[0,1]$,
we have
$(\sum_i(1-\widehat w_i^{[k]}))^2\le k\sum_i(1-\widehat w_i^{[k]})$.
Using the fixed derivative bounds, the nuisance suprema and the sum over
responses, we have
$\|\mathsf D_k\|_{(k)}^2\le C^k\eta_f^4N^{-d}\mu^{-(k-2)}$.
The corresponding factorial sum is at most $C\eta_f^4\mu^2N^{-d}$.

For the Taylor remainder, the band estimate with lower cutoff below one
gives $\int\|E_{k,\nu_k}\|_1^2\le C^k$, also for omitted targets.
Combining this with \eqref{eq:LB-response-error}, we find
$\|\mathsf Q_k\|_{(k)}^2\le C^k k^{4q_{\rm rsp}+4}
\eta_f^{4q_{\rm rsp}+4}$. The corresponding factorial sum is at most
$C\eta_f^{4q_{\rm rsp}+4}\mu^2$, which is absorbed using
$\eta_f^{4q_{\rm rsp}}N^d\le1$. The assertion follows from these two
bounds, the pair bound and the squared triangle inequality.
\end{proof}
\section{Activity and aliases}\label{sec:LB-alias}

Fix $K$. We work in $\mathsf{Reg}(K)$ and take $M$ sufficiently large.
In this regime, we have
\begin{equation}\label{eq:LB-hierarchy}
 D=O(M),\qquad H=O(M^2),\qquad
 M/K\le\omega\le KM,\qquad \log N\ge M^2/K.
\end{equation}
These bounds imply that $\ell_\circ>1$, $\mu H\le1$, and
$\omega>2\log H$ eventually. A nonzero fine band requires
$\nu_r>\ell_\circ$, which implies that
\[
 r<2+\frac{dc_0D}{\omega-\log H}\le R,
\]
for some fixed integer $R$ depending on $K$ and the model constants.
Eventually $R\le M$, so every fine guard is legal.
We choose a sufficiently large fixed $C_{\rm act}$ and define
\begin{equation}\label{eq:LB-alias-constants}
 q_a=C_{\rm act}DH(\mu H)^{a/d},\qquad a=1,2.
\end{equation}
Since $\log q_a\le O(\log M)-aM/(Kd)$, we may choose $M_0(K)$ so that
\begin{equation}\label{eq:LB-row-small}
 \begin{gathered}
 M\ge\max\{M_1',R\},\qquad \tau_M\le c_0,\qquad
 \mu H\le1,\\ \log H\le M,\qquad
 q_2\le q_1\le\tfrac12,\qquad D^2q_1\le1
 \end{gathered}
\end{equation}
and all pair bounds hold. Further, any fixed multiple of
$De^{c_0(D+1)}$ is at most $N$ eventually. This proves that the rows are
well defined.

\begin{lemma}\label{lem:activity}
The centered update measure has $B_{\rm loc}\le C_BN^2e^{\tau_MM}$
for a fixed $C_B$, as asserted in \cref{prop:local}(ii).
\end{lemma}
\begin{proof}
An ordinary target-$r$ band with guard $m$ and representation cost
$\mathcal C$ has variation at most $C^rD^re^{\tau_Mm}\mathcal C$.
For $r\ge3$, using the interpolation cost and
$\min\{\nu_r,\ell_\circ\}^2\le\ell_\circ\nu_r$, we have
\[
 B_r^{\rm co}\le CD^2e^{\tau_MD}N\ell_\circ q_1^{r-2},\qquad
 B_r^{\rm fi}\le CD^2e^{\tau_MM}N^2q_2^{r-2}.
\]
Since $e^{\tau_MD}\ell_\circ/N\le1$, the geometric sums of these bounds
are at most $CD^2N^2(q_1+e^{\tau_MM}q_2)\le CN^2e^{\tau_MM}$.
The singleton costs at most $CDe^{c_0(D+1)}\le N$.
It remains to add the pair bound from \cref{lem:new-pair} and to include
the fixed factor from the reference centering. No factor depending exponentially on $M$ is added
to the activity.
\end{proof}

\begin{lemma}\label{lem:alias}
The aliases vanish through count $M$ and, for a fixed $C_6$,
\[
 \sum_{k=M+1}^D\frac{(C_\sharp\mu)^k}{k!}
 \|\mathsf R_k^{\rm al}\|_{(k)}^2
 \le C_6\eta_f^4e^{2\tau_MM}N^{4-d}
       \frac{(C_6\mu)^{M+1}}{(M+1)!}.
\]
This is \cref{prop:local}(vi).
\end{lemma}
\begin{proof}
The pair alias is bounded by \eqref{eq:new-pair-alias}.
All other aliases come from nonzero fine bands with $3\le r\le R\le M$,
so none of them occurs for $k\le M$. Let $E$ be any symmetric matrix
kernel. By the uniform response and selector bounds, we have
\begin{equation}\label{eq:LB-subset-lift}
 \left\|\mathfrak A_k^{r,M}\sum_{|S|=r}\prod_{i\in S}p_i\,
 \mathcal B_{E(\mathbf u_S)}\mathsf P\right\|_{(k)}
 \le C^k\eta_f^2e^{\tau_MM}\|E\|_{L^2(\ell^1)}.
\end{equation}
Because the response and selector bounds hold uniformly in the
nuisance values at each tuple, they bound the supremum in
\eqref{eq:LB-norm}. Summing over subsets and responses and integrating
over the unused coordinates, we obtain the factor
$\binom kr k^2 3^{k/2}|Q|^{(k-r)/2}\max\{1,p_+\}^k$, which is absorbed
in $C^k$.
For a nonzero fine band, interpolation gives
\[
 \|E_r^{\rm fi}\|_{L^2(\ell^1)}^2
 \le C^r\ell_\circ^{4-d}(1+\log\ell_\circ)^{r-2}
 \le C^rN^{4-d}e^{C M},
\]
where we used $r\le R$, $D=O(M)$ and $\log H\le M$.
With at most $R$ such targets, their sum satisfies
\[
 \sum_{k=M+1}^D\frac{(C_\sharp\mu)^k}{k!}
 \|\mathsf R_k^{\rm al}-\mathsf U_k\|_{(k)}^2
 \le C\eta_f^4e^{2\tau_MM}N^{4-d}e^{C M}
       \frac{(C\mu)^{M+1}}{(M+1)!}.
\]
Here we used $\sum_{k\ge M+1}x^k/k!\le e^xx^{M+1}/(M+1)!$ and $\mu<1$.
We absorb $e^{C M}$ into a fixed constant raised to the power $M+1$, and
we combine the result with the pair alias by the squared triangle
inequality. This changes only the constant in the factorial tail, not
the activity exponent.
\end{proof}
\section{A positive path of priors}\label{sec:LB-heaps}

In this section, we use the centered marks $(E,\pi,\mathsf a^\circ)$
and the updates of \cref{sec:LB-local}, with $M\ge M_1'$, $n_h\ge4$ and
vacuum \(\vartheta_{\rm vac}=\vartheta^1_{\rm vac}\). Recall that updates
at disjoint patches commute, that their weights are state independent,
and that $|\mathsf a^\circ|\le B_{\rm loc}$. We fix the constants
\begin{equation}\label{eq:LB-prior-constants}
 \Delta=3^d,\qquad \varrho=2^{-(\Delta+1)},\qquad
 c_{\rm pos}=\frac1{2+8\Delta^2},\qquad
 C_{\rm mgf}=\frac{2^{2d}(p_+-p_-)^2}{8}.
\end{equation}
We assign polynomial weights to finite marked histories of local
updates, and we identify histories that differ only in the order of
disjoint updates. The next lemma proves positivity and a differentiation
identity. Conditional annihilation then keeps the law of the raw density
fixed, and bounded differences control its mass. The interval on which
positivity holds depends on $B_{\rm loc}$, but not on the number of
patches.

\begin{lemma}[Positive realization]\label{lem:positive-realization}
If $\delta B_{\rm loc}/\varrho\le c_{\rm pos}$, there are a standard
Borel history space $\mathfrak H$, a jointly measurable raw-state map
$\vartheta(\mathfrak h)$ and probability measures $\nu_t$,
$0\le t\le\delta$, such that, for every bounded state observable $G$
with measurable history and mark evaluations,
\begin{equation}\label{eq:LB-weak-transport}
 \frac d{dt}\int G\,d\nu_t=\int\mathcal AG\,d\nu_t.
\end{equation}
The product rule holds if $G$ is also continuously differentiable in
$t$ in supremum norm. The law of
$\mathfrak m=\int p_{\vartheta(\mathfrak h)}$ is independent of $t$.
Moreover
\begin{equation}\label{eq:LB-mass-mgf}
 \int\mathfrak m\,d\nu_0=1,\qquad
 \int e^{u(\mathfrak m-1)}\,d\nu_0
 \le e^{C_{\rm mgf}u^2h^d}\quad(u\in\mathbb R).
\end{equation}
For fixed $n\ge1$, define
\begin{equation}\label{eq:LB-tilted-prior}
 G_n=\int\mathfrak m^n\,d\nu_0,\qquad
 d\sigma_t=G_n^{-1}\mathfrak m^n\,d\nu_t.
\end{equation}
These are probability measures with time-independent mass law, and
whenever $\varepsilon\ge8C_{\rm mgf}nh^d$,
\begin{equation}\label{eq:LB-mass-tail}
 \sigma_t(|\mathfrak m-1|>\varepsilon)
 \le2\exp\left\{-\frac{\varepsilon^2}
                         {8C_{\rm mgf}h^d}\right\}.
\end{equation}
\end{lemma}

\begin{proof}
We call two patch labels dependent if they are equal or adjacent, and
independent otherwise. We identify finite words that differ by exchanges
of consecutive independent labels, and we call an equivalence class a
\emph{shape}. The pieces of a shape are indexed by label and occurrence.
We order dependent pieces as in a representative word and take the
transitive closure. Shapes are the elements of the free partially
commutative monoid of Cartier and Foata \cite{CartierFoata1969}, and
with this order they are heaps of pieces in the sense of Viennot
\cite[Section~3]{Viennot1986}.
Incomparable pieces have independent labels, and any compatible ordering
represents the same shape. To see this, we move the first piece of one
ordering to the front of another, passing only incomparable pieces, and
repeat. Applying the marked updates therefore gives a well-defined raw
state. As the shapes are countable, the disjoint union of their finite
products of $E$ is standard Borel. We choose a canonical word for each
shape. The finite update recursion along this word makes the density,
field, mass and likelihoods jointly measurable, also after a mark is
appended.

In the cover graph, two pieces are joined when one covers the other in
this partial order. Each piece has at most $2\Delta$ neighbors. Indeed,
above and below it there is at most one covering piece of each dependent
label, since equal labels are comparable. Maximal pieces have distinct
labels. Removing a maximal piece is the inverse of appending its label
and mark, and both operations are measurable finite deletions and
insertions of coordinates.

We give each shape $\mathfrak g$ weight $\varrho^{|\mathfrak g|}$
and, conditionally on the shape, we give its pieces independent $\pi$
marks. This defines a finite reference measure. Indeed, for fixed
multiplicities $(n_j)$, the interleavings of every adjacent pair of
label chains determine the shape, and their number is at most
\[
 \prod_{\{j,k\}:j\sim k}2^{n_j+n_k}
 \le2^{(\Delta-1)\sum_jn_j}.
\]
Summing over the multiplicities, we see that the normalizer is bounded
by $\prod_j\sum_{r\ge0}(\varrho2^{\Delta-1})^r<\infty$, where
$\varrho2^{\Delta-1}=1/4$. Let $\tau_0$ be the normalized reference
law, with the empty shape included.

For an upper set $U$ of a history $\mathfrak h$, let
\[
 v_t(U)=\operatorname{Leb}\{(s_x)_{x\in U}\in[0,t]^U:
                    s_x<s_y\text{ whenever }x<y\},\qquad
 v_t(\varnothing)=1,
\]
and define
\begin{equation}\label{eq:LB-density}
 \varphi_t(\mathfrak h)=
 \sum_{U\text{ upper in }\mathfrak h}
 \varrho^{-|U|}v_t(U)\prod_{x\in U}\mathsf a^\circ(e_x).
\end{equation}
Since $v_t(U)=t^{|U|}v_1(U)$, this is a finite measurable polynomial.
The weights $v_1(U)$ are the volumes of the order polytopes of the
induced posets on $U$; see \cite[Section~1, Corollary~4.2]{Stanley1986}.

To prove positivity, we let $F(S)$ be the sum in \eqref{eq:LB-density}
restricted to upper sets whose maximal pieces lie in $S\subseteq R$,
where $R=\max(\mathfrak h)$. Connected components of an upper
set are upper sets, since every upward cover chain stays inside it. Two
upper sets intersect exactly when they share a maximal piece of the
history. Disjoint upper sets are incomparable, and their order-volume
weights multiply. Separating the component $P$ containing $x\in S$, we
have
\[
 F(S)=F(S\setminus\{x\})+
 \sum_{\substack{P\text{ connected upper}\\x\in\max P\subseteq S}}
 b_t(P)F(S\setminus\max P),\qquad
 |b_t(P)|\le\epsilon^{|P|},\quad
 \epsilon=\frac{tB_{\rm loc}}\varrho.
\]
We prove by induction on $|S|$ that $F(S)>0$ and
$1/2\le F(S)/F(S\setminus\{x\})\le3/2$. We adapt the deletion induction
of Dobrushin \cite{Dobrushin1996} for hard-core partition functions, as
credited and presented by Scott and Sokal \cite[Section~5.1, Theorem~5.1]{ScottSokal2005},
to the connected upper sets of a history.
The induction hypothesis on proper subsets gives
$F(S\setminus\max P)/F(S\setminus\{x\})\le2^{|\max P|-1}$. A connected set of $k$ vertices containing $x$ is
the vertex set of a depth-first walk of length $2(k-1)$ along a spanning
tree, so there are at most $(2\Delta)^{2(k-1)}$ choices. Combining these
bounds, we conclude that
\[
 \left|\frac{F(S)}{F(S\setminus\{x\})}-1\right|
 \le\sum_{k\ge1}(8\Delta^2)^{k-1}\epsilon^k
 =\frac\epsilon{1-8\Delta^2\epsilon}\le\frac12.
\]
At the stated threshold, the denominator is positive. Multiplying the
ratios, we find
$2^{-|R|}\le\varphi_t(\mathfrak h)\le(3/2)^{|R|}$.

Next, we differentiate the order volumes. A coordinate can reach the
upper face $s_x=t$ only when $x$ is maximal, and removing it leaves the
order volume on the remaining pieces. Because the maximal pieces of an
upper set are maximal in the whole history, we have
\[
 \partial_t\varphi_t(\mathfrak h)=\varrho^{-1}
 \sum_{x\in\max(\mathfrak h)}\mathsf a^\circ(e_x)
                                  \varphi_t(\mathfrak h\setminus x).
\]
Appending a specified label adds one independent $\pi$ mark and
multiplies the reference weight of the shape by $\varrho$. Changing
variables between a history and the history obtained by deleting a
maximal piece, we obtain
\[
 \frac d{dt}\int G\varphi_t\,d\tau_0
 =\int(\mathcal AG)\varphi_t\,d\tau_0.
\]
Here, on histories, $\mathcal A$ appends a marked update, and on state
observables it is given by \eqref{eq:LB-generator}. Since
$|R|\le|\mathcal J|$, the bounds on $\varphi_t$ and its derivative are
uniform for fixed $n$. We then use dominated convergence to justify this
identity, its product rule and continuous differentiability in total
variation.

Let $H$ be a bounded observable of the raw density. For a fixed incoming state and patch, its updated value depends on
the mark only through $e_{\rm d}$, and \eqref{eq:LB-annihilation} then implies
that $\mathcal AH=0$. Taking $H=1$, we find that
$d\nu_t=\varphi_t\,d\tau_0$ is a probability measure, and general $H$
shows that the law of the raw density is invariant. As
$a_M\le\mathfrak m\le b_M$, \eqref{eq:LB-tilted-prior} then defines
probability measures with a constant normalizer and an invariant mass
law.

It remains to prove the mass bound under $\tau_0=\nu_0$.
We condition on a shape. Its marks are independent, and by the
centering identity \eqref{eq:LB-row-centered}, every fresh density
update has expected output one at every point of its patch, for any
incoming value. Starting from the vacuum, we find that the conditional mean of
the mass is one. We group all marks of each label into one block, so that the
blocks are independent. Changing one block changes the final density
only on its patch, because all later density updates act pointwise. The
resulting change in mass is at most $c=2^dh^d(p_+-p_-)$. Revealing
these blocks in turn gives centered martingale differences whose
conditional ranges have length at most $c$. The second derivative of the
conditional log moment-generating function of each difference is its
variance under exponential tilting, which is at most $c^2/4$.
Integrating twice, we bound each log moment-generating function by
$u^2c^2/8$. There are $h^{-d}$ blocks, and averaging over the shape
proves \eqref{eq:LB-mass-mgf}. This is the method of bounded differences
of McDiarmid \cite[Lemma~5.8 and Theorem~6.7]{McDiarmid1989}, applied
conditionally on the shape.

We abbreviate $Z=\mathfrak m-1$ and $a=C_{\rm mgf}h^d$.
By Jensen's inequality, we have $G_n\ge1$. Further,
$\mathfrak m^n\le e^{nZ}$.
For $\varepsilon\ge8na$, the exponential Markov inequality with
$u=\varepsilon/(2a)-n\ge0$ gives
\[
 \sigma_0(Z>\varepsilon)
 \le e^{-u\varepsilon}\int e^{(n+u)Z}\,d\nu_0
 \le e^{n\varepsilon-\varepsilon^2/(4a)}
 \le e^{-\varepsilon^2/(8a)}.
\]
On the event $Z<-\varepsilon$, we have $\mathfrak m^n\le1$, and the
same exponential bound gives
$\sigma_0(Z<-\varepsilon)\le e^{-\varepsilon^2/(4a)}$. Adding the two
tails and using the invariance of the mass law, we obtain
\eqref{eq:LB-mass-tail} for every $t$.
\end{proof}
\section{From local updates to risk}\label{sec:LB-fisher}

Throughout this section, we work with the constructed rows and assume
that $M\ge M_1'$, $n_h\ge4$, $0<\eta_f=c_fh^s\le\eta_0$ and
$C_wc_f\le L_s$. We set
\begin{equation}\label{eq:LB-priors-constants}
 c_m=\frac1{2p_+},\qquad
 \mathcal A_{\mathfrak m}
 =\{\mathfrak h:|\mathfrak m(\mathfrak h)-1|\le c_m/M\},
\end{equation}
and define the full local score energy
\begin{equation}\label{eq:LB-score-energy}
 \mathcal E=\sum_{k\ge0}\frac{(C_\sharp\mu)^k}{k!}
                         \|\mathsf R_k\|_{(k)}^2,
 \qquad C_\sharp=\frac1{p_-^2c_q},\qquad \mu=nh^d.
\end{equation}
Here the numerator and the norm are given by \eqref{eq:LB-numerator}
and \eqref{eq:LB-norm}, and the series is finite by the full-response
bound \eqref{eq:LB-full-response}.

\begin{proposition}[A lower bound from local cost]\label{lem:score-risk}
Suppose $c_m/M\ge8C_{\rm mgf}\mu$, and put
\begin{equation}\label{eq:LB-bad-mass}
 \epsilon_n=2\exp\left\{-\frac{c_m^2}{8C_{\rm mgf}M^2h^d}\right\}.
\end{equation}
Then \eqref{eq:LB-cost-risk} holds with a constant $c_{\rm tr}>0$
depending only on the model constants.
\end{proposition}

\begin{proof}
Set
\[
 I=\sqrt{3^dh^{-d}\mathcal E},\qquad
 \mathcal C=1+B_{\rm loc}+I,\qquad
 c_* =\min\{1,\varrho c_{\rm pos},\rho/\eta_0^2\},\qquad
 \delta=c_*/\mathcal C.
\]
Then $\delta B_{\rm loc}/\varrho\le c_{\rm pos}$ and
$\eta_f^2\delta\le\rho$. \Cref{lem:positive-realization} applies and
gives the probability priors \eqref{eq:LB-tilted-prior} on $[0,\delta]$.
Let $V_t=v-\eta_f^2t$, and let $\mathbb E_t$ denote prior expectation
followed by expectation under the normalized observation law at
$V_t$. The choice of the ternary neighborhood and the bound on the
variance decrease ensure that $V_t\in J_V$. By \cref{lem:legality}, every state in $\mathcal A_{\mathfrak m}$ is
legal, and the mass tail bound implies that
$\sup_{t\le\delta}\sigma_t(\mathcal A_{\mathfrak m}^c)\le\epsilon_n$.

For bounded observables $T$ of the data, the mass tilt cancels the
likelihood normalization exactly:
\[
 \mathbb E_tT=G_n^{-1}\int\nu_t(d\mathfrak h)
       \int T(\xi)\mathsf U(\vartheta(\mathfrak h),V_t;\xi)
                         \,d\Lambda_n(\xi).
\]
The inner integral is continuously differentiable in supremum norm
for fixed $n$, since the response factors are affine in $V_t$ and
$\Lambda_n(\mathfrak X)=3^n$. Using the product rule in \eqref{eq:LB-weak-transport}, we obtain
\begin{equation}\label{eq:LB-observable-derivative}
 \frac d{dt}\mathbb E_tT=\mathbb E_t[TR],\qquad
 R=\frac{\mathcal A\mathsf U-\eta_f^2\partial_V\mathsf U}{\mathsf U}
   =\sum_j\mathsf r_j,\qquad \mathbb E_tR=0.
\end{equation}
Here the score identity is \eqref{eq:LB-score-sum}, and the centering
follows from \cref{lem:score}.

At a fixed raw state, we group the observations according to their
count $k$ in $Q_j$. Using the chart ratio, $dx=h^ddu$, and the bound of
one on the probability that the remaining observations lie outside the
patch, we have
\[
 \E_{\vartheta,V}\mathsf r_j^2
 \le\sum_{k=0}^n\binom nk\Bigl(\frac{h^d}{m_\vartheta}\Bigr)^k
 \int_{Q^k}\sum_{\mathbf y}\frac{\mathsf R_k^2}{\mathsf L_k}\,d\mathbf u
 \le\mathcal E.
\]
In the second inequality, we used $m_\vartheta\ge p_-$,
$\mathsf L_k\ge(p_-c_q)^k$, $\binom nk\le n^k/k!$ and the uniform bound
in \eqref{eq:LB-norm}, whose supremum is taken before integration. By
\cref{lem:score}, disjoint patches have orthogonal full scores. Each
label has $3^d$ dependent labels, including itself, and applying
$2ab\le a^2+b^2$ to the pairs of dependent labels yields
\begin{equation}\label{eq:LB-residual-second-moment}
 \mathbb E_tR^2\le3^d\sum_j\mathbb E_t\mathsf r_j^2
                  \le3^dh^{-d}\mathcal E=I^2.
\end{equation}
The integrated score is at most $\delta I\le c_*$, and the separation
is $\eta_f^2\delta$. Since $D_V\le v_+-v_-$, \cref{lem:path-risk}
implies that
\[
 \mathfrak r_n^2\ge
 \frac{c_*^2\eta_f^4}{(2+c_*)^2\mathcal C^2}-(v_+-v_-)^2\epsilon_n.
\]
Since $\mathcal C^2\le3^{d+1}(1+B_{\rm loc}^2+h^{-d}\mathcal E)$,
this proves \eqref{eq:LB-cost-risk} with
$c_{\rm tr}=c_*^2/[3^{d+1}(2+c_*)^2]$.
\end{proof}

\begin{lemma}[Cost of the local construction]\label{lem:raw-energy}
Let $(M,D,N,\mu,\eta_f)\in\mathsf{Reg}(K)$ with $M\ge M_0(K)$,
$\mu=nh^d$ and $\eta_f=c_fh^s$. Set $\mathcal B=N^2e^{\tau_MM}$.
There is a fixed $C_G\ge1$, depending only on the model and design
constants and $K$, such that
\begin{equation}\label{eq:LB-raw-energy}
 1+B_{\rm loc}^2+h^{-d}\mathcal E
 \le C_G\{\mathcal B^2+\mathfrak E_1+\mathfrak E_2+\mathfrak E_3\},
\end{equation}
where
\begin{equation}\label{eq:LB-terms}
\begin{aligned}
 \mathfrak E_1&=n^2h^{d+4s}N^{-d},\\
 \mathfrak E_2&=n h^{4s}e^{2\tau_MM}N^{4-d}
                    \frac{(C_G\mu)^M}{(M+1)!},\\
 \mathfrak E_3&=n h^{4s}\mathcal B^2
                    \frac{(C_G\mu)^D}{(D+1)!}.
\end{aligned}
\end{equation}
\end{lemma}

\begin{proof}
The score numerators at counts zero and one vanish. For $2\le k\le D$, we square the
decomposition of \cref{prop:local} into defect, alias and field terms,
at the cost of a factor three. For $k>D$, using
\eqref{eq:LB-full-response}, the volume $2^{dk}$ and the $3^k$ response
vectors, we have
$\|\mathsf R_k\|_{(k)}^2\le(3\cdot2^dC_{\rm hi}^2)^k
\eta_f^4(B_{\rm loc}+1)^2$. Let
$c_{\rm h}=3\cdot2^dC_{\rm hi}^2C_\sharp$. Using
$\sum_{k\ge a}x^k/k!\le e^xx^a/a!$, we find
\[
\begin{aligned}
 \mathcal E\le{}&3(C_5+C_7)\eta_f^4\mu^2N^{-d}\\
 &+3C_6\eta_f^4e^{2\tau_MM}N^{4-d}
                          \frac{(C_6\mu)^{M+1}}{(M+1)!}\\
 &+e^{c_{\rm h}}\eta_f^4(B_{\rm loc}+1)^2
                          \frac{(c_{\rm h}\mu)^{D+1}}{(D+1)!}.
\end{aligned}
\]
We multiply by $h^{-d}$ and use $h^{-d}\mu=n$,
$h^{-d}\mu^2=n^2h^d$ and $\eta_f=c_fh^s$.
As $B_{\rm loc}\le C_B\mathcal B$ and $\mathcal B\ge1$, the claimed
bound holds if we take $C_G$ to be the maximum of
\[
 1,\ C_6,\ c_{\rm h},\ 1+C_B^2,\ 3(C_5+C_7)c_f^4,\quad
 3C_6^2c_f^4,\ e^{c_{\rm h}}c_{\rm h}c_f^4(C_B+1)^2.
\]
All these constants are fixed before $n$.
\end{proof}


\section{Parameters and proof of \texorpdfstring{\cref{thm:lower}}{Theorem 6.1}}\label{sec:LB-proof}

We choose the parameters by balancing the squared activity with the
interpolation contribution $\mathfrak E_1$, which leads to
$N^{d+4}=n^2h^{d+4s}e^{-2\tau_MM}$. The remaining contributions
are factorial tails. The choice $M\asymp\sqrt{\log n}$ makes these tails
smaller than the principal cost, while the correction $-\log M$ in
$\omega$ produces the logarithmic power in the lower bound. We define
\[
\begin{gathered}
 G=\frac{d-4s}{d},\qquad \gamma=\frac{G}{d+4},\qquad
 \theta=\frac{d\tau}{2(s-1)},\\
 m_*^2=\frac{8\gamma(s-1)}{d\tau},\qquad C_D=\frac{d+8}{4},
 \qquad c_f=\min\{1,L_s/C_w\}.
\end{gathered}
\]
Fix $K$ such that $\theta$ and $\gamma/m_*^2$ lie strictly inside
$[1/K,K]$. At this fixed $K$, we take the constants of \cref{prop:local}
and then $C_G$ from \cref{lem:raw-energy}. Finally, we fix a constant
$C_\omega>A+2$, where
\[
 A=1+\log C_G+\frac{d^2-16s}{d(d+4)}\theta+\frac{2d\tau}{d+4}.
\]
All these constants are fixed before $n$. With $L=\log n$, we set
\begin{equation}\label{eq:LB-parameters}
\begin{gathered}
 M=\lceil m_*\sqrt L\rceil,\qquad D=\lceil C_DM\rceil,\qquad
 \omega_0=\theta M-\log M+C_\omega,\qquad
 n_h=\lceil(ne^{\omega_0})^{1/d}\rceil,\\
 h=1/n_h,\quad\mu=nh^d,\quad\omega=\log(1/\mu),\quad
 \eta_f=c_fh^s,\quad
 N=\left(n^2h^{d+4s}e^{-2\tau_MM}\right)^{1/(d+4)},\\
 \mathcal B=N^2e^{\tau_MM},\qquad R_{\rm low}=\eta_f^2/\mathcal B.
\end{gathered}
\end{equation}

\begin{lemma}\label{lem:saddle}
For all sufficiently large $n$, the tuple
$(M,D,N,\mu,\eta_f)$ belongs to $\mathsf{Reg}(K)$, and the local
construction satisfies every hypothesis of \cref{lem:score-risk}.
Moreover,
\[
\begin{gathered}
 1+B_{\rm loc}^2+h^{-d}\mathcal E\le3C_G\mathcal B^2,\\
 R_{\rm low}\asymp n^{-\lambda}e^{-\kappa\sqrt L}L^{(s-1)/(d+4)},
 \qquad \epsilon_n=o(R_{\rm low}^2).
\end{gathered}
\]
\end{lemma}

\begin{proof}
Rounding upward gives
\begin{equation}\label{eq:LB-omega}
 \omega_0\le\omega\le\omega_0+d\log2,\qquad
 \omega=\theta M-\log M+C_\omega+O(1),
 \qquad\tau_MM=\tau M+O(1).
\end{equation}
Since $h^d=e^{-\omega}/n$, we find from the definition of $N$ that
\begin{equation}\label{eq:LB-N}
 (d+4)\log N=GL-(1+4s/d)\omega-2\tau_MM
 =GL+O(\sqrt L).
\end{equation}
We see that $\omega/M\to\theta$ and $\log N/M^2\to\gamma/m_*^2$,
which verify the regime bounds at the fixed $K$. Also, $h\le n^{-1/d}$
eventually, and since $q_{\rm rsp}\ge d/(4s)-1$, we have
\[
 \eta_f^{4q_{\rm rsp}}N^d
 \le Cn^{(d-4s)/(d+4)-4q_{\rm rsp}s/d}
 \le Cn^{-4\gamma}\longrightarrow0.
\]
All fixed thresholds on $M,n_h,\eta_f$ hold eventually, and the mass
condition $c_m/M\ge8C_{\rm mgf}\mu$ follows from $M\mu\to0$.
The model bound $C_wc_f\le L_s$ holds by construction, and
\cref{prop:local} supplies annihilation and reference centering.

For the cost, we have $\mathfrak E_1=\mathcal B^2$ exactly.
With $c_t=(d^2-16s)/\{d(d+4)\}$, direct substitution gives
\begin{equation}\label{eq:LB-T2}
\begin{aligned}
 \log\frac{\mathfrak E_2}{\mathcal B^2}
 &=4\gamma L+(c_t-M)\omega+
 \left(\frac{2d\tau_M}{d+4}+\log C_G\right)M-\log(M+1)!\\
 &=4\gamma L-M(\omega+\log M)+AM+O(\log M)
 \le(-C_\omega+A+o(1))M.
\end{aligned}
\end{equation}
Here $A$ is independent of $C_\omega$, while the remainder may depend
on that subsequently fixed constant. In the final inequality, we used
$\theta M^2\ge4\gamma L$ and $\omega\ge\omega_0$. This shows that the
alias tail tends to zero relative to $\mathcal B^2$.
For the high-count tail, we have
\[
 \log\frac{\mathfrak E_3}{\mathcal B^2}
 \le GL-D\omega+O(D)
 =(G-4\gamma C_D)L+o(L)=-4\gamma L+o(L).
\]
The cost bound follows from \cref{lem:raw-energy}.

For the risk scale, we compute
\begin{equation}\label{eq:LB-saddle}
\begin{aligned}
 \log R_{\rm low}
 &=-\lambda L-\frac{2(s-1)\omega+d\tau_MM}{d+4}+O(1)\\
 &=-\lambda L-\frac{2d\tau}{d+4}M
   +\frac{2(s-1)}{d+4}\log M+O(1)\\
 &=-\lambda L-\kappa\sqrt L+\frac{s-1}{d+4}\log L+O(1).
\end{aligned}
\end{equation}
Finally, $M^2h^d\le CL/n$ implies that
$\epsilon_n\le2e^{-cn/L}=o(R_{\rm low}^2)$.
\end{proof}

\begin{proof}[Proof of \cref{thm:lower}]
We take the parameters \eqref{eq:LB-parameters} for all sufficiently
large $n$. Combining \cref{lem:saddle} with \cref{lem:score-risk}, we
obtain
\[
 \mathfrak r_n^2\ge\frac{c_{\rm tr}}{3C_G}R_{\rm low}^2
 -(v_+-v_-)^2\epsilon_n\ge\frac{c_{\rm tr}}{6C_G}R_{\rm low}^2
\]
eventually. Taking square roots and using \eqref{eq:LB-saddle}
proves the stated lower bound.
\end{proof}
